% =====================================================================
%  Harald Høffding, RELIGIØSE TANKETYPER.
%  København: Gyldendalske Boghandel, Nordisk Forlag, MCMXXVII (1927).  165 pp.
%
%  GENERATED; DO NOT EDIT BY HAND.  Built by ebook-method/religioese-tanketyper/
%  (rebuild.sh) from the scan, the collation decisions and the patches there.
%
%  SOURCE: the National Library of Norway's scan (nb.no item
%  7cfe6d92ba67c5ca2c2e33faa1f19da7, URN:NBN:no-nb_digibok_2012041608093),
%  ~/bibliotek/Høffding, Harald/1927-religioese-tanketyper.pdf (SCANS.tsv).
%  PDF page = printed page + 4.
%
%  EDITION CHECK (2026-10-17): one edition (Gyldendal 1927; NB's record 1927,
%  165 s.; the title page MCMXXVII; the Fortale dated 18 April 1927).  The SAGA /
%  Lindhardt og Ringhof e-book (ISBN 9788726363050, "Copyright 1927, 2020") is
%  the same text (same Fortale, same fourteen pieces in the same order, the 1927
%  orthography); it was used only as a voting witness in the collation.
%
%  HOW ACCURACY WAS ESTABLISHED (2026-10-17; ebook-method/religioese-tanketyper/, README entry there):
%  (a) Draft: NB's own OCR (ALTO) laid over the page images (nb_scan.py), turned into LaTeX by rule (draft.py), with
%  running heads, sheet signatures and the junk above the openings dropped; notes placed at their marks (16 = the
%  e-book's 16).
%  (b) Word collation (collate/check.py, punctuation included) against an INDEPENDENT tesseract reading of the same
%  images: 854 disagreements; 18 settled by check.py's own rules, 731 by rules.py (the e-book's vote where it agrees
%  with exactly one reading, 699 of them; mark/quote/dash rules), 70 by hand where three witnesses agree (NB's l/J/n/..
%  for I and „, the German umlaut rule, junk, headings; tagged hand-3wit), 35 at the image.  OPEN 0; italics OPEN 0.
%  (c) Blind crops (b1: 105 crops on 7 sheets, two readers and the main session, controls 14/14): the open checks, the
%  e-book's differences the collation could not see, letterspacing candidates, and an AUDIT of 40 rule and hand
%  decisions (40/40 right).
%  (d) Structure and all 16 footnotes read WHOLE at the image (s1: 72 crops, two readers): title page, verso, Indhold,
%  the fourteen openings, the six dividers (three stars) and the eleven rules, verse, doubtful paragraph starts.
%  (e) The whole text compared word by word with the e-book (ebookdiff.py): 165 differences, every one not of style
%  settled at the image or by the e-book where both OCRs agree with it.  Sweeps: unknown words, glued words, specks,
%  straight quotes, quote balance, dropped runs.
%  (f) A blind sampled page check (pagecheck.py): 10 pages from five strata (pp. 9, 17, 46, 55, 77, 101, 128, 129, 153,
%  163), 2,924 words: 0 word errors, 1 punctuation error (a straight quote, mended with its class), 95 % upper bound
%  1.0 word errors per 1000 words.
%  (g) The translators' 61 odd spots: 31 read at 400-dpi line crops: 1 fault (p. 34), 3 divisions whose hyphen the
%  layer lost, the verse on p. 125; 11 misprints of the print, marked PRINTED AS IS.
%  NOT done: a whole-book reading at the image.
%
%  PP. 27-28 ARE NOT FROM THE PRINT.  The NB copy lacks the leaf (PDF 30-33 = pp. 26, 29, 26, 29: shot twice, read at
%  the image; a former reader's pencil note on p. 17 reads „Side 27-28 …“, partly legible).  Their text, „[Grund]værdi sig
%  igen“ to „jeg har [kunnet]“, is set from the SAGA / Lindhardt og Ringhof e-book (2020, ISBN 9788726363050) and NOT
%  read at a print (Hans, 2026-10-17; KB-VISIT-TODO.md).  The turn from p. 27 to p. 28 is not marked.
%
%  EDITORIAL: diplomatic.  Printed as is, and marked: „gør afhængig“ (p. 43), „Ækvalensforhold“ (p. 61), „Han sidste
%  Ord“ (p. 79), „sammme“ (p. 85), „Kirkegaard“ (p. 105), „nøvrig“ (p. 111), „væsenlige“ (p. 112), „Lakaris“
%  (p. 119; „Laskaris“ elsewhere), „Frygt.’“ (p. 121), „Halucinationer“ (p. 140), „Krafte“ (p. 141).  Also as printed:
%  „Neiiendams“, „Boccacio“, „Pomponazzo“, „Marannos“, „Ticho“, „Brunschwicg“, „Jaqueline“, „Goetze“, „Lyssind“.
%  The Indhold's „Platon om Kunst og Religion“ and „De Vises Sten“ against the headings (both kept).  A former reader's
%  pencil (slashes between sentences, crosses, underlining) is not transcribed; where a slash covers a stop, the stop
%  was read at the image.  Running heads (folio and short title) and sheet signatures dropped.  The ø rule (oefix.py):
%  no site.  The print's emphasis is italic (\textit); it has no letterspacing in the text.  Page markers \opage{N} at
%  the page turns (NB's ALTO lines); a word divided across a page turn is joined with %.
% =====================================================================
\documentclass[12pt,a4paper,openany]{book}
\usepackage[utf8]{inputenc}
\usepackage[T1]{fontenc}
\usepackage{libertinus}
\usepackage{textalpha}
\usepackage[danish]{babel}
\usepackage{setspace}
\usepackage[margin=1.2in]{geometry}
\usepackage{fancyhdr}
\setlength{\headheight}{28pt}
\addtolength{\topmargin}{-13.5pt}
\usepackage{hyperref}
\hypersetup{pdftitle={Religiøse Tanketyper (1927)}, pdfauthor={Harald Høffding}, hidelinks}
\pagestyle{fancy}
\fancyhf{}
\fancyhead[LE,RO]{\thepage}
\fancyhead[RE]{\textit{Harald Høffding: Religiøse Tanketyper}}
\fancyhead[LO]{\textit{\leftmark}}
\fancypagestyle{plain}{\fancyhf{}\fancyhead[LE,RO]{\thepage}}
\setstretch{1.3}
\usepackage{marginnote}
\usepackage{graphicx}
\DeclareUnicodeCharacter{0254}{\reflectbox{c}}  % ɔ: the old Danish id-est mark (a reversed c)
\newcommand{\starmark}[1]{\ifnum#1>10 \ifcase\numexpr#1-10\relax\or*\or**\or***\fi\else\arabic{footnote}\fi)}   % \footnote[1] ... = 1) 2) ...; [11], [12] = the print's *) **)
\renewcommand{\thefootnote}{\starmark{\value{footnote}}}   % the print numbers its notes 1) 2) per page, a few *)
\newcommand{\opage}[1]{\marginnote{\footnotesize\textit{[#1]}}}
\newcommand{\apage}[1]{\marginnote{\footnotesize\textit{[#1]}}}
\newcommand{\whiteline}{\par\vspace{\baselineskip}\par}
% the printed Indhold: an essay or chapter line (title, page), a section line
\newcommand{\ientry}[2]{\par\smallskip\noindent #1\dotfill\,#2\par}
\newcommand{\isub}[2]{\par\noindent\hspace*{2em}#1\dotfill\,#2\par}
% the headings (no \addcontentsline: the print has its own Indhold, and a footnote in a heading breaks a toc entry)
\pdfstringdefDisableCommands{\renewcommand{\footnote}[2][]{}}
\newcommand{\prefacehead}[1]{\chapter*{#1}\markboth{Fortale}{Fortale}}
\newcommand{\sectionhead}[1]{\par\medskip\begin{center}#1\end{center}\smallskip}
\newcommand{\essayhead}[2][]{\clearpage\begin{center}{\Large #2}\end{center}\if\relax\detokenize{#1}\relax\markboth{#2}{#2}\else\markboth{#1}{#1}\fi\medskip}   % [running head] when the title carries a note
\newcommand{\essaysource}[1]{\begin{center}{\small #1}\end{center}\medskip}
\newcommand{\subhead}[1]{\par\medskip\begin{center}#1\end{center}\smallskip}
\newcommand{\asterism}{\par\medskip\begin{center}*\hspace{2.5em}*\\[-2pt]*\end{center}\medskip}   % three stars, two above one (pp. 10-16)
\newcommand{\piecerule}{\par\medskip\begin{center}\rule{0.2\textwidth}{0.4pt}\end{center}\medskip}   % the short rule at the end of a piece
\newcommand{\romanhead}[1]{\par\medskip\begin{center}#1.\end{center}\smallskip}

\begin{document}

% ---------------------------------------------------------------------
%  FRONT MATTER (scan PDF 7-12).
% ---------------------------------------------------------------------
\frontmatter
\begin{titlepage}
\centering
\vspace*{2cm}
{\large H A R A L D \ \ H Ø F F D I N G}\par
\vspace{1.5cm}
{\Huge R E L I G I Ø S E\\[8pt] T A N K E T Y P E R}\par
\vspace{1cm}
{[\textit{fleuron}]}\par
\vfill
\rule{6cm}{1.2pt}\\[-8pt]\rule{6cm}{0.4pt}\par
\vspace{6pt}
{\small GYLDENDALSKE BOGHANDEL · NORDISK\\ FORLAG · KØBENHAVN · MCMXXVII}\par
\end{titlepage}
\thispagestyle{empty}
\vspace*{\fill}
\begin{center}
{\footnotesize COPYRIGHT 1927 BY GYLDENDALSKE\\ BOGHANDEL NORDISK FORLAG · COPENHAGEN\\ ---\\ PRINTED IN DENMARK\\ GYLDENDALS FORLAGSTRYKKERI\\ KØBENHAVN}
\end{center}
\vspace*{\fill}
\clearpage

% the Indhold (p. 5, PDF 9), read at the image 2026-10-17 (s1-02-01): four groups set apart by white space
\chapter*{I N D H O L D}
\markboth{Indhold}{Indhold}
\opage{5}%
\noindent\hfill{\small Side}\par
\ientry{Fortale (1927)}{7}
\ientry{Filosofiske Synspunkter (1923)}{9}
\ientry{Om min Forfattervirksomhed (1923)}{17}
\medskip
\ientry{Platon om Kunst og Religion (1924)}{41}
\ientry{Giordano Bruno (1922)}{50}
\ientry{Spinoza (1927)}{59}
\ientry{Pascal og Kierkegaard (1923)}{70}
\ientry{Sibbern og Kierkegaard (1923)}{98}
\ientry{Filosofi og kirkelig Rationalisme (1922)}{112}
\ientry{En moderne Buddhist (1926)}{117}
\medskip
\ientry{„Bakkehus og Solbjerg“ (1924)}{126}
\ientry{„De Vises Sten“ (1923)}{136}
\ientry{„For aaben Dør“ (1926)}{147}
\ientry{„Den grønne Mark“ (1926)}{154}
\medskip
\ientry{Shakespeare: „Stormen“ (1926)}{158}
\clearpage

% --- p. 7 ---
\setcounter{footnote}{0}%

\opage{7}\prefacehead{FORTALE}

Denne Bog indeholder en Række Foredrag og Afhandlinger
fra de seneste Aar. Jeg har følt Trang til at
prøve endnu engang at belyse Spørgsmaalet om Forholdet
mellem Tanke og Tro, der har sysselsat mig fra min Ungdom.
I min „Religionsfilosofi“ har jeg givet den Besvarelse
af Spørgsmaalet, jeg har været istand til at give. Der er
efter min Overbevisning kun tre Veje, ad hvilke det religiøse
Problem kan belyses, en erkendelsesteoretisk, en psykologisk
og en etisk Vej, alle tre Veje hver for sig igen belyste
af, hvad Religionshistorien kan lære os. I den Bog,
jeg her forelægger, har jeg, som Titelen antyder, fremdraget
nogle Typer paa Tanker om Religion, saaledes som de til
forskellige Tider er fremtraadt hos Tænkere og Digtere.

Som Indledning til de historiske Fremstillinger giver jeg
i de to første Stykker en kort Redegørelse for min egen filosofiske
Opfattelse, der naturligvis ligger til Grund baade for
Valget af Typer og for Karakteristiken af dem. Det første
Stykke er en Forelæsning, jeg efter Indbydelse fra Universitetet
holdt i dets Festsal nogle Dage efter min 80 Aars Fødselsdag,
og i hvilken jeg henviste til, at min Filosoferen væsentlig
var bestemt ved Ønsket om at udfinde, hvorledes
man menneskelig maa stille sig dels overfor Kampen for
Tilværelsen, dels overfor den aandelige Arbejdsdeling, der
bl. a. synes at tvinge Videnskab og Religion ind paa helt forskellige
Veje. Det andet Stykke giver en Oversigt over min
Forfattervirksomhed, der kan vise, i hvilket Forhold mit religionsfilosofiske
Arbejde staar til mine andre Arbejder.

Derefter følger Karakteristiker af Tænkere, der har hævdet
videnskabelig Forskens Ret ved Forstaaelsen af det reli%
% --- p. 8 ---
\setcounter{footnote}{0}%
\opage{8}giøse Liv, saavel som af Tænkere, der har ment, at der var
en skarp, en uovervindelig Modsætning mellem intellektuel
og religiøs Interesse. Karakteristikerne af Spinoza og Pascal
udgør de vigtigste Afsnit i Bogen. Det skyldes en Opfordring
fra et fransk filosofisk Tidsskrift, at jeg kom til at anstille
en Sammenligning mellem Pascal og Kierkegaard, hvilken
sidste med en vis Ret kan betegnes som den danske Pascal.
Den sammenlignende Metode har her vist sig gavnlig til Belysning
af vor store religiøse Tænker. Allerede tidligere har
jeg anvendt denne Metode i en lille Afhandling om Kierkegaard
og Nietzsche (i tredie Bind af mine mindre Arbejder).
Ogsaa Afhandlingen om Sibbern og Kierkegaard er Exempel
paa Anvendelse af denne Metode. Og Artiklen om Arthur
Pfungst's Forsøg paa at praktisere Buddhisme midt i
europæisk Kulturliv kan ses under samme Synspunkt, idet
dens Emne danner en Analogi til, hvad Kierkegaard kaldte
„moderne Kristendom“.

Det har i høj Grad interesseret mig, at nogle af vore Digtere
har søgt at gøre sig Rede for den inderlige Forening af
Poesi og Religion, der stedse vil gøre sig gældende, hvergang
en Tro udvikler sig ad Livets egen Vej.

Naar jeg til Slutning lader Shakespeares Prospero faa
Ordet, er det, fordi jeg i ham ser den digteriske Repræsentant
for det Højsind, som paa Renaissancens Tid naturlig
kom til at staa som den ypperste Karakteregenskab baade
i Forholdet til Naturens store Horizonter og skiftende Skæbner,
og i Forhold til de Mennesker, vi møder i Kampen for
Tilværelsen. Her er en Form for Tro, som stadig vil staa
sin Prøve.

Carlsberg d. 18. April 1927.

\hfill\textit{Harald Høffding.}

\mainmatter
% --- p. 9 ---
\setcounter{footnote}{0}%

\opage{9}\essayhead{FILOSOFISKE SYNSPUNKTER}

\essaysource{Forelæsning, holdt paa Universitetet i Anledning af min \\ 80 Aars Fødselsdag.}

Det er en stor Ære, Universitetet viser mig ved at opfordre
mig til at holde en Forelæsning i Anledning af
min 80 Aars Fødselsdag. Paa bedre Maade kunde man ikke
hædre og glæde mig.

Det er nu 64 Aar siden, jeg første Gang gik op ad Universitetstrappen.
Jeg hørte allerede i min Skoletid Carsten
Hauch's Forelæsninger over Folkeviserne, Ewald, Oehlenschläger
og Shakespeare. Saa kom mine Studenteraar, i
hvilke de modstridende Opfattelser, der gjordes gældende
af Filosofer og Teologer, gav mig nok at tænke paa. Og for
52 Aar siden holdt jeg min første Forelæsning som Privatdocent.
Siden den Tid har der været en stadig Vekselvirkning
mellem min Forelæsningsvirksomhed, min personlige
Udvikling og min Forfattervirksomhed. Der er en naturlig
Trang til at meddele, hvad man har syslet med i Lønkammeret,
og gennem Forsøg paa Meddelelse kan Tanken
blive klarere og bestemtere. Jeg har ofte sandet, hvad
Zeuthen engang fortalte mig som sin Erfaring, at baade
Betydningen og Vanskeligheden af en Tanke ofte først ret
kommer frem for os, naar vi forsøger at udtale den. Men
naar jeg tænker tilbage paa min Universitetsvirksomhed,
mindes jeg dog først og fremmest med Taknemmelighed
den Modtagelighed, der vistes overfor, hvad jeg kunde yde,
indenfor de forskellige Kredse af Tilhørere --- den snevre
Kreds af dem, der mere indgaaende vilde orientere sig i
Filosofi, den større Kreds af unge Studenter i deres første
Aar, og den undertiden nok saa store Kreds, der hørte mine
% --- p. 10 ---
\setcounter{footnote}{0}%
\opage{10}Aftenforelæsninger. Og saa bragte Universitetsvirksomheden
mig gode kollegiale Forhold til andre Universitetslærere,
hvoraf ikke sjeldent et varmt personligt Forhold udviklede
sig. Den allerstørste Del af mit Liv har været knyttet
til Universitetet.

\asterism

Som Emne for denne Forelæsning har jeg valgt et Blik
paa de Problemer, der især satte min Eftertanke i Bevægelse,
paa de Tankemotiver, der især har været bestemmende
for mit Standpunkt.

\asterism

Da jeg begyndte min mere selvstændige Filosoferen,
øvede Darwinismen en overvældende Indflydelse paa Sindene.
For mit Vedkommende blev min Eftertanke især sat
i Bevægelse ved Darwin's stærke Betoning af, at selv de
herligste Livsformer maa „kæmpe for Livet“, det vil sige,
gøre et Arbejde, fyldestgøre visse Betingelser for at kunne
bestaa og udvikle sig. Darwin har med Urette i selve Kampen
for Livet set en Aarsag til nye Arters Opstaaen. Men
det forringer ikke Betydningen af det Billede, han har givet
os af Naturen. Naturen fremtræder ikke mere som en Idyl,
men som en Verden fuld af Omskiftelser, Brydninger og
Kampe.

Det Spørgsmaal maatte da naturligt opstaa: hvorledes
gaar det saa det aandelige Liv, som jo dog skal bestaa og
udvikle sig indenfor Naturen, og som er underkastet dennes
Love? Hvorledes kan da under denne Omskiften, der
blot bestemmes ved Livets og Forholdenes ubønhørlige
Krav, de aandelige Goder, Arbejdet i Sandhedens og Godhedens
Tjeneste bevare deres Betydning?

Det Svar, jeg mente at finde, og som jeg i en Række
Skrifter søgte at begrunde, gik i Korthed ud paa, at det
aandelige Liv i Verden selv er en uafbrudt Kamp for at
% --- p. 11 ---
\setcounter{footnote}{0}%
\opage{11}bestaa --- at det er en Virken, et Arbejde, hvad enten Arbejdet
øves med Bevidsthed eller ikke.

Jeg fandt som det aandelige Livs Karaktermærker to
Træk, der netop var af Betydning, naar det gjaldt om at
hævde sig overfor en Omverden. Det ene Træk er Aabenheden,
Modtageligheden overfor det Nye, overfor alle Forandringer.
Det andet er, at det Nye bringes i Sammenhæng
med vore tidligere Erfaringer, saaledes at der bliver
en stadig Vekselvirkning, gennem hvilken snart tidligere
Erfaringer berigtiges gennem de nye, snart de nye
fortolkes ved Hjælp af de tidligere. Dertil udfordres et indre
Arbejde, en Kamp for at overvinde det Kaotiske og Mangfoldige,
der kan sprede, saavel som det Tilvante, der kan
sløve og blinde. Som det aandelige Livs Grundkarakter
staar da en sammenfattende Energi. Energi er Evne til at
udføre et Arbejde; Arbejde forudsætter Modstand; og paa
det aandelige Omraade bestaar Modstanden i Spredthed og
Sløvhed og i den ofte saa skarpe Modsætning mellem det
Nye og det Tilvante. Jo fyldigere og mere omfattende et
Indhold af Iagttagelser og Erindringer, af Tilskyndelser og
af Tilbøjeligheder der optræder, des større Energi kræves
der for at holde sammen paa det og forme det til en Helhed.

Paa dette Punkt blev et Studium af Modsætningen mellem
sygt og sundt Sjæleliv af stor Betydning for mig. Sindssygdom
ytrer sig netop som Sammenhængsløshed eller som
enkelte Elementers ensidige Overvægt. Naar Genier ikke
sjeldent er sindssyge, hænger det sammen med, at der hos
dem fremvælder en saa stor Fylde af Drifter, Tanker og
Billeder, at der kræves en meget stor Energi til at ordne
dem og lede dem. Det er ingen Kunst at være Herre over
sig selv, naar der ikke er meget at være Herre over.

\asterism

Medens Psykologien har til Opgave at vise, hvorledes
det aandelige Liv faktisk udfolder sig som et Arbejde for
% --- p. 12 ---
\setcounter{footnote}{0}%
\opage{12}at bestaa, er det Erkendelsesteoriens Opgave at undersøge,
hvilke Tankeformer der ligger til Grund, naar der ikke
blot spørges om vore virkelige Tanker, men om vore gyldige
Tanker, det vil sige, de Tanker, der kan føre til sikker,
atter og atter bekræftet Orientering i Tilværelsen udenom
os og i os selv. Ogsaa Videnskab er Kamp mod et Kaos
og mod det Tilvante, et Arbejde for at komme ud over det
Spredte og det Blinde. Det er Filosofiens Opgave at paavise
de Grundtanker („Kategorier“), som gør sig gældende
i dette Arbejde, baade naar der spørges, og naar der svares.
Ofte er for Filosofien Spørgsmaalene interessantere end
Svarene. Saadanne Grundtanker fødes og dør som andre
Tanker. Det gælder da, støttet til Videnskabernes Historie,
at vise, hvorledes baade deres Fødsel og deres Død betinges
ved de Krav, den stadig fortsatte Stræben efter Sandhed
stiller. I Tankelivet, som overalt, er det ikke de bestemte
Former, der bestaar, men de Kræfter, der fører til Formernes
Opstaaen og selv kan vedblive at virke under nye Former.

Den Analogi, der saaledes fremtræder mellem Videnskab
og Personlighed, idet det paa begge Omraader gælder Kamp
mod Kaos, har for mig været af stor Betydning. Det er ikke
fremmede Magter, vi giver os i Vold, naar vi vier vort Liv
til videnskabeligt Arbejde. Videnskaben er Slægtens Fællesarbejde
for at opnaa Orientering i Tilværelsen, ligesom det
er den enkelte Personligheds Opgave at orientere sig i sin
lille Kreds. Og selve Arbejdet i den lille Kreds kan ikke
øves uden Delagtighed i det store Fællesarbejde, i Tilegnelse
eller i Medvirken. En Afslutning paa dette Arbejde kan
ikke tænkes, saalænge Tilværelsen frembyder nye Sider
for os. Og det kunde jo ogsaa være, at Tilværelsen selv ikke
nogensinde er færdig en Gang for alle!

\asterism

Skarpest maatte Læren om Livet som en Kamp for Tilværelsen
gribe ind paa det etiske Omraade. Det er da ogsaa
% --- p. 13 ---
\setcounter{footnote}{0}%
\opage{13}blevet sagt, at den voldte „den største Krise, Etiken til
nogen Tid har haft at bestaa“. Gælder det i Livet overhovedet
om andet end om at skaffe sig Midler til Selvopholdelse?
--- Overfor et saadant Spørgsmaal har jeg især
gjort gældende, at Selvopholdelse altid betinges ved det Udviklingstrin,
hvorpaa man nu engang staar. Der er altid et
bestemt Niveau, man vil hævde, og som man ikke vil synke
ned under. At det gælder i økonomisk Henseende, viser
daglig Erfaring. Men det gælder ogsaa i etisk Henseende.
De aandelige Goder, man har Adgang til, vil man ikke udelukkes
fra; der kæmpes netop i vore Dage i nødlidende
Egne mangen haard Kamp af denne Grund. Det Trin, man
staar paa, vil betinge den Maade, paa hvilken Kampen
føres. Omvendt kan den Maade, paa hvilken Kampen føres,
vise, at man i Virkeligheden ikke stod paa saa højt et Trin,
som man mente. --- Men naar der nu gives saare mange forskellige
Trin, og ikke blot Handlinger, men ogsaa Domme
om Handlinger betinges ved dem, hvorledes er saa en almengyldig
Etik mulig? Her saa jeg ingen anden Mulighed
end at tage sit Stade paa et af disse Trin og undersøge,
hvilke etiske Domme der maa fældes ud fra det. Paa bar
Grund kan man ikke filosofere, lige saa lidt i Etiken som
paa andre Omraader. For mit Vedkommende fandt jeg i
den Humanitetsfølelse, der udviklede sig i den antike Oldtids
Slutning, uddybedes ved Kristendommen og i det 18de
Aarhundrede frigjordes fra dogmatiske Skranker, et Udgangspunkt,
det vilde være frugtbart at benytte, for saa
dernæst at sammenligne de Resultater, man kom til, med
de Resultater, man maatte komme til ud fra andre Synspunkter.
En forudsætningsløs Etik gives der ikke.

\asterism

Foruden det naturhistoriske Motiv, der gjorde sig gældende
for mig i de hidtil antydede Spørgsmaal, var der ogsaa
et kulturhistorisk Motiv, som satte min Eftertanke i
Bevægelse. Dette var især Tilfældet ved det religiøse Pro%
% --- p. 14 ---
\setcounter{footnote}{0}%
\opage{14}blem. Brodden i det religiøse Problem var for mig fremkommen
ved Arbejdets Deling paa det aandelige Livs Omraade.
Paa de mest primitive Trin, vi kender, staar Religion
som eneste Form for aandelig Kultur. Det forholder
sig paa lignende Maade, naar en stor Religion staar i sit
første Udspring, saaledes i Buddhismen og Kristendommen
i deres oprindelige Former. Da er Religion Et og Alt,
bestemmer alle Tanker og alle Opgaver. Senere hen forsøges
det maaske --- i størst Stil af Katolicismen i det 13de
Aarhundrede --- at forene al Tidens Kultur til et stort Hele,
saaledes at religiøse Idéer giver de afsluttende Synspunkter
og løser de sidste Gaader. Men paa intet af disse Trin
staar vi nu. Religion er bleven et enkelt Element i Aandslivet,
hverken det eneste eller det altbeherskende. Videnskab,
Kunst og Moral staar paa deres egne Ben, og Samfundslivet
udvikler sig efter sine egne Love. Religionen er
ikke længere Ildsøjlen, der leder Menneskehedens Vandring
gennem Ørkenen; den er snarere Ambulancen, som
følger bagefter og opsamler de Mødige og de Lidende. Den
er Trøsteren, ikke Vejviseren. Ogsaa dette er stort. Og der
er ingen anden aandelig Magt til Rede, som kan overtage
denne Opgave. Men det betager ikke Problemet dets Brod.

Om nogetsteds saa gælder det her, at hvad der bestaar
i Verden, den aandelige som den materielle, ikke er de Former,
under hvilke Livet til given Tid fremtræder, men de
virkende Kræfter, der har skabt de gamle Former og kan
skabe nye. For mit Vedkommende har jeg ment som Kærnen
i Religionen paa alle Trin at finde en Trang og Villie
til at fastholde Livsværdiernes Bestaaen ud over de Grænser,
indenfor hvilke menneskelig Villie kan virke for dem.
Man kan gennem Religionernes Historie aflæse, hvad Menneskene
til forskellige Tider har sat højest. Og den Trang,
der her har virket, kan fremdeles virke, selv efter at de
aandelige Vilkaar i væsentlig Grad har forandret sig.

Der gives, som jeg engang har udtalt, et Standpunkt, paa
hvilket der gennem Livets Blanding af Tragedie og Komedie,
af Sejr og Nederlag, er vundet en Grundstemning, som
% --- p. 15 ---
\setcounter{footnote}{0}%
\opage{15}paa een Gang har Karakteren af Alvor --- formedelst det
Store, der er oplevet, og den Smerte, det ofte har kostet
--- og af Ironi overfor ethvert Forsøg paa i Billeder eller
Begreber at ville udtrykke, hvad der i sin Fylde og ved
sine Modsætninger sprænger enhver Form, man forsøger
at slaa fast. Den Livsanskuelse eller Livsfølelse, som saaledes
opstaar, fastholder Forvisningen om en stor Sammenhæng
i Tilværelsen, hvad Værdier saavel som hvad andre
Emner og Elementer angaar, og kan derfor vende sig spøgende
mod Livets Tilskikkelser, fordi de staar magtesløse
overfor det Højeste og det Bedste, saavel som mod alle
Forsøg paa at udsige det Uudsigelige. --- Om man vil give
en saadan Tro Navn af Religion eller ikke, er ligegyldigt.
Vi mangler maaske et træffende Ord for, hvad der her
kan foreligge. Har vi i mange Tilfælde for mange Ord, saa
har vi her for faa. Hvad de gamle Religioner angaar, vil
deres blivende Værdi, fra et Standpunkt som det antydede,
væsentlig bestaa i den store Poesi, gennem hvilken de har
formaaet at udtrykke Livets Modsætninger, dets Højdepunkter
og dets Dybder. Det er da ogsaa gennem deres
Kultus og den Kunst, der er traadt i denne Kultus' Tjeneste,
at de endnu den Dag idag har saa stor Indflydelse paa
Menneskers Sind.

\piecerule

Jeg har her berørt nogle af de Tankemotiver, der har
gjort sig gældende i min Filosoferen. Om Filosofiens Betydning
har jeg aldrig tvivlet. Jeg har været i Tvivl om
Midler og Veje, endnu oftere om mine Evner til at arbejde
i Filosofiens Tjeneste og fyldestgøre de intellektuelle og
personlige Krav, Filosofien stiller til sine sande Dyrkere.
Men at der her er Opgaver og Arbejder, der stadig maa
tages op, derom har jeg aldrig tvivlet, saa højt der end
var fra mit Stade op til Filosofiens store Skikkelser.

Der staar paa mit Arbejdsbord et Fotografi af en Buste
af Platon, som en af vore Kunsthistorikere har foræret
mig. Naar jeg om Morgenen sætter mig til mit Arbejde
% --- p. 16 ---
\setcounter{footnote}{0}%
\opage{16}foran dette Billede, føler jeg mig ofte uværdig til at regne
mig til samme Lav som Platon. Platonismens Tid er forbi,
og alle „ismer“ vil forsvinde. Men den Tankedrift, der førte
Platon til, hele sit lange Liv igennem indtil han døde 80
Aar gammel (en passende Tid at dø i), i stedse fornyet
Stræben at søge tilbage til Erkendelsens og Livets Grundtanker,
den vil ikke dø. Og den vil kunne røre sig selv i
den beskedneste Virken, naar denne er forbunden med
Ærlighed og med intellektuel Redelighed.

\asterism

Da jeg for 53 Aar siden havde forsvaret min Afhandling
for Doktorgraden, sluttede jeg, efter Skik og Brug,
men ikke blot fordi det var Skik og Brug, med at udtale
de bedste Ønsker for Universitetet. Dermed vil jeg ogsaa
slutte iaften. Maatte dette Universitet, hidtil Danmarks eneste
videnskabelige Højskole, vedblive at virke gennem Tiderne
og bære Vidne om vort Folks Evne og Villie til at
yde Bidrag til Menneskeslægtens intellektuelle Kultur!

% --- p. 17 ---
\setcounter{footnote}{0}%

\opage{17}\essayhead{OM MIN FORFATTERVIRKSOMHED}

I min første Ungdom levede jeg mig ind i to Tankerækker,
der senere skulde stille mig de alvorligste Problemer.
Allerede i min Skoletid havde jeg tænkt mig Teologi som
mit fremtidige Studium. Men Læsningen af Platon's Symposion
i det sidste Skoleaar opfyldte mig med dyb Beundring
og bragte mig et Vidnesbyrd om, at der ogsaa paa
rent menneskelig Grund kunde opstaa høje og betydningsfulde
Tanker. I Løbet af mine Studenteraar fortsatte jeg
Studiet af Platon, og dertil kom saa Studiet af den tyske
Filosofi fra Kant til Hegel. Modsætningen mellem Teologi
og Filosofi blev af de filosofiske Professorer, der dengang
virkede ved Universitetet, Rasmus Nielsen og Brøchner, af
hver paa sin Vis, meget skarpt betonet. Jeg følte mig da
mere og mere delt. Jeg vandrede fra det teologiske Auditorium
til det filosofiske, og tilbage igen; de laa hvert i sin
Ende af den lange Universitetsgang. Det var en Tid fuld
af Pine, og jeg var ofte ved at fortvivle. Men den Forfatter,
jeg fordybede mig mest i, var Søren Kierkegaard, og ved
hans Hjælp øjnede jeg tilsidst en ny Mulighed. Hans vældige
Kamp mod Folkekirken, og overhovedet mod moderne
Kristendom, gav endnu Genlyd. Men det var dog
især hans filosofiske Skrifter, som fængslede mig. Jeg lærte
af ham, at religiøs Overbevisning tilsidst var personlig Erfarings
Sag og beroede paa personlig Trang, og derved blev
den Opgave stillet mig, at prøve, om den kristelige Tros
Indhold virkelig svarede til min egen inderste Trang. Svaret
kom først efter en lang Kamp. En Tidlang forsøgte jeg
at holde min personlige Tro og min filosofiske Interesse
helt ude fra hinanden. Kierkegaard betragtede jo al viden%
% --- p. 18 ---
\setcounter{footnote}{0}%
\opage{18}skabelig Erkendelse som „uvæsentlig“ i Forhold til den
religiøse Trang. Men denne „uvæsentlige“ Erkendelse fik
nu stedse mere „væsentlig“ Betydning for mig. Dog har
Kierkegaard ogsaa paa mere positiv Maade forberedt min
senere Udvikling. Han hævdede, at det aandelige Livs formale
Ejendommelighed bestod i den Enhed og Koncentration,
med hvilken Livsindholdet sammenholdes, og at Maalestokken
for et aandeligt Livs Værdi bestaar i Forholdet
mellem Indholdets Fylde og den Energi, hvormed Koncentrationen
øves. Denne Tanke, der i Filosofiens Historie især
repræsenteres af Leibniz og Kant, danner Grundlaget for
den Karakteristik, Kierkegaard har givet af de forskellige
Maader at tage Livet paa, hans saakaldte „Stadier“. Betydningen
heraf var mig naturligvis ikke fra først af klar.
Først ved et kritisk Studium af Kant og af Positivismen,
og især under Forarbejderne til min Fremstilling af Psykologien,
gik dens store teoretiske og praktiske Betydning
op for mig. Jeg følte mig nu stedse mere som Filosof, og
efterat jeg 1870 havde taget Doktorgraden med en Afhandling
om „Den antike Opfattelse af Menneskets Villie“, begyndte
jeg som Privatdocent at læse over Filosofi.

Efterat jeg havde læst over forskellige Afsnit af Filosofiens
Historie (især over „efterhegelsk Filosofi i Tyskland“
og over „den engelske Filosofi i vor Tid“) holdt jeg i Efteraaret
1874 en Forelæsning over „Indledning til Filosofien“,
i hvilken jeg i al Korthed udviklede den Opfattelse af de
filosofiske Problemer og deres indbyrdes Forhold, der havde
fæstnet sig hos mig. Og saa begyndte min egentlige Forfattervirksomhed,
idet jeg stillede mig den Opgave at behandle
hvert Problem for sig. Det korte Program, jeg havde
givet i den omtalte Forelæsning, maatte naturligvis under
Udarbejdelsen i det Enkelte og ved de nye Studier, jeg
foretog, paa mange Maader ændres. Der var et stort Arbejde
at gøre paa de forskellige Omraader. Naar jeg nu ser
tilbage, kan jeg inddele mit Arbejde i de følgende 40 Aar
i tre Perioder. Fra 1875 til 1887 stod Psykologi og Etik i
første Række; fra 1887 til 1895 var jeg især beskæftiget
% --- p. 19 ---
\setcounter{footnote}{0}%
\opage{19}med den nyere Filosofis Historie, og fra 1895 af har jeg
været optagen af Religionsfilosofi og Erkendelsesteori. Erkendelsesteori
blev mere og mere for mig den centrale Disciplin
i Filosofien. Filosofien er for mig en Stræben efter
at sammenfatte det menneskelige Tankearbejde under visse
almene Synspunkter, en Stræben, der forberedes indenfor
de enkelte Videnskaber, naar disse bliver sig deres ledende
Tanker bevidst.

Min „Psykologi paa Grundlag af Erfaring“ udkom i Efteraaret
1882. Bogens Hovedtanke, eller rettere den Hypotese,
den søger at begrunde og bekræfte er, at Sjælelivets
Grundlov er Syntesens, den sammenfattende Virkens Lov.
Begrundelsen af denne Hypotese finder jeg dels i den Maade,
paa hvilken Bevidstheden gaar over til Ubevidsthed,
og omvendt vaagner fra denne, dels i en Analyse af nogle
vigtige sjælelige Fænomener og Processer (Erindring, Sammenligning,
Sorg, Tvivl, Beslutning). Og Bekræftelsen finder
jeg ved Undersøgelse af de Betingelser, under hvilke
Sansefornemmelsen opstaar, saaledes som de belyses af
Fechner og den af ham grundlagte Experimentalpsykologi,
fremdeles ved Undersøgelse af Forestillingsforbindelse,
Fantasi og Følelseslivets og Villieslivets Hovedformer. Syntetisk
Energi, Villen i Ordets videste Betydning er for mig
det centrale Udtryk for Sjælelivet. Der gøres, saalænge
Sjælelivet bestaar, et Arbejde for at holde sammen paa et
mere eller mindre rigt Indhold. Følelser er Symptomer paa
denne Villen, og Erkendelsen er fra først af dens Tjener,
siden maaske dens Herre. I de senere Udgaver er det forsøgt
at gennemføre dette Synspunkt med større Konsekvens.

Grundtanken i Bogen har successivt udviklet sig hos mig,
fra min første Ungdoms Kierkegaardsstudier gennem mine
Kantstudier og gennem Studier over det syge Sjæleliv i
Sammenligning med det sunde. --- Til den franske Oversættelse
af Bogen har Pierre Janet skrevet en Fortale, og
da, som sagt, en Sammenligning mellem sygt og sundt Sjæ%
% --- p. 20 ---
\setcounter{footnote}{0}%
\opage{20}leliv havde været af stor Betydning for Udarbejdelsen af
min Hypotese, var det mig en særlig Glæde, at Bogen kunde
fremtræde paa et Hovedsprog under en fremragende
Psykopatologs Ægide.

Jeg ser meget godt Bogens Mangler (maaske kender jeg
dem ikke alle). Men jeg tror, at dens Hovedtanker er berettigede
og stedse vil fremkomme paany, maaske i Former,
vi nu ikke kan forestille os. Kilderne til psykologisk Indsigt
er jo saa mangfoldige (Selviagttagelse, Experimenter,
Fysiologi, Psykopatologi, Litteratur- og Religionshistorie
o. s. v.), at det allerede nu er vanskeligere at benytte dem
alle, end det var, da jeg i Aarene 1876---1882 udarbejdede
min Psykologi.

For min Virksomhed som Universitetslærer, især efterat
jeg 1883 var bleven Professor, var den psykologiske Betragtningsmaade,
jeg havde indøvet, af stor Betydning.
Overfor Tidens aandelige Bevægelser var der her et Standpunkt,
som kunde møde de forskelligste Retninger med
intellektuel Sympati, uden dog at foregribe de afsluttende
erkendelsesteoretiske og religionsfilosofiske Tanker. Min
gamle Lærer Sibbern (som jeg havde hørt i mit første Studenteraar,
men som jeg først senere ret fik Forstaaelse for)
udtalte, at ligesom der for Botanikeren ikke gives noget
Ukrudt, saaledes gives der for Psykologen intet aandeligt
Fænomen, som ikke kan være Genstand for Iagttagelse og
Interesse. Før vi fælder en Dom om et aandeligt Fænomens
intellektuelle eller etisk-religiøse Værdi, maa vi erkende og
forstaa det. Altfor ofte gaar man den omvendte Vej og
stempler som Dumhed eller Synd, hvad der dog har sin
naturlige Rod i Livets Sammenhæng og ved rigtig Behandling
maaske kan naa op i Sandhedens og Godhedens Verden.

Det næste Problem, der optog mig, var det etiske. Allerede
1876 havde jeg udgivet et lille Skrift derom (Om
Grundlaget for den humane Etik). Men det forekom mig
nu, at Forudsætningerne for en Etik ikke var tilstrækkelig
% --- p. 21 ---
\setcounter{footnote}{0}%
\opage{21}kritisk undersøgte her. I min Psykologi havde jeg gjort en
skarp Forskel mellem Kendsgerning og Værdi (Gyldighed).
Det etiske Problem stillede sig nu for mig paa følgende
Maade. Ligesom Logik og Matematik beror paa almenmenneskelige
Forudsætninger, og ligesom Erfaringsvidenskaben,
f. Ex. Psykologien, hviler paa Aarsagsloven, der er
alment anerkendt, gives der saa ogsaa alment anerkendte
Forudsætninger, paa hvilke Etiken kan bygge, og hvilke er
de? Det kostede mig ikke lidet Arbejde, før jeg kom til den
Overbevisning, at naar man stiller tilstrækkelig strenge Fordringer
til Begrundelsen, gives der ingen Etik, som paa
rationel Maade kan paatvinges Alle, saaledes som Tilfældet
er med logiske, matematiske og naturvidenskabelige Sætninger. Psykologisk
Iagttagelse og historisk Erfaring viser,
at menneskelige Handlinger bedømmes ud fra højst forskellige
Forudsætninger. Selv om etiske Standpunkter kan være
enige i deres Resultater, vil de dog ofte gaa forskellige Veje
i deres Begrundelse. Ved logisk Analyse kan vi komme tilbage
til de for et enkelt bestemt etisk Standpunkt ejendommelige
Forudsætninger, men selve disse Forudsætninger
hænger sammen med psykologiske og historiske Forhold,
der ikke paa ganske samme Maade er tilstede paa
andre Standpunkter. Jeg fandt visse Hovedtyper, ved hvilke
etiske Dommes Forudsætninger kan bestemmes --- fra Øjebliksstandpunktet,
Individualitetsstandpunktet, Familie-,
Klasse- og Nationalitetsstandpunktet til det universelle Humanitetsstandpunkt.
Det var mig umuligt at udlede disse
forskellige Standpunkter af „Menneskets Idé“ og at tilskrive
noget af dem enestaaende Nødvendighed.

Etikens Stilling kan til en vis Grad sammenlignes med
Astronomiens Stilling før Newton. Da var ogsaa i Astronomien
forskellige Systemer mulige. Jeg véd ikke, om Etiken
vil faa sin Newton; hidtil er han i hvert Tilfælde ikke kommen.
Ethvert Forsøg paa en Etik maa derfor være klart
over sine psykologiske og historiske Forudsætninger. Hvad
mig selv angaar, fandt jeg et betydningsfuldt Grundlag for
en Etik i den universelle Sympati, den almene Humanitets%
% --- p. 22 ---
\setcounter{footnote}{0}%
\opage{22}følelse, der udviklede sig i den græsk-romerske Verden,
især efter Alexanders Tid, og saa, fordybet og udvidet ved
den kristelige Menneskekærlighed, dukkede op igen i det
18. Aarhundredes Idéer. Min Opgave var nu at analysere og
uddybe dette Grundlag og ud fra det drage Slutninger med
Hensyn til menneskelige Handlemaaders og Institutioners
Værdi. Hvad der kunde begrundes ud fra dette Standpunkt
maatte saa sammenlignes med de Slutninger, der kunde
drages ud fra andre Standpunkter (Individets Selvopholdelse,
Familie, Nation o. s. v.). „Menneskets Idé“ udtrykker
sig nu engang for os i forskellige, dels samtidige dels successivt
optrædende Typer. Her foreligger, hvad Fichte ud
fra sit idealistiske Synspunkt har kaldt en projectio irrationalis. De specielle Vurderinger bestemmes tilsidst ved et
Grundmotiv, en Grundværdi. Man kan opstille en borgerlig
eller social Moral uden at spørge om Motiverne, men det
Etiske angaar nu engang ikke blot Handling, men ogsaa
Sindelag. Og hvad Motiverne angaar, maa der skelnes mellem
Vurderingsmotiv (Grundværdi) og Handlingsmotiv. De
falder ikke altid sammen. Saaledes kan Menneskeligheden
som Vurderingsmotiv afføde specielle Handlingsmotiver
som Retsfølelse, Tapperhed, Tjenstvillighed o. s. v. Vurderingsmotivet
begrunder tilsidst den Værdi, der tillægges de
forskellige Handlingsmotiver. Hvad der i ét etisk System
er Handlingsmotiv, er maaske i et andet System Vurderingsmotiv,
og omvendt.

Den specielle Udarbejdelse af min Etik foretoges paa
Grundlag af et psykologisk, historisk og nationaløkonomisk
Materiale, som jeg i Aarenes Løb havde samlet. Saa meget
som muligt gik jeg ind paa de specielle menneskelige Livsomraader.

Efter min Overbevisning behøves der hverken nogen
overnaturlig eller nogen metafysisk Sanktion af Etiken.
Den menneskelige Natur, saaledes som vi kender den fra
Erfaringen, frembyder tilstrækkelige Vurderings- og Handlingsmotiver.
Min Etik, som udkom 1887, blev dels angreben
fra teologisk Side, dels fra visse antiteologiske Sider.
Der er Antiteologer, som mødes i broderlig Forening med
% --- p. 23 ---
\setcounter{footnote}{0}%
\opage{23}Teologerne i den Paastand, at en Etik uden overnaturlig
Sanktion og uden metafysisk Grundlag er umulig. Overfor
begge Parter hævder jeg, at selv naar man har opgivet den
gamle Tro, er alvorligt og positivt Arbejde for etisk Overbevisning
og Livsførelse ikke blot muligt, men nødvendigt.

--- Det er nu 35\footnote[1]{1927 er det 40 Aar siden,} Aar siden min Etik udkom, og dens væsentligste
Tanker holder jeg endnu fast. Bogen finder ogsaa
stadig Læsere baade i Udlandet og herhjemme. --- Kun paa
enkelte Punkter har jeg i senere Udgaver gjort Ændringer
og Tilføjelser.

Efter Fuldendelsen af min Psykologi og min Etik stillede
jeg mig den Opgave at give en Fremstilling af den nyere
Filosofis Historie. Denne Opgave havde en naturlig Sammenhæng
med den Plan om at foretage en Undersøgelse af
de filosofiske Hovedproblemer i deres indbyrdes Forhold,
som jeg allerede havde fattet 1874, og som stærkt beskæftigede
mig.

Som jeg ovenfor har udtalt, spillede græsk Tænken og
Livsanskuelse en stor Rolle i min personlige Udvikling. Jeg
havde ved den vundet en dyb Overbevisning om Videnskabens
og det humane Livs selvstændige Betydning. Og i min
Etik havde jeg paastaaet, at vor Etik er den græske Etik,
uddybet og udvidet ved de følgende Tiders Aandsliv. Overhovedet
ligger græsk Tænken ved Roden af al europæisk
Filosofi, og ethvert dybere gaaende Studium maa stedse
paany gaa tilbage til den. Jeg har holdt Laboratorieøvelser
over Platon's „Stat“ og over Aristoteles' „Etik“. Men jeg
maatte begrænse mig, og for Problemernes Stilling i vor Tid
var naturligvis de sidste tre Aarhundreders Filosofi af største
Betydning.

I min Bearbejdelse af den nyere Filosofis Historie lagde
jeg større Vægt paa Problemerne end paa Systemerne, og
jeg søgte mere at karakterisere end at referere. Gennem
Karakteristiken kommer man paa naturlig Maade over i
Kritiken. Til Karakteristiken hører ogsaa en Undersøgelse
af Forholdet mellem Filosofens Personlighed og hans Lære,
% --- p. 24 ---
\setcounter{footnote}{0}%
\opage{24}og flere af de smaa Biografier i min Bog har jeg ret skrevet
con amore. Desuden har jeg lagt Vægt paa Filosofiens Forhold
til andre Videnskaber, og derfor indtager i min Fremstilling
Mænd som Kopernikus, Galilei, Newton, Darwin og
Robert Mayer en fremragende Plads, idet deres Forskninger
foranlediger betydningsfulde Ændringer i de filosofiske
Problemstillinger. --- Bogen udkom 1894---1895. Paa enkelte
Punkter havde jeg tidligere gjort monografiske Forarbejder
(om Kant's filosofiske Udviklingsgang og om Kierkegaard);
andre Punkter har jeg senere taget op til monografisk
Bearbejdelse (Rousseau, Spinoza). I Sammenligning
med de fleste tyske Fremstillinger har jeg lagt større Vægt
paa fransk og engelsk Filosofi i det 19. Aarhundrede. Jeg
havde her et Forarbejde i min 1874 udkomne Bog „Den
engelske Filosofi i vor Tid“. --- Senere har jeg skrevet om
„Moderne Filosofer“ (1903) og om „Henri Bergson's Filosofi“
(1914)\footnote[1]{Senere har jeg (se nedenfor) skrevet om Platon og Spinoza.} ---

Efter Udgivelsen af mit historiske Arbejde koncentrerede
mine Studier sig om Erkendelsesteori og Religionsfilosofi.
Et af Forarbejderne til min Erkendelsesteori ligger i en
Afhandling om „det psykologiske Grundlag for logiske
Domme“ (1898). Men foreløbig var jeg især optaget af det
religiøse Problem, efterat det en Tidlang ikke havde interesseret
mig. Jeg havde ikke tænkt mig at foretage en speciel
Undersøgelse af det, ud over de Antydninger, der var givne
i min Psykologi og i min Etik. Men nu var Læsning af Biografier
bleven min Yndlingslæsning, og jeg kom derved ogsaa
til at beskæftige mig med Augustin's „Confessiones“.
De interesserede mig overordentlig, baade psykologisk og
religionshistorisk. Og da jeg nu fordybede mig i tiende Bog,
hvor den store Kirkelærer opkaster det Spørgsmaal: Qvid
est quod amo, quum deum amo? (Hvad er det, jeg elsker,
naar jeg elsker Gud?) --- blev det mig klart, at der laa en
vigtig Kilde til Forstaaelse af religiøse Fænomener i Selvvidnesbyrd
af Personligheder, der, alt imedens de lever helt
% --- p. 25 ---
\setcounter{footnote}{0}%
\opage{25}og fuldt i deres Religion, har Trang og Evne til, i Eftertankens
Form at klare sig, hvad de oplever i den religiøse
Grebetheds Tider. Jeg samlede da en Række af lignende
Vidnesbyrd --- fra indisk Literatur, fra Middelalderens Mystik, fra Luther og fra Pascal --- og analyserede dem psykologisk.
Jeg var naturligvis klar over, at man ikke ad denne
Vej kan løse det historiske Problem om Religionens Opstaaen.
Her er jeg ofte bleven misforstaaet, især af dem,
der mener, at naar de har „gendrevet“ Hottentotternes Religion,
har de „gendrevet“ al Religion. Overhovedet lader
Spørgsmaalet om Religionernes første Opstaaen sig ligesaalidt
direkte besvare som Spørgsmaalet om Sprogets første
Opstaaen. Men ad den Vej, jeg slog ind paa, kunde man
maaske opdage, hvad der holder Religion oppe, --- hvad
der, naar den engang er bleven til, giver den dens Magt og
Betydning. Og dette er for Filosofen, der vil undersøge Problemets
nærværende Stilling, vigtigere end Spørgsmaalet
om Kultusformers og Dogmers første Opstaaen. Jeg saa
naturligvis ogsaa, at Studiet af Religionernes subjektive og
individuelle Side maatte forbindes med et Studium af Religionshistorien.
Og her søgte jeg især at efterspore, hvilken
Maalestok Religionshistorikerne anlægger, naar de kalder
en Religion højere eller lavere end en anden. Ad begge Veje
søgte jeg at danne mig en begrundet Formodning, en Hypotese
om Religionens Grundlag i dens forskellige Former.
Og dermed mente jeg ogsaa at kunne finde Grundlaget for
en Forstaaelse af Religionens Stilling i vore Dage og for
Formodningen om dens Fremtid. --- I en Tale ved en Universitetsfest
1897 gav jeg en foreløbig Karakteristik af mit
Standpunkt, og i en Forelæsning, som jeg i to Semestre
(1899---1900) holdt for Studerende af alle Fakulteter, gav
jeg en udførlig Fremstilling af min Opfattelse. I det følgende
Aar (1901) udkom saa min Religionsfilosofi.

Enhver kirkelig Teologi har et Indbegreb af religiøse
Forestillinger, som den tillægger bogstavelig Gyldighed. Indenfor
en og samme Religion kan der dog være forskellige
Meninger om, hvor Grænsen skal drages mellem det Billed%
% --- p. 26 ---
\setcounter{footnote}{0}%
\opage{26}lige og det Bogstavelige. Men Filosofien maa underkaste
dette Spørgsmaal en principiel Undersøgelse, idet den sammenligner
religiøse Forestillinger med de Forestillinger, der
dannes paa videnskabelig Grund og efter videnskabelig Metode.
Den første Undersøgelse indenfor Religionsfilosofien
vil derfor være erkendelsesteoretisk. En Analyse af de Hovedbegreber,
med hvilke Religionen opererer, førte mig til
det Resultat, at disse Begreber er dannede ved Analogi, og
det en Analogi, der ikke --- som de Analogier, der i Videnskaben
tjener som Arbejdsanvisninger --- kan finde Bekræftelse
ved nøje Iagttagelse. Religionens Sprog er et Billedsprog,
der stammer fra en Tid, da man ingen Forskel gjorde
mellem Billede og Tanke. En saadan Forskel kan først gøres,
naar der har udviklet sig en strengt prøvende Videnskab.
Det er ikke saa meget Videnskabens Resultater som
den videnskabelige Metode, den videnskabelige Maade at
stille Spørgsmaal paa og at danne Begreber paa, som her
er af Betydning. Religionens Spørgsmaal kan hverken stilles
eller besvares af Videnskaben, og Religionens Svar
er ikke Svar paa Videnskabens Spørgsmaal.

Det er altsaa forskellige Interesser, der fremtræder i
Religion og Videnskab. Hvad er da nu den Interesse, vi
kalder den religiøse? Dette Spørgsmaal forsøger Religionsfilosofiens
psykologiske Del at besvare. Min Bestræbelse
gaar ud paa at betragte Religionen som et psykologisk
Fænomen. Opgaven er at udfinde, i hvilken Egn af Sjælelivet
Religionen har sine Rødder, --- hvilken Trang den kan
tilfredsstille, --- og saa at undersøge Betingelser for denne
Trangs fortsatte Bestaaen. Det er her ikke nok at godtgøre,
at samme Dogmer læres og samme Kultus øves. Man maa
undersøge, hvorledes disse Dogmer og denne Kultus virker
paa de forskellige Individualiteter, og især, hvorpaa det beror,
at Mennesker atter og atter vender tilbage til dem. Det
viser sig da, at baade Dogmer og Kultusformer hænger
meget nøje sammen med, hvad der for Menneskene paa det
givne Trin er det højeste Værdifulde, Genstanden for deres
Ønske og for deres Stræben. Her melder Begrebet Grund%
% --- p. 27 ---
% pp. 27-28 are MISSING from the NB scan (the leaf was skipped; pp. 26 and 29 were shot twice, PDF 30-33).  The text below,
% from „Grund-|værdi“ to „jeg har“, is taken from the SAGA / Lindhardt og Ringhof e-book (2020, ISBN 9788726363050) of the
% same 1927 text, NOT read at a print; the e-book's own OCR faults are unchecked here, and the turn to p. 28 is not known
% (Hans, 2026-10-17; listed in KB-VISIT-TODO.md for a reading at a print).
\setcounter{footnote}{0}%
\opage{27}værdi sig igen, og min religionsfilosofiske Hypotese gaar da ud paa, at
Religionens Væsen har sin Grund i en Trang til eller et Ønske om, at se
Grundværdiens Bestaaen sikret, ogsaa ud over de Grænser, indenfor
hvilke menneskeligt Arbejde kan virke for denne Bestaaen. Naar det nu
kommer dertil, at de religiøse Forestillinger ikke kan fastholdes i
deres bogstavelige Betydning, saa maa Trangen til at fastholde det
Værdifuldes Bestaaen søge Udtryk i Symboler, i Livspoesi, d. v. s. i en
af Livet selv fremvoxende Poesi. Idealet vilde være, at enhver enkelt
frembragte sine egne Symboler. Men det vil ogsaa fra den Enkeltes
Standpunkt være berettiget, ved uvilkaarlig Tydning at genfinde de
Tanker, hans Livserfaring har ført ham til, i de overleverede religiøse
Forestillinger. Langt flere end de, der er sig det klart bevidst, har
gjort Overgangen fra Dogmetro til Livspoesi. Man tager af den religiøse
Overlevering, hvad man behøver og lader det øvrige ligge. Dette gør
enhver Kirke og ethvert religiøst Parti. Ad den Vej er moderne
Kristendom bleven til en anden Religion end Urkristendommen, idet den
faktisk har ladet væsentlige Forudsætninger, der afgørende bestemmer de
ældste Kristnes Etik, falde, for at lempe sig efter de nye historiske
Forhold.

Det er for mig stedse mere en rent psykologisk Interesse, der fører mig
til Beskæftigelse med det religiøse Problem. Jeg finder dér, hvor man
kan vinde et Indblik i dybt religiøse Personligheders indre Liv, en
Rigdom i Stemningen, en Storhed i Fantasien, en Lidenskab i Villien,
der ikke blot er en tiltrækkende og betagende Genstand for Psykologen
som saadan, men ogsaa foranlediger ham til at spørge, om Kræfter som
disse kan forsvinde af det menneskelige Aandsliv, uden at nye sjælelige
Kræfter danner en virkelig Erstatning for dem. Ideen om psykiske
Ækvivalenter, der allerede dukker op hos Platon, fører til at søge nye
Kilder til Begejstring og indre Grebethed. Og denne Søgen er egentlig
selv religiøs, da den motiveres ved en Trang til at se Værdierne
bevarede.

Undertiden har jeg fremsat det religiøse Problem paa følgende Maade. Er
Religion et Organ eller er den en Svulst? Hvis den er et Organ, saa kan
det være livsfarligt at operere den bort, selv om den skulde høre til
de 107 overflødige Organer, vi ifølge Metchnikoff endnu besidder.
Spørgsmaalet kan kun være, om den vedblivende kan faa Næring, og om
dens Funktioner kan bevares gennem Øvelse, eller om den dør bort af
Mangel paa begge Dele. Hvis den er en Svulst, vil et Forsøg paa at
operere den bort kunne foranledige et saadant Blodtab, at Livet staar i
Fare; sikrere vilde det være at underbinde den, saa at den kan dø af
Mangel paa Næring. --- I begge Tilfælde er det fremtidig Erfaring, der
maa bringe Afgørelsen.

Denne sidste Betragtning fører over til den etiske Del af
Religionsfilosofien, som dels undersøger Muligheden af en Konflikt
mellem Etik og Religion (nemlig hvis Troen paa Værdiernes Bestaaen
afholder fra at arbejde for at frembringe og bevare Værdier), dels
drøfter det Spørgsmaal, hvorvidt Troen paa Værdiens Bestaaen kan yde
Styrke og Magt til etisk Arbejde. ---

Der gives ikke andre religionsfilosofiske Synspunkter end de tre
angivne, det erkendelsesteoretiske, det psykologiske og det etiske. En
særlig religionsfilosofisk Metode gives der ikke. ---

I en Afhandling, jeg har kaldt „Forord og Efterskrift til min
Religionsfilosofi“, og som findes i andet Bind af mine „Mindre
Arbejder“, har jeg forsøgt at belyse min Stilling til det religiøse
Problem som Filosof og som Menneske. Ogsaa i „den store Humor“, der
siden skal omtales, belyses denne Stilling. Som Filosof har jeg kun at
gøre med Videnskabens og Livets almene Principper; som Menneske har jeg
at gøre med den specielle Nuanceren og Anvendelse af disse Principper,
hvortil min Livsgerning og mine personlige Erfaringer har ført mig. Jeg
tror, at Filosofen og Mennesket hos mig er traadt hinanden nærmere i
Løbet af Livet, og at Ingen af dem har taget Skade deraf. --- Jeg
henviser særlig til Afsnittet „Tragik og Humor“ i Humorbogen, som det
klareste og dybeste Udtryk, jeg har
% --- p. 29 ---
\setcounter{footnote}{0}%
\opage{29}kunnet give for det Omraade, hvor min Filosofi og min
personlige Livserfaring møder hinanden.

Erkendelsesproblemet havde indirekte gjort sig gældende
i alle mine Skrifter. I Psykologien havde jeg meget skarpt
markeret Forskellen mellem det psykologiske og det erkendelsesteoretiske
Synspunkt. I Etiken havde den strenge
Fordring om videnskabelig Begrundelse ført mig til Tvivl
om en rationel Etik. I Religionsfilosofien laa erkendelsesteoretiske
Synspunkter til Grund for Benægtelsen af Religionens
intellektuelle Betydning. I Indledningen til „Den
nyere Filosofis Historie“ havde jeg opstillet Erkendelsesproblemet
som det centrale Problem overfor det psykologiske,
det etisk-religiøse og det kosmologiske (metafysiske)
Problem. I et Universitetsprogram (1902) gav jeg en foreløbig
Undersøgelse af de filosofiske Hovedproblemer i deres
indbyrdes Forhold. Efter disse og andre Forarbejder udkom
saa 1910 „Den menneskelige Tanke, dens Former og
dens Opgaver“.

Bogen opsøger først Tankelivet i dets uvilkaarlige Former.
Allerede i det uvilkaarlige Tankearbejde kan de Tankeformer
spores, der i den videnskabelige Tænken uddannes
videre, specialiseres og gennemføres med større Nøjagtighed
og Konsekvens. Efterat der ved Undersøgelser af
de uvilkaarlige Tankefunktioner har vist sig Mulighed for
en Erkendelse, gælder det at finde de bestemte Former,
ved hvilke virkelig Videnskab kan vindes. Mulighedens
Rige er større end Virkelighedens. Her gælder Wallenstein's
Ord: „Eng ist die Welt, und das Gehirn ist weit“. Det er
Videnskabernes Historie, som viser os de Tankeformer, i
hvilke sand Erkendelse kan vindes. Disse Former, de saakaldte
Kategorier, er ikke de samme til alle Tider. Det var
Kant's skæbnesvangre Vildfarelse, at han troede én Gang
for alle at kunne opstille et Kategorisystem. Kategorierne
opstaar sporadisk, men kan saa bagefter samles og ordnes,
f. Ex. saaledes, at der skrides frem fra de mest omfattende
til de speciellere. Jeg begynder min Kategorilære med Be%
% --- p. 30 ---
\setcounter{footnote}{0}%
\opage{30}grebet Syntese og det korrelate Begreb Relation; disse
to Begreber aabner Rækken af de fundamentale Kategorier,
der lægger sig for Dagen i al Tænken, den førvidenskabelige
saavel som den videnskabelige. Derfra gaas over til de
formale Kategorier (Identitet; Grund og Følge o. a.), der
ligger til Grund i Logik og Matematik, og derfra til de
reale Kategorier (Aarsag, Totalitet, Udvikling), som ligger
til Grund i al Erfaringsvidenskab, og derfra igen til de
ideale Kategorier (Værdibegreberne), der lægger sig for
Dagen i Æstetik, Etik og Religionsfilosofi. Overgangen fra
de reale Kategorier til de ideale finder jeg deri, at enhver
Vurdering forudsætter et Forhold til en Totalitet (hvad
enten denne udgøres af et enkelt Øjeblik eller et helt Personlighedsliv,
af et enkelt Individ eller et større eller mindre
Samfund). --- Ligesom Kategorierne opstaar, kan de
ogsaa dø; saaledes betragter jeg f. Ex. „Substans“ som en
hendøende Kategori.

Kategorilæren kan kaldes en deskriptiv Erkendelsesteori
og staar som saadan mellem Psykologien og Videnskabshistorien
paa den ene Side, Drøftelsen af de filosofiske Problemer
paa den anden Side. Problemerne opstaar ved Kategoriernes
Forhold til de faktisk givne Emner, der ikke er
blotte Produkter af visse intellektuelle Funktioner. Tydeligst
viser dette sig ved Erkendelsesproblemet. Kan vore
Erkendelsesformer afledes af de givne Emner, eller har
de, saa meget de end betinges og specialiseres ved Emnerne,
deres sidste Grund i den menneskelige Aands Struktur
(som dog ikke behøver at være uforanderlig)? Her
vil, som jeg forsøger at vise, den af Kant grundlagte Kriticisme
stedse hævde en Førsteplads i den erkendelsesteoretiske
Diskussion, fordi den kan anvende den analytiske
eller regressive Metode til Udfindelse af Erkendelsens Former og Principer. En saadan Undersøgelse ligger baade til
Grund for Dogmatisme, Empirisme, Pragmatisme og Postulatsteori.
Kun Analyse af vor Erkendelse og dens Betingelser
kan lære os, hvad der skal opstilles som Dogme,
som Erfaring, som Økonomi eller som Postulat. 

% --- p. 31 ---
\setcounter{footnote}{0}%
\opage{31}I et af sine Breve omtaler William James mig som
„god Pragmatist og Pluralist“. Det beror paa en Misforstaaelse.
Pragmatisk („økonomisk“) Betragtning af Videnskaben
har for mig kun historisk Interesse, nemlig som
Belysning af de Motiver, der har ført til Videnskab, ikke
som Begrundelse af Videnskabens Gyldighed. Og blandt
disse Motiver giver jeg den rent intellektuelle Interesse den
første Plads, medens de fleste Pragmatister ikke udtaler sig
tilstrækkelig tydelig paa dette Punkt. Og Pluralismen er for
mig en Modstand, der skal overvindes; den stiller en Opgave;
men er ingen Løsning og gør (som i vor Literatur
Herbert Iversen klart saa) al Løsning, endog al Stræben
efter Løsning, umulig. --- James, som jeg fra mit Ophold
i Amerika erindrer med hjerteligt Venskab, har sluttet for
meget af den Enighed, med hvilken vi begge stod overfor
det Absolutes Metafysik. Dog havde jeg i Amerika to Gange
Lejlighed til i James' Nærværelse at udvikle mine filosofiske
Synspunkter. De to Foredrag sammenarbejdedes, efter
James' Forslag, til en Afhandling i et amerikansk filosofisk
Tidsskrift. De er trykt paa Dansk i tredie Bind af
mine „Mindre Arbejder“ under Titel „En filosofisk Bekendelse“.
---

Det kosmologiske (metafysiske) Problem opstaar, naar
det Spørgsmaal opkastes, om der kan vindes en videnskabelig
Verdensanskuelse. Vanskeligheden herved er den, at
hele Tilværelsen aldrig er given som Emne for vor Tænken.
Stedse er det en enkelt Del eller Side af Tilværelsen,
vi lægger til Grund; „hele Tilværelsen“ opfattes saa i Analogi
med det saaledes til Grund lagte „Urfænomen“. Verdensanskuelsernes
Historie er saaledes Urfænomenernes
Historie. Indenfor vor Erfaring kan Analogi være et godt
Erkendelsesmiddel, idet vi ved dens Hjælp belyser det ene
Erfaringsomraade ved Hjælp af det andet. Men Tilværelsen
som Helhed ligger ikke afsluttet for os, og selv om dette
kunde tænkes, vilde vi dog ikke kunne karakterisere den
gennem Sammenligning med andre Emner. I Valget af Urfænomen,
af den Del eller Side af Tilværelsen, fra hvilken
% --- p. 32 ---
\setcounter{footnote}{0}%
\opage{32}man slutter til Helheden, og i den Maade, paa hvilken
Analogien anlægges og gennemføres, kundgør sig uvilkaarlig
Tænkerens individuelle Personlighed, og ethvert Forsøg
paa en Verdensanskuelse vil derfor ligge paa Grænsen
mellem Videnskab og Kunst. Det kosmologiske Problem
opfattes da ogsaa mere og mere som et Grænseproblem, og
Behandlingen af det vil kun have filosofisk Værdi, naar
den forberedes ved grundigt erkendelsesteoretisk Arbejde.
Men der vil sandsynligvis altid gives Aander, hvis intellektuelle
Expansionstrang hverken kan tilfredsstilles ved
religiøse Billeder eller ved videnskabelige Begreber.

I Afsnittet om Vurderingsproblemet giver „Den menneskelige
Tanke“ dels en rent principiel, dels en udførligere
Fremstilling af de Synspunkter, der ligger til Grund for
min Etik og min Religionsfilosofi. ---

Naar jeg i al Korthed skal angive min Plads overfor de
filosofiske Retninger, maa jeg først lægge Vægt paa, at Erkendelsesproblemet
stedse mere og mere er blevet det centrale
Problem for mig. Egentlig forfølger jeg i min Behandling
af andre Problemer tilsidst dette Problem i dets forskellige
Forgreninger. Kant's store Spørgsmaal om Muligheden
af Erkendelse, særlig af Erfaring, har ogsaa været
afgørende for mig. Jeg har lært meget af Positivismen,
baade i dens franske og i dens engelske Form, og min Opfattelse
kunde kaldes en kritisk Positivisme, fordi jeg lægger
stor Vægt paa Erfaringens Indhold, alt imedens jeg
ser et Problem i selve Begrebet Erfaring. Og jeg har lært
meget af Spinoza, men min Enhedslære, min Hævden af
Verdens Enhed og Lovmæssighed er ikke dogmatisk som
Spinoza's; den er en kritisk Monisme. --- Men alle saadanne
-ismer har meget lidt at sige.

Jeg har forsøgt at filosofere saaledes, at jeg fastholdt
Filosofi som et selvstændigt videnskabeligt Omraade. Den
virker for Klarhed og Selvbesindelse i den aandelige Verden,
med aaben Sans for alle Sider af Livet, i den Overbevisning,
at kun klar Tænken kan lade Livets Fylde komme
til dens Ret, skønt denne Tænken selv ikke kan træde
% --- p. 33 ---
\setcounter{footnote}{0}%
\opage{33}istedetfor den Livsfylde, indenfor hvilken den søger at
orientere sig. Jeg har en inderlig Overbevisning om Filosofiens
Betydning som et uundværligt og stadig virkende
Element i det aandelige Liv.

I mit 72. Aar tog jeg min Afsked fra Universitetet og
kunde nu ganske hellige mig Udarbejdelsen af forskellige
Ideer og Udkast, der endnu ikke havde antaget bestemte
Former. Kort forud havde Videnskabernes Selskab valgt
mig til Beboer af en „Æresbolig“, som Grundlæggeren af
Carlsbergfondet havde stillet til Selskabets Disposition. I
en skøn Bolig med en stor Have fandt jeg et godt Arbejdssted.
I de Aar, jeg har boet her, er det tre Arter af Arbejder,
der har lagt Beslag paa mig: psykologiske, erkendelsesteoretiske
og filosofihistoriske.

Det havde længe været mit Ønske at udføre et Arbejde
paa den beskrivende og analyserende Psykologis Omraade,
hvor der efter min Opfattelse endnu var meget at gøre.
Den experimentale Psykologi havde efter sin Natur hovedsagelig
beskæftiget sig med Sjælelivets elementære Fænomener.
Men det er af en stor Betydning for det aandelige
Liv, at der ogsaa indenfor de højere, mere komplicerede
og totaliserede sjælelige Tilstande, i hvilke netop vort Sjælelivs
inderligste Tendenser fremtræder, eftervises klar og
bestemt Sammenhæng. Og det gælder ikke blot om at finde
Sammenhæng. Det er ogsaa vigtigt at karakterisere de
Forskelligheder og Modsætninger, der her kan gøre sig
gældende. Jeg havde særlig længe beskæftiget mig med Begrebet
Humor. I min Psykologi havde jeg ganske kort beskrevet
Humor som en særlig Form af Følelsen ved det
Latterlige, og jeg havde fremstillet Haan og Humor som
de store Modsætninger paa dette Omraade. De fleste Latterfænomener
ligger mellem disse to Modsætninger. Nu ønskede
jeg at undersøge Humor nærmere, og jeg kom nu til
at betragte den som en af de Følelser, i hvilke Livet kan
ligesom opsummere sig for den Enkelte, --- som en Art
højere Livsfølelse, der er bestemt ved Forholdet mellem
% --- p. 34 ---
\setcounter{footnote}{0}%
\opage{34}et alvorligt Grundlag (der kan være af forskellig Beskaffenhed)
og Livets mangehaande Onder, Smaating og Urimeligheder.
I daglig Sprogbrug forstaar man ved Humor
alslags mere eller mindre vittig Spøg; det er, hvad jeg
kalder „den lille Humor“. Men jeg vilde beskrive „den
store Humor“, der forudsætter et dybt Erfaringsgrundlag i
Følelseslivet, medens Forgrunden finder Udtryk i Latter,
eller dog i Smil. Det var tillige min Opgave at bestemme
Forholdet mellem Humor som Livsfølelse og den religiøse
Følelse. Her havde Kierkegaard i sin ivrige Stræben efter
at hævde skarpe Forskelle stillet Humoren paa en altfor
lav Plads. Det Afsnit, i hvilket jeg behandler dette Punkt,
nemlig Kapitlet om „Tragik og Humor“, betragter jeg selv
som det vigtigste i Bogen og tillige som det klassiske Udtryk
for min egen Totalfølelse overfor Livet. --- Metoden
i denne Bog er, ligesom i min Religionsfilosofi, psykologisk-historisk.
Jeg fandt bl. a. her en velkommen Lejlighed til
at udtale min Beundring for to af de største Skikkelser,
jeg kender, den ene tilhører Tankens Verden, den anden
Poesiens Verden: Sokrates og Shakespeare. --- I Bogens
sidste Afsnit („Humor og Filosofi“) mener jeg at have henledet
Opmærksomheden paa en vigtig Sag, nemlig Forholdet
mellem en Forskers Sindstilstand, medens han er optagen
af intellektuelt Arbejde, og medens han bevæger sig
i det daglige Livs Forhold. Man kan ikke uden videre
slutte fra en Forskers videnskabelige Karakter til hans
Sindstilstand under det daglige Livs smaa og store Begivenheder.

Et andet psykologisk Arbejde udgav jeg 1918 under Titel
af „Oplevelse og Tydning“. Min Ven Théodore Flournoy i
Genève havde nogle Aar iforvejen udgivet nogle meget interessante
Optegnelser af en Dame, der havde spurgt ham
til Raads angaaende sin sjælelige Tilstand. Der kundgjorde
sig i hendes indre Liv en skarp Modsætning mellem Oplevelser
af den mest ophøjede og ideelle Art og andre Oplevelser
af den mest sanselig-lidenskabelige Art. Ofte antog
% --- p. 35 ---
\setcounter{footnote}{0}%
\opage{35}hine første Oplevelser Karakteren af Ekstaser. I sine sunde
Øjeblikke stod den psykologisk forstaaende Dame frit og
klart overfor disse højest forskellige Oplevelser. Navnlig
havde hun ikke den hos Mystikere, især i tidligere Tider,
hyppige Tilbøjelighed til at betragte ekstatiske Oplevelser
som Aabenbaringer af Bibelens og den kirkelige Traditions
guddommelige Skikkelser, og heller ikke i selve Visionerne
vævede saadanne Forestillinger sig ind. I denne Henseende
dannede hun en interessant Modsætning til den spanske
Nonne Sta Teresa (16. Aarh.), med hvem jeg under Forarbejderne
til min Religionsfilosofi havde beskæftiget mig,
og i hvis Skrifter jeg nu paany fordybede mig. Min Afhandling
søger at vise, hvor forskellig Oplevelser, der rent
psykologisk er af samme Art, tydes under forskellige historiske
Forudsætninger og af forskellige Personligheder. Jeg
fandt her et Exempel til Belysning af Grundtanken i min
Religionsfilosofi. Og tillige fandt jeg her en Analogi af erkendelsesteoretisk
Interesse. Ti ogsaa Videnskaben er en
Tydning af givne kvalitative Oplevelser, og heller ikke her
maa Oplevelse og Iagttagelse forvexles med Tydning, eller
omvendt. --- I Bogen indeholdes desuden forskellige, især
religionsfilosofiske Afhandlinger, paa hvilke Titlen nogenlunde
passer.

Mine erkendelsesteoretiske Arbejder fra de senere Aar
bestod i speciellere Undersøgelser af forskellige i „Den
menneskelige Tanke“ opstillede Begreber. Først behandlede
jeg Begrebet Totalitet (1917). Dette Begreb spiller en Rolle
ved al Eftertanke. Enten gaar Eftertanke ud fra en given
Totalitet, hvis Elementer den søger at udrede, eller ogsaa
søger den af en given Mangfoldighed at danne en ny Totalitet.
Hvor Totaliteter af sig selv frembyder sig, opløser vi
dem; men naar der er givet et Kaos, stræber vi at forvandle
det til en Totalitet. Mit Motto var et Epigram af Goethe,
i hvilket han kalder Livet „einen wahren Schein“ og „ein
ernstes Spiel“. Totalitet (Atom, Organisme, Personlighed,
Samfund) er et Skin og fører til en blot Tankeleg, naar man
% --- p. 36 ---
\setcounter{footnote}{0}%
\opage{36}i den ser noget én Gang for alle afsluttet, eller naar man
i den ikke ser en Opgave, men en Forklaring. Men der er
Sandhed i Skinnet, og Alvor i Legen, naar vi søger de Love,
efter hvilke mangfoldige Elementer samvirker til et Hele.

Denne Tankegang blev fortsat i en Afhandling om Begrebet
Relation (1921). Det er i de to Afhandlinger de to
første og mest fundamentale Kategorier, Syntese og Relation,
jeg behandler. Baade naar vi analyserer Totaliteter,
og naar vi søger at danne Totalitet af en Mangfoldighed,
gælder det at paavise bestemte Forhold --- indenfor den
givne Totalitet eller mellem de mangfoldige Elementer. En
Totalitet forstaaes ved Forholdet mellem dens Dele, og et
Kaos bliver overvundet ved Paavisning af bestemte Forhold
mellem de tilsyneladende isolerede Elementer. Vi tænker
i Relationer. Vor Tanke er som en Passer, hvis ene Ben
nu staar fast, medens der søges et Sted for det andet Ben,
og vor Erkendelse svarer til Forholdet mellem de to Stillinger.
Kan vi ikke finde nogen Stilling for det andet Ben,
ender vi i et Problem, maaske i et uløseligt Problem. I Afhandlingen
paavises dette nærmere for de forskellige Grupper
af Kategorier og for de forskellige Hovedproblemer. Lejlighedsvis
udtales, at hvis Einstein's Hypotese er rigtig, ligger
der i den et nyt Eksempel paa Relationsbegrebets Betydning,
en Videreførelse af visse ved Kopernikanismen
fremkaldte Tanker. Ganske særlig udvikles tilsidst Relationsbegrebets
Betydning for Etiken\footnote[1]{Senere har jeg (i Videnskabernes Selskabs filosofiske Meddelelser) behandlet endnu et Begreb, der spillede en Rolle i „Den menneskelige Tanke“, men trængte til videre Udførelse, „Begrebet Analogi“. (1922).}.

Ved mine filosofihistoriske Arbejder fra de senere Aar
har jeg kunnet gøre Brug af Optegnelser og Samlinger, som
stammede fra Forberedelserne til mine Laboratorieøvelser.
Jeg havde ofte gennemgaaet Spinoza's Hovedværk med
ældre Studerende. Jeg vilde kunne sige med Lessing, at
naar jeg skulde kalde mig efter Nogen, vilde jeg kalde mig
efter Spinoza. Naturligvis ikke i den Forstand, at jeg skulde
være en Tilhænger af hans System. Men til flere af hans
% --- p. 37 ---
\setcounter{footnote}{0}%
\opage{37}Grundtanker er jeg kommen ad min egen Vej, og det dybe
filosofiske Sindelag, der besjæler hans Tænken, gør ham til
et stort Forbillede for alle Tænkere. Jeg ønskede nu at lette
Forstaaelsen af hans Hovedværk ved at give en sammenhængende
Analyse af dets Indhold og derigennem en Karakteristik
af hans Filosoferen. „Spinoza's Ethica. Analyse
og Karakteristik“ udkom 1918, ligesom de to erkendelsesteoretiske
Afhandlinger i Videnskabernes Selskabs Publikationer.
---

Af andre filosofihistoriske Arbejder fra disse Aar nævner
jeg „Bemærkninger om den platoniske Dialog Parmenides“
(Videnskabernes Selskabs filosofiske Meddelelser.
1920). Mit Resultat af en Undersøgelse af denne mærkelige
Dialog er, at den i første Del givne Kritik af den
platoniske Idelære ikke blot rammer den, men al systematisk
Metafysik, og at den anden Del indeholder saa meget
ufærdigt og uklart, at maaske her et efterladt Udkast er
udgivet af uforstaaende Disciple. Hvad der her interesserede
mig, var den forskellige Stilling, som Begrebet „det
Pludselige“ indtager i Idelæren, hvor Springet fra den ene
Ide til den anden staar som ubegribeligt, og i Læren om
Eros i „Symposion“ (Platon's vigtigste Bidrag til Psykologien).
Logisk betragtet bestaar der hos Platon en skarp
Forskel mellem de forskellige Ideer, men for Platon som
Psykolog stod en lovmæssig bestemt Udviklingsgang fra
sanselig Higen til det højeste Trin af aandeligt Liv som
mulig, og „Pludseligheden“ udtrykker kun Fuldendelsens
Modsætning til Forberedelsen. Platon som Psykolog har
korrigeret ham selv som Erkendelsesteoretiker.\footnote[1]{Senere har jeg sammesteds udgivet en Afhandling om „Platons Bøger om Staten“.}

Et tredie historisk Arbejde fremkaldtes ved Opfordringen
til at skrive en Bog om „Ledende Tanker i det 19.
Aarhundrede“ til den Samling af Værker over det 19. Aarhundredes
Kultur, som Gyldendals Forlag og et svensk Forlag
udgiver paa Dansk og Svensk. Jeg opstillede her den
Hypotese, at det 19. Aarhundrede var Erfaringsbekræftel%
% --- p. 38 ---
\setcounter{footnote}{0}%
\opage{38}sens, Specialiseringens og den tekniske Anvendelses Tid.
Det 17. og det 18. Aarhundredes store teoretiske og praktiske
Ideer underkastedes Erfaringens, Experimentets og
Livets Prøvelse. Det 19. Aarhundredes moralske Karakter
er dets Troskab mod Fortidens store Tanker. Den sande
Troskab bestaar ikke i at sværme i den ideale Verden, der
aabner sig for os, men i med Udholdenhed at omsætte den
til Virkelighed og gennemarbejde den i det Enkelte. Som
Motto kunne man bruge Opfordringen til Gudesønnerne i
„Faust“:

\begin{quote}
Und was in schwankender Erscheinung schwebt,\\
befestiget mit dauernden Gedanken!
\end{quote}

Jeg søgte at vise, at Kant's Filosofi ligger til Grund for
Aarhundredets to Hovedretninger, Romantik og Positivisme,
og at Kritiken af disse to Hovedretninger i Aarhundredets
sidste Del førte til Genoptagelse af Erkendelsesproblemet
som det centrale filosofiske Problem, --- at Naturvidenskabens
største Bedrifter (Paavisningen af Energiens
Bestaaen og af det organiske Livs naturlige Udvikling) var
Gennemførelse af Tanker, der allerede var udtalte i 17. og
18. Aarhundrede, --- og at Aarhundredets største sociale
Fænomener, de store Arbejderorganisationer, var en Gennemførelse
af det 18. Aarhundredes Frihedstrang, da Frihed
konsekvent ogsaa maa være Frihed til Association. Paa
det religiøse Omraade fremtraadte et analogt Træk i Fordringen
om, at religiøse Forestillinger, i hvilken Retning
de saa gaar, skal være Udtryk for personlig Erfaring og
for de enkelte Individers ærlige Tilflugt.

Af en fremragende dansk Naturforsker er det blevet indvendt
mod min Opfattelse, at den undervurderer Naturvidenskabens
Betydning. I den Anledning maa jeg gøre opmærksom
paa, at min Opgave er at tale om de ledende
Tanker, og disse peger afgjort tilbage til de foregaaende
Aarhundreder. De tre sidste Aarhundreder danner et storartet
Hele, hvis Perioder supplerer hverandre. Om Fremtiden
vil blive en Fortsættelse af det aandelige Arbejde un%
% --- p. 39 ---
\setcounter{footnote}{0}%
\opage{39}der Ledelse af de samme Grundtanker, eller om der vil
opdukke helt nye Principer for Teori og Praxis, kan kun
Fremtiden selv lære. Tankearbejdet kan være lige stort
i Gennemførelsens som i Grundlæggelsens Periode.

Det har været mit Livs store Lykke, at jeg har kunnet
hellige mig til filosofiske Studier. Hvad har da Filosofien
været for mig? --- Overfor den Tvivl og den Taage, der
plagede mig i min første Ungdom, har jeg vundet Klarhed
over, i hvilken Retning de store Opgaver for Tænken og
Handlen ligger, og hvorledes der maa arbejdes paa deres
Løsning (hvad enten denne er mulig eller ikke). I selve
Tankearbejdet har jeg fundet den bedste Hjælp mod Mismod
og tungt Sind. Som Spinoza har jeg erfaret, at Vanskeligheder
og Bekymringer forsvandt, naar jeg blot kunde
fordybe mig i Eftertanke (modo possem penitus deliberare!).
Men Filosofien har dog ikke blot været mig et Husmiddel.
Den kan nu engang ikke være et Husmiddel, naar
den blot er et Husmiddel. Jeg har vundet Ærefrygt for
Tankens Verden, hvorsomhelst den aabenbarer sig under
Tidernes Kamp og Uro og peger hen imod de store Opgaver
for vor Stræben. Jeg tror mere paa Idealer end paa
Ideer. Naar vi ikke kommer ud over stadig Stræben, kunde
dette maaske finde sin Forklaring deri, at Tilværelsen selv
ikke er færdig. Om den Del af Verden, der rører sig i Menneskelivet
og i Mennesketanken, gælder det i hvert Tilfælde.

Jeg har den Tro, at Filosofi ogsaa i Fremtiden vil staa
som et nødvendigt Element i aandelig Dannelse. Den betinger
en aandelig Horizont, som ingen speciel Uddannelse,
intet specielt Fag kan give, men til hvilken ethvert Fag kan
give sit Bidrag. Og en Ting gives der, som særlig maa betones,
og som jeg kan bekræfte af lang Livserfaring. Filosofisk
Eftertanke har en befriende Indflydelse overfor Vaner,
Fordomme og Overleveringer. Den skaffer Plads til et personligt
Forhold til Livet, et Førstehaandsforhold. Det er
en Grundtanke i min Religionsfilosofi, at en ægte Livstro
kun voxer frem af egen Erfaring. Men et Tankearbejde
% --- p. 40 ---
\setcounter{footnote}{0}%
\opage{40}kan være den nødvendige Forudsætning for, at vi virkelig
naar til at bygge paa selve Livets Grund. En Livstro forudsætter
ikke, at alle Gaader er løste. Ligesom vort Mod og
vor Villie stedse paany maa opbydes til Kamp mod Taage
og Kaos i os og om os, saaledes vil atter og atter Tankens
Mod være nødvendigt til Arbejde paa Tilværelsens store
Gaader, der stedse fremtræder i nye Former\footnote[1]{I Afhandlingen „Erkendelsesteori og Livsopfattelse“ (Vid. Selskabs filosofiske Meddelelser 1926) har jeg, som i et filosofisk Testamente, udtalt min Opfattelse af Forholdet mellem Forskningen og Livstroen.}.

Hvad der bestaar i Tankens som i Livets Verden, det er
ikke de Former, i hvilke Tanke og Liv fremtræder til given
Tid, men de Kræfter, ved hvilke disse Former er blevne
til, og som for en Tid har lagt sig til Ro i dem. Til disse
Kræfter maa man gaa tilbage, for at befæste sit Haab til
Fremtiden. Men selv om man har den Tro, at der stedse
paany vil opstaa aandelige Værdier, har man dog ikke Ret
til at tro, at vi forud skulde kunne vide, hvilke disse Værdier
vil være.

\piecerule

% --- p. 41 ---
\setcounter{footnote}{0}%

\opage{41}\essayhead[Platons Opfattelse af Kunst og Religion]{PLATONS OPFATTELSE AF KUNST \\ OG RELIGION, SAALEDES SOM DEN FREM- \\ TRÆDER I BØGERNE OM STATEN\footnote[1]{Et Foredrag, der er bygget paa min i Videnskabernes Selskabs filosofiske Meddelelser (1924) fremkomne Afhandling: \textit{Platon's Bøger om Staten.} Analyse og Karakteristik.}}

Platon fremtræder i sine Skrifter som Tænker, som Lovgiver
og som Profet. De forskellige Interesser, der saaledes
rører sig hos ham, kommer ikke sjelden i Strid med hverandre.
Ikke blot staar Bøgerne om Staten af denne Grund
i en vis Modsætning til andre af hans Værker, men indenfor
selve Bøgerne om Staten anlægges der forskellige Synspunkter,
særlig hvad Opfattelsen af Kunst og Religion angaar.
Som Lovgiver taler han især i tredie og fjerde Bog,
og som Digter fremtræder han i den kunstneriske Udformning
af Værkets enkelte Dele. Det filosofiske Grundlag,
han mente at have fundet, har gjort det muligt for ham at
sammenarbejde Udkast, der aabenbart er forfattede til forskellige
Tider og paa forskellige Stadier af hans Udvikling,
til en Helhed og derved skabe det i Aandshistorien saa betydningsfulde
Værk, der kaldes Platon's „Stat“. --- Vi skal,
som sagt, her holde os til hans Opfattelse af Kunst og Religion.

1. Bøgerne om Staten handler om Begrebet Retfærdighed,
og Platon mener, at det er lettest at belyse dette Begreb
ved at undersøge, hvorledes Retfærdigheden maa
fremtræde i en fuldkommen Stat. I den Beskrivelse, han
forsøger at give af en saadan Stat, bliver der et Spørgsmaal,
der fanger hele hans Interesse, hvorledes Herskerne i en
saadan Stat bør opdrages. Foreløbig staar han her paa det
% --- p. 42 ---
\setcounter{footnote}{0}%
\opage{42}gammelgræske Stade, hvor der --- foruden gymnastisk Uddannelse
--- kræves en aandelig Udvikling gennem Indlevelse
i den nationale Digtning, i hvilken Folkets Gudetro
har fundet et fyldigt Udtryk. En Vurdering af Digtningen
bliver nødvendigvis en Vurdering af de religiøse Forestillinger,
der fremtræder i den, og Platon opstiller da, under
skarp Kritik af Homer og Hesiodos, visse teologiske Normer
eller Forbilleder. Man kan, mener han, ikke lade Digterne
have frit Spil, fordi de Unge saa let vilde faa urigtige
Forestillinger om Guderne, hvad Lovgiveren ikke kan tilstede.
Platon siger ikke ligefrem, hvorfra han véd, hvilke
de rigtige Forestillinger om Guderne er, men det fremgaar
tydelig nok af hans Drøftelse, at de teologiske Normer eller
Forbilleder bestemmes ved hans egen Filosofi. Der er tre
Hovednormer.

a. Guddommen er god og kan ikke være Aarsag til noget
Ondt. Digterne maa derfor ikke tillægge Guderne onde
Handlinger. Guderne maa ikke fremstilles som Aarsager til
Ulykker og Lidelser, uden forsaavidt saadanne kan blive
Midler til Rensning og Forbedring. Men nu gives der ---
lærer Platon baade her og i andre Skrifter --- mere Ondt
end Godt i Verden, og Guddommen kan derfor kun være
Aarsag til meget Lidt. Det virkelige Onde i Verden maa
have en anden Aarsag end Guddommen. --- Ved at forfølge
denne Tankegang kom Platon i sit Alderdomsværk „Lovene“
til Antagelse af en ond Verdenssjæl, en Antagelse,
som man har sat i Forbindelse med en vis Indflydelse af
Zarathustralæren, der skal kunne spores hos Platon i hans
sidste Aar. --- En særlig Anvendelse af denne første Norm
er Forbudet mod at lære, at Guderne ligger i Strid med
hverandre. Det vilde skade den Endrægtighed, der bør herske
mellem Statens Herskere! Ikke engang allegorisk maa
saadanne mytologiske Forestillinger benyttes, da unge Mennesker
ikke kan skelne mellem Allegori og bogstavelig
Sandhed.

b. Guddommen bør fremstilles som uforanderlig. Kunde
en Gud forandre sig, vilde Forandring være Et med For%
% --- p. 43 ---
\setcounter{footnote}{0}%
\opage{43}ringelse, Nedværdigelse. En Gud maa derfor ikke fremstilles
i Menneskeskikkelse. Og en Gud vil hade Løgn, aldrig
dølge Sandheden; Sanddruhed er en Art af Uforanderlighed.
Og ved Løgn maa da ikke blot forstaas den ydre
Løgn, den, der ytrer sig i Ord, men den indre Løgn, som
skyldes Uvidenhed, og som kun Selvbedrag hindrer os i at
opdage. Det er kun en vanfør Sjæl, der skammer sig mere
ved frivillig end ved ufrivillig Usandhed.

c. Den tredie Norm er, at Guderne ikke maa fremstilles
som underlagte stærke Sindsbevægelser, som Dødsfrygt, Latter,
Sorg, Elskovstrang. Der skal skildres som det Højeste en
% PRINTED AS IS (read at the image, odd spots o1)
Sindsligevægt, en Selvbeherskelse, der ikke gør afhængig af
ydre Kaar, men er den rette Tapperhed overfor, hvad Livet
kan bringe. --- Og Statens Herskere skal efterligne, hvad
der saaledes hævdes som den rette Opfattelse af Gudernes
Væsen. Herskerne har den ene, store Opgave at sørge for
Statens Frihed og Sikkerhed. Hverken i deres Fantasi eller
i deres personlige Færd maa de efterligne tragiske eller
komiske Digtere, overhovedet optræde i anden Karakter
end deres egen. Platon nægter her, hvad han i „Symposion“
havde udtalt, at det er samme Mands Sag at digte Tragedie
og Komedie. Det er Arbejdsdelingens Nødvendighed, som
har faaet Magt over ham og nu faar en etisk Anvendelse.
Af Digterne foretrækkes Epikerne for Dramatikerne, fordi
hine ikke som disse efterhaanden gør sig til Et med forskellige
Karakterer og derved spalter sig. --- Det vil ikke
være muligt at bringe Overensstemmelse mellem Platons
Udtalelser i det her omtalte Punkt. Han havde selv, og det
ikke blot som Digter, Blik for Livets sammensatte Karakter.
I en senere Dialog (Filebos) taler han om Livets Tragedie
og Komedie. Men han fik atter og atter Betænkelighed
ved, at en Enkelt i sin Kunst, som Digter eller Skuespiller,
udtrykte denne Livets sammensatte Natur i sin Fantasi
eller i sin Optræden.

2. Allerede af denne etisk-pædagogiske Omtale af Kunsten
i tredie og fjerde Bog ser man, hvilke strenge Fordringer
Platon stiller til den herskende Klasse i Samfundet, skønt
% --- p. 44 ---
\setcounter{footnote}{0}%
\opage{44}han foreløbig blot vil reformere den gammelgræske Opdragelse.
Men i sin nærmere Udvikling af Herskernes Opgave
kommer han til at skærpe Fordringerne, saaledes at
han ikke blot gaar ud over den traditionelle Pædagogik,
men kommer til at stille helt nye Krav, idet han kræver en
streng videnskabelig Disciplin for dem, der skal raade i
Staten. Herskerne kan, mener han, kun sørge for Statens
Frihed og Sikkerhed, naar de under alle Omskiftelser og
Fristelser bevarer Tanken om Statens egentlige Væsen i
Modsætning til Krav og Opgaver, som kunde føre til at forfalske
den rette Forfatning. Vogterne maa derfor have Evne
til at forsvare og fastholde Tanken om noget, der staar
fast og bevarer sin Enhed med sig selv under alle Forandringer.
Men denne Evne uddannes kun (hævder Platon i
Slutningen af femte Bog) gennem videnskabeligt, særlig
matematisk og filosofisk, Studium, og i sjette og syvende Bog
gaar han saa nærmere ind paa Ordningen af dette Studium.
Naar man ofte har fundet en stor Modsætning mellem
tredie, fjerde og femte Bog paa den ene Side, sjette og syvende
Bog paa den anden Side, er det, fordi man har overset,
at der er et Slægtskab mellem det praktiske Skøn (f. Ex.
om en Forfatnings Bestaaen i Enhed med sig selv) og den
strenge Gennemførelse af Identitetsbegrebet, som er det Karakteristiske
for al Videnskab. Naar Platon skærper Fordringerne
til Herskerne ved at kræve videnskabelig Indsigt
af dem, er det den naturlige Overgang fra blot Lighed til
fuldstændig Identitet, han gør. Her ligger det Enhedsbaand,
der knytter det vidunderlige Værks forskellige Dele sammen,
selv om de er forfattede til forskellige Tider.

Naar Platon saa, efterat han saaledes har skærpet Kravene
til den højeste Dannelse indenfor Samfundet, igen (i
tiende Bog) vender tilbage til at drøfte Kunstens Betydning,
betoner han stærkere end før to Synspunkter, der allerede
i hans første Omtale af Kunsten gjorde sig gældende,
nemlig Kunstens Forhold til Sandhed og Virkelighed og
Kunstens Væsen som Efterligning. Og han tager ikke blot
Hensyn til Poesi, men ogsaa til Billedkunst. 

% --- p. 45 ---
\setcounter{footnote}{0}%
\opage{45}Medens han i sin første Omtale af Kunsten stillede den
i Forhold til Teologi og Moral, er det nu dens Forhold til
Videnskaben, der staar i første Linie. Og han stiller da
Kunsten langt under Videnskaben, fordi den vel giver et
Billede af noget Virkeligt, men et ufuldkomment Billede,
der staar lige saa meget tilbage for den virkelige Ting eller
Handling, som disse hver for sig staar under det Almenbegreb
eller den almene Lov, under hvilken de hører. Et
Digt om en Heltegerning staar langt under selve Heltegerningen.
Bedrifter skal man ikke digte om. I denne Sammenhæng
kommer Platon igen tilbage til Faren ved Efterligning
af Sindsbevægelser og Lidenskaber. Han mener, at
naar Kunstneren ynder den Art Emner, er det, fordi det er
lettere at fremstille meget lidenskabelige Naturer end energiske
og harmoniske Naturer. Den, hvis Tanke kan bevæge
sig i Ideernes evige Verden, vil i hvert Tilfælde ikke have
Tilbøjelighed til at dvæle ved de Følelser, menneskelige
Forhold kan frembringe, da de taber al Betydning overfor
den rene Sandhed.

I sin første Omtale af Kunsten (i tredie og fjerde Bog)
taler Platon som Lovgiver; i den nye Omtale, som tiende
Bog bringer, taler han væsentlig som Filosof, og Overgangen
er som sagt sket gennem Omtalen af Videnskaben i
sjette og syvende Bog. I Modsætning til den Sandhed, videnskabelig
Forsken kan bringe, synker Kunsten ned til
at være Efterligning af sanselige Ting. Det er Platons Idealisme,
der gør ham til Realist eller Impressionist i hans
kunstneriske Opfattelse. Han kommer derved til at miskende
Kunstens Evne til at forme individuelle Skikkelser,
danne Helhed, hvor Virkeligheden kun giver Fragmenter
eller Tilløb. Ved denne Evne til at forme givne Elementer
til et Hele med individuel Karakter giver Kunsten mer,
end Videnskaben, der maa holde sig til almene Begreber,
kan give, og forsaavidt maa Kunsten siges at staa over Videnskaben.
Kunsten overvinder Kaos bedre, end Videnskaben
formaar at gøre det. --- Ingen har bedre forstaaet dette
end Platon selv, der i sine Dialoger har givet os saa mange
% --- p. 46 ---
\setcounter{footnote}{0}%
\opage{46}prægtige Skildringer af Personligheder, af Scener og af
Samtaler, Skildringer, der fremtræder med plastiske Helheders
Præg. Ved at tale mod Kunsten taler Platon da
egentlig imod sig selv. Han erklærer da ogsaa tilsidst, at
han vilde blive glad, hvis nogen kunde gendrive ham, da
han af egen Erfaring kender Kunstens tryllende Magt.

Naar Aristoteles fandt Tragediens Betydning deri, at den
vækker Følelser som Frygt og Medlidenhed paa en saadan
Maade, at vi befries for den Pine og Bundethed, som disse
Følelser medfører, naar det er Livets virkelige Forhold,
der fremkalder dem, saa har man (Finsler: Platon und die
aristotelische Poetik) deri maaske med Rette set en direkte
Indsigelse mod Platon's Opfattelse af Tragedien. Aristoteles
har teoretisk mere Sans for det Individuelle end Platon og
har derfor kunnet hævde Betydningen af at leve sig ind i
digteriske Skikkelser. En saadan Indleven havde for ham
en „rensende“ Virkning paa Sindet.

3. Første Gang, Platon i „Staten“ taler om Religionen,
er det, som vi har set, i Anledning af Digtekunsten, der i
Athen spillede saa stor en Rolle ved Opdragelsen, da de Unge
gennem den fik deres Forestillinger om Guderne bestemte.
Kritiken af Digtekunsten blev en Kritik af Teologien. Den
Vægt, han senere lægger paa videnskabelig Erkendelse, fører
til fra en ny Side at bedømme de overleverede religiøse
Forestillinger. Det er karakteristisk, at Spørgsmaalet om
Kultus, der religionshistorisk viser sig at være saa afgørende
et Element i al folkelig Religion, ikke nærmere drøftes. Han
overlader det til „den Fortolker, som vi har faaet fra vore
Fædre, og som tronende paa Jordens Midtpunkt giver Svar
til alle Mennesker“. Det er altsaa Apollon, der skal raade
for Kultus, og hvis Lovgiverne er kloge, indser de, at de
ikke forstaar sig paa at ordne Kultus, men gør bedst i at
overlade alt Herhenhørende til den fædrene Gud. Hovedspørgsmaalet
bliver, hvorledes han paa Grundlag af sin
Filosofi maa stille sig til de religiøse Forestillinger. Selv
benytter han Forestillingerne om Guder som personlige Væsener,
men der kan neppe være Spørgsmaal om, at hans
% --- p. 47 ---
\setcounter{footnote}{0}%
\opage{47}Opfattelse af Guderne er betinget ved hans Lære om „Ideerne“,
d.v.s. om de højeste Tanker, der gør os Tingene i
Verden forstaaelige. Det afgørende Sted findes dog ikke i
„Staten“, men i den noget senere Dialog „Fajdros“, hvor
det siges, at medens vi Mennesker kun i enkelte Øjeblikke
kan hæve os op til Ideernes Verden, saa lever Gud stadig
i denne Verden --- og det er herpaa, hans Guddommelighed
beror. Begrebet „guddommelig“ betinger altsaa for Platon
Begrebet Gud. Og han ved, hvad det Guddommelige er ---
det er de Tanker, der tilsidst gør Verden forstaaelig, og
som derved tillige hæver over alt Sanseligt og Forgængeligt.
Nyplatonikeren Plotinos hentyder til denne Udtalelse af
Platon og føjer til, at det, som gør Gud guddommelig, maa
være den højeste Gud. I nyere Tid vilde man udtrykke dette
saaledes, at Begrebet „guddommelig“ er et Værdibegreb,
og at Menneskene naturlig overfører paa deres Guder de
værdifuldeste Egenskaber, de kender. Vi kommer da til at
sande Goethe's Ord:

\begin{quote}
Daher der Völker löblicher Gebrauch,\\
dass jeglicher das Beste was er kennt,\\
er Gott, ja seinen Gott benennt.
\end{quote}

Hvis Platon har en Teologi, saa er den i hvert Tilfælde
afhængig af hans Filosofi.

Men i en vis Uafhængighed af hans Forestillinger om
Guderne staar hans Udødelighedstro. Paa den store Rejse
til Nordafrika og Syditalien, han foretog efter Sokrates'
Død, og som fik saa stor Indflydelse paa hans Aandsliv,
lærte han ikke blot Videnskabens (særlig Matematikens)
Betydning at kende, men han har ogsaa levet sig ind i den
pytagoreiske Lære om Sjælevandringen. Ikke blot i „Staten“,
men baade i tidligere Dialoger („Gorgias“, „Menon“,
„Fajdon“) og i senere („Fajdros“) kommer denne Lære
frem. Den har dannet den religiøse Baggrund for hans
Tankeliv. Man har endog ved Sammenstilling af, hvad han
i forskellige Dialoger udtaler om dette Spørgsmaal, fundet,
at der her for Platon's Aand har staaet et stort Fantasibillede,
en antik Comedia Divina. Platon optræder her af%
% --- p. 48 ---
\setcounter{footnote}{0}%
% CHECK p. 48: notes paired with the marks in order (the layer misread a number)
\opage{48}gjort som Profet. I Slutningen af „Staten“ fortæller han
om en Mand, der faldt i et Slag og kom til Sjælenes Rige,
hvor han blev Vidne til, hvorledes Sjælene efter deres Domfældelse
og den derefter følgende Løn eller Straf igen sendtes
ned til Jorden, idet det saa beroede paa deres eget Valg,
hvilke Livskaar de stedtes i. Manden selv kom tilbage til
Jorden som „Forkynder af det Hinsidige“ --- ligesom Orfeus
i et gammelt Digt kom til Underverdenen, ikke for,
som i den senere Myte, at hente Eurydike, men for at forkynde
Menneskene deres hinsidige Skæbne.

Forestillingen om Sjælevandringen og om Sjælenes Genkomst
forekommer hos nogle af de mest primitive Folk.
Platon kunde forsaavidt sige ligesom Myers, den moderne
Spiritismes betydeligste Repræsentant, at han vedkender
sig Slægtskabet med „Stenalderens arme Tænkere“.\footnote[1]{Smlgn. min Karakteristik af Myers i \textit{Oplevelser og Tydning.} p. 25---33.} Men
i de mest primitive Former, i hvilke den optræder, har
Forestillingen om Sjælevandringen egentlig ingen etisk Betydning.
En saadan Betydning fik den derimod i Indien og
i Grækenland, da haarde Skæbner havde knækket det oprindelige
Livsmod.\footnote[2]{Smlgn. min Religionsfilosofi. \textsuperscript{3} p. 97---100.} Hos Platon tjener Sjælevandringslæren
dels til at begrunde Antagelsen af, at Sjælenes Antal er
begrænset (saa at de samme Sjæle maa komme igen, hvis
Menneskelivet ikke skal høre op), dels til at begrunde Vigtigheden
af at føre Livet her i Verden paa rette Maade, saa
at man kan være forberedt til at træffe et godt Valg, naar
man skal begynde Livet igen. Dette Valg bestemmes nemlig
ved Erfaringen fra det tidligere Liv, --- og det viste sig,
at de Sjæle, der havde friske Erindringer fra deres Jordeliv,
traf det bedste Valg, hvorimod de, der var belønnede
med et Ophold i Himmelen, ofte tog fejl! Himmelen er altsaa
ikke nogen god Skole. --- Platons Mytologi har altsaa
en etisk Kærne. Det kommer an paa at føre Livet i Sandhed
og Harmoni og med et mandigt Sind. Platon kommer
egentlig ikke ud over Sokrates' frejdige Tro paa, at et godt
% --- p. 49 ---
\setcounter{footnote}{0}%
\opage{49}Menneske kan intet Ondt vederfares. Ligesaalidt som Platon's Teologi
bringer noget Nyt i Forhold til hans Filosofi,
idet Guder saavel som Mennesker maa underordne sig under
Ideernes Højhed, ligesaalidt giver Sjælevandringslæren
ham helt nye etiske Motiver. Han staar her omtrent som
senere Spinoza og Kant, hvis Etik var uafhængig af deres
religiøse Forestillinger. Platon's Stilling var den, at Retfærdigheden
kan fastholdes under ydre Modstand og uden
Støtte fra Guder og Mennesker, men at der ingen Grund er
til at tro, at den under sædvanlige Forhold ikke skulde
finde Støtte baade fra Guder og Mennesker. Platon nærmer
sig til, hvad Stoikerne siden tænkte, at den etiske Karakter
i sit Væsen er uafhængig af den ydre Hjælp og af
ydre Modstand.

\piecerule

% --- p. 50 ---
\setcounter{footnote}{0}%

\opage{50}\essayhead{GIORDANO BRUNO}

Den Tid i Europas Kulturhistorie, man plejer at kalde
Renaissancens, Genfødelsens Tid, omfatter det fjortende,
femtende og sextende Aarhundrede. Den aabnes af de tre
store Digtere, Dante, Boccacio og Petrarca; efter dem kommer
de lærde Humanister og derefter igen de banebrydende
Naturforskere og Filosofer: Pomponazzo og Machiavelli,
Leonardo og Telesio, Bruno, og senere Galilei og Campanella.
Her forberedtes den Verdens- og Livsanskuelse, som i
sine store Træk kan siges at være ejendommelig for den
nyere Tid. To nordiske Astronomer, Kopernikus (om hvis
Herkomst Tyskere og Polakker strides) og Danskeren
Tycho Brahe, gav Grundlaget, og begge hyldedes da ogsaa
med Begejstring af den Mand, vi her skal beskæftige os
med, Bruno, Renaissancens typiske Repræsentant i Filosofien,
hos hvem alle den mærkelige Tids Traade væves sammen.
Bruno saa i den nye Astronomi Mulighed for en ny
Verdensanskuelse, som han fremstillede med Begejstring,
og hvis vigtigste Konsekvenser han saa. Støttet til et af
Videnskabens store Vendepunkter sprængte hans Tanke
Fortidsoverleveringens snevre Grænser og paaviste nye Problemer,
hvis Løsning skulde blive Fremtidens Sag. Det hører
til det Store hos Bruno, at han er klar over, at Arbejdet
for at vinde Sandhed aldrig kan afsluttes.

Bruno er født 1548 i Nola i Syditalien. Hans Døbenavn
var Filippo, men da han 15 Aar gammel gik i et Dominikanerkloster
i Neapel, fik han Navnet Giordano (efter Dominicus'
nærmeste Efterfølger). Det var et skæbnesvangert
Skridt, han her gjorde, bestemmende for hans hele Livs%
% --- p. 51 ---
\setcounter{footnote}{0}%
\opage{51}gang. Sandsynligvis har det overvejende været Lyst til at
studere i Ro, der har ført ham til Klosteret, men det var
det foreløbige Resultat af disse Studier, der, da han var
28 Aar, fik ham til at flygte fra Klosteret. Der svævede en
Anklage for ikke mindre end 130 Kætterier over hans Hoved.
Tvivl om Treenighedsdogmet har været det mest graverende;
men det bebrejdes ham ogsaa, at han ingen Helgenbilleder
vilde have i sin Celle, kun et Krucifix. Han
flygtede op i Norditalien, og her har han levet i nogle Aar
som Pogelærer og som Lærer i Astronomi for unge Adelsmænd.
Det er sandsynligvis i disse Aar, at hans Verdensanskuelse
er bleven grundlagt gennem et gentaget Studium
af Kopernikus' berømte Værk (der var udkommet 1543),
hvis Idéer han ved første Læsning ikke havde kunnet gaa
ind paa. I sine senere Skrifter priser han ham højt: „Du
højsindede Kopernikus, hvis Værk allerede gjorde Indtryk
paa mig i mine ganske unge Aar, skønt jeg mente, det stred
baade mod Sansning og Fornuft!“

Lad os stanse lidt her. Hvorfor stred Kopernikus' Lære
om Jordens Bevægelse „mod Sansning og Fornuft“?

Ja, Sanserne viser os jo, at det er Solen og Stjernerne,
der bevæger sig, medens Jorden, vi lever paa, staar stille!
--- Kopernikus' Svar herpaa var, at vi ikke har Ret til at
anse vort Standpunkt og Synspunkt, det jordiske, for det
eneste mulige, og han viser ad matematisk Vej, at for en
Beskuer paa Solen vilde Jorden bevæge sig, ligesom alle
Planeterne, medens Solen vilde staa stille. Og tillige viser
han, at dette Synspunkt fører til den simpleste Opfattelse,
fri for alle de Hjælpehypoteser, man i Tidens Løb havde
maattet udforme for at gennemføre det jordiske Standpunkt
som det eneste rigtige. Vi forholder os da --- som
allerede før Kopernikus den tyske Filosof Cusanus havde
sagt --- til Stjernerne, som de Søfarende, der sejler forbi
Kysterne. Kopernikus bygger allerede paa den store Relativitetslov,
der lige op i den allernyeste Tid spiller saa stor en
Rolle i Videnskaben. --- Denne kopernikanske Opfattelse
har Bruno ikke blot gjort til sin, men han er med Dristig%
% --- p. 52 ---
\setcounter{footnote}{0}%
\opage{52}hed gaaet store Skridt videre, især paa to Punkter. Kopernikus
ansaa Fixstjernehimlen (den „ottende Sfære“) for
ubevægelig og mente, at i den sad alle Fixstjerner fast, alle
altsaa i lige Afstand fra Jorden. Og han antog endnu den
gamle Lære, at Planeterne sad hver i sin roterende Krystalsfære,
i forskellig Afstand fra Jorden. Bruno viser, at vi ikke
har Ret til at sætte Grænser for Verdensrummet og at antage,
at alle Fixstjerner er lige langt fra os. Fjerne Skibe ude
paa Havet synes jo at være i samme Afstand fra os! Det er
et Punkt, der trænger til nærmere Undersøgelse; men naar
man blindt tror paa den gamle Opfattelse, kommer man
aldrig til at undersøge det! --- Man ser, at Brunos dristige
Tænkning udtrykkelig kræver, at der søges Bekræftelse
ved Erfaring. --- Hvad Planetsfærerne angaar, fandt hans
Tvivl Bekræftelse ved Tycho Brahe's Undersøgelser om
Kometernes Bevægelser. Ifølge disse Undersøgelser maatte
Kometerne gaa tværs igennem Planetsfærerne og for længe
siden have slaaet disse Krystalhimle istykker. Bruno lovpriser
derfor i sine senere Skrifter Tycho Brahe som den
ypperste af Tidens Astronomer (Ticho Danus, nobilissimus
atqve princeps astronomorum nostri temporis).

Dette, hvad Sansningen angaar. Naar „Fornuften“ havde
gjort Vanskeligheder, var det, fordi Afstandene indenfor
Verdenssystemet nu blev saa uhyre store, at Tanken ikke
syntes at kunne faa et samlet Hele ud deraf. Jorddiameterens
Længde skulde jo være forsvindende i Forhold til Jordens
Afstand fra Solen! Bruno viser, at Fornuften i Virkeligheden
ikke har nogen Trang til et afsluttet Verdensbillede.
Tværtimod fører, naar vor Fornuft er vaagen, den ene
Tanke over i den anden, ligesom vi i Talrækken stedse kan
føje Tal til Tal uden Grænse. Bevisbyrden paahviler den,
der antager, at Verden er begrænset. Skulde Naturen være
fattigere end vor Tanke? Og hvorfor skulde der være absolut
Hvile noget Sted? Hvorfor skulde Jorden og Fixstjernerne
være Undtagelser fra den stadige Bevægelse, Erfaringen
overalt viser os, naar vi ser nøje til? --- Allerede som Barn
havde Bruno lært at se, at Opfattelsen beror paa Stand%
% --- p. 53 ---
\setcounter{footnote}{0}%
\opage{53}punktet. Fra det vin- og skovbevoksede Bjerg Cikada ved
Nola tog Vesuv sig ud som en nøgen Masse, der laa i Taage;
men da han engang kom til Vesuv, var det omvendt Cicada,
der laa nøgen og taaget. Ud fra denne Erfaring forstod han
Kopernikus, ja forstod tillige, at man ikke kunde stanse,
hvor han stansede. Naar Jorden er en Planet, saa er vi jo
selv i Himmelen. Og naar man er i Forlegenhed med, hvor
Guds Bolig skal være, da den faste Himmel, der hidtil stod
som hans Bolig, ikke længere er der, saa lærer Bruno, at
Guddommen er den uendelige Kraft eller Sjæl, der rører
sig i Alt, ogsaa i Menneskets Sjæl, ja endog længere inde i
vor Sjæls Inderste, end vi selv kan trænge! Ingen Tanke
kan omfatte Guddommens Væsen, da vi ved nærmere Eftertanke
over det føres ud i det uendelige. Religionens Betydning
er praktisk, moralsk. Bibelen vil ikke lære os Astronomi
og Filosofi, men virke paa Menneskers Sind og Villie,
og de religiøse Forestillinger har kun symbolsk Betydning.
Bruno lærer ikke, at hvad der er sandt i Filosofien, kan
være falsk i Teologien; men han lærer, at en og samme
Sandhed er udtrykt paa to Maader, i Filosofiens Tanker
og i Religionens Billeder. ---

Vi forlod Bruno i Norditalien. Af en eller anden Grund
drog han videre mod Nord og kom til Genève. Her aflagde
han Munkekutten og optraadte nu i Vams og med Kaarde.
Han synes at være traadt over til Protestantismen, skønt
han nægtede det senere; idetmindste er han bleven exkommuniceret
af det protestantiske Præsteskab og har gjort
Afbigt. Han har fra denne Tid af været ugunstig stemt overfor
Protestantismen. Han fandt dér mange Paver istedetfor
én, og han savnede den store Omsorg for Syge og Fattige,
som han stadig beundrede hos Katolicismen. Dette gjorde
ham hjemløs. Han flakker nu om i Frankrig, England og
Tyskland, uden at vi altid kender Aarsagerne til, at han
skifter Sted. Aarsagen har ofte været Krige og politiske
Uroligheder, men ligesaa ofte hans eget urolige Sind, hans
fordringsfulde Optræden og hans Tankegangs Fjernhed fra
Tidens videnskabelige og kirkelige Forestillinger.

% --- p. 54 ---
\setcounter{footnote}{0}%
\opage{54}Han læste som Professor i Toulouse og senere i Paris. En
stor Rolle ved hans Undervisning spillede en Teori om Hukommelseskunst,
man dengang nærede store Forventninger
til, og som han har brugt som psykologisk Indledning til
den egentlige Filosofi. Den franske Konge Henrik III lod
ham kalde til sig og synes at have ment, at han sad inde
med en hemmelig Visdom, der kunde løse alle Gaader. Baade
fra Toulouse og Paris har Bruno indledt Forhandlinger
om at komme tilbage til Italien i Forsoning med Kirken.
Han har ment, at Kirken kunde lade sig nøje med hans
symbolske Tydning af dens Dogmer. Men der blev svaret
ham, at han strax skulde vende tilbage til sit Kloster; saa
kunde man tale om det Andet.

Opholdet i London (fra 1583 til 1585) var Brunos lykkeligste
Tid. Han boede i den franske Gesandts Hus og færdedes
i aandfulde og fornemme Kredse. Sine mest dybsindige
og glimrende Værker har han forfattet i disse Aar.
Medens han før og siden skrev paa Latin, er de sex Bøger,
han forfattede i London, skrevet paa Italiensk. Han forsvarede
den kopernikanske Teori og udviklede de filosofiske
Idéer, den store Udvidelse af Verdensbilledet fremkaldte
hos ham. I en snevrere Kreds fandt han Interesse for sine
Idéer; ogsaa Dronning Elisabeth satte Pris paa Samtaler
med ham, og han priser hende i høje Toner som Amfitrite,
Havdronningen, eller som Diana, den kongelige Jomfru.
Men hos de Lærde og i større Kredse fandt han ingen Forstaaelse,
og han giver drastiske Skildringer af den lærde
Dumhed og Vigtighed, han mødte paa sin Vej. England var
ham idethele for barbarisk i Sæder og Levevis. Allerede i
Frankrig følte han sig i et halvbarbarisk Land, men i England
fik han et stærkere Indtryk i den Retning.

Da den franske Gesandt blev kaldt tilbage, fulgte Bruno
med ham. I Paris holdt han en stor Disputats om sine Teorier.
En enkelt Sætning blandt dem, han opstillede, har Interesse
for, hvad der for Øjeblikket interesserer Naturforskere
og Filosofer. Han gør opmærksom paa, at naar Tiden
maales efter Himmellegemernes Bevægelser,  og naar disse
% --- p. 55 ---
\setcounter{footnote}{0}%
\opage{55}tager sig forskellig ud, baade hvad Hurtighed og Retning
angaar, fra hver enkelt Klode, saa faar hver Stjerne sin
egen Tid, forskellig fra de andre Stjerners.

Efter i flere Aar at have flakket om fra det ene tyske
Universitet til det andet slog han sig en Tid ned i Frankfurt,
hvor han forfattede en hel Række Skrifter paa Latin,
i hvilke han gav en mere systematisk Fremstilling af sin
Lære end i de halvt digteriske italienske Skrifter. Midt i
denne Virksomhed bryder han pludselig op og drager til
Italien (1591). Man synes i Frankfurt at have været betænkelig
ved at have denne Kætter i sin Midte, og tillige havde
en ung Venetianer, Mocenigo, der interesserede sig for hans
Skrifter, indbudt ham til at bo i hans Hus. Det har været
Længselen efter Fædrelandet, Lede ved Barbariet og
Uimodtageligheden i de nordiske Lande for de nye Tanker,
der har ført ham tilbage. Og han havde ikke opgivet Haabet
om Forsoning med den gamle Kirke, uden at han behøvede
at opgive sin Overbevisning. Det skulde blive en
stor Skuffelse. Reaktionen i Italien var i fuld Gang, og Inkvisitionen
var aarvaagen.

Mocenigo blev skuffet over sin Lærer og mistænkte ham
for at holde sine vigtigste Tanker skjult for ham. Da han
blev klar over, at hans Gæst var en arg Kætter, meldte han
ham efter sin Skriftefaders Raad for Inkvisitionen, og
Bruno blev fængslet af den (1592). I et Aar sad han fængslet
i Venedig. I Forhørene optraadte han først med stor
Frejdighed, fortalte sit Levned og udviklede sine filosofiske
Idéer. Han mente, at han kun indirekte havde angrebet
Kirkelæren, skønt han indrømmede, at han ikke havde
overholdt Skikkene, og han udviklede, at han i Kirkelæren
saa et symbolsk Udtryk for, hvad han selv ad anden Vej
var kommen til. Men man lod sig ikke nøje hermed, og
hans Haab om Forlig med Kirken bristede. Da faldt han
sammen og erklærede sig villig til fuldstændig Tilbagekaldelse.
Hvad Udfaldet af Processen i Venedig vilde være
blevet, kan ikke vides. Men nu greb den romerske Inkvisition
ind, og for at gøre Paven til Behag udleverede Venedigs
% --- p. 56 ---
\setcounter{footnote}{0}%
\opage{56}Regering ham til Rom (1593). En hel ny Proces begyndte.
Man gennemstuderede hans Skrifter. Det synes, som man
har haabet at slaa et stort Slag ved at faa denne store Kætter
til udtrykkelig at fornægte sine videnskabelige Lærdomme,
ikke blot sine teologiske Kætterier. Blandt hans
Dommere var den ivrige Kætterforfølger Kardinal Sanseverino,
der havde fejret Bartolumæusnatten som en Glædesdag,
og Bellarmin, der senere ogsaa var med til at domfælde Galilei,
og som maaske netop ved Bruno's Proces fik Øjet op for
de farlige Konsekvenser af Kopernikus' hidtil upaaagtede
Lære. Hvad der er sket i de syv Aar, i hvilke Bruno sad
fængslet i Rom, véd man ikke; Akterne siges at være forsvundne
fra Vatikanets Arkiv. Kun de allersidste Forhandlinger
kender man, da en tysk Lærd skrev dem af, medens
Pius IX i 1849 en Tid var fordreven fra Rom. Da vi igen
hører om Bruno, staar han fast, nægter at kalde nogetsomhelst
tilbage. Den 9. Februar 1600 blev han saa degraderet,
exkommuniceret, og overgiven til den verdslige Øvrighed
(med den sædvanlige hykleriske Anmodning om at undgaa
Blodsudgydelse). Ved denne Lejlighed har Bruno sagt: „I,
som fordømmer mig, er mere bange end jeg“. Den 19.
Februar blev han brændt, og hans Aske kastet i Tiberen.

En tysk Lærd, der var gaaet over til Katolicismen, var
tilstede ved Henrettelsen og har beskrevet den i et Brev.
Han slutter sin Beskrivelse med de Ord: „Nu kan han gaa
hen til en af de andre Verdener, han troede paa, og fortælle,
hvorledes man i Rom behandler Folk som ham!“ ---
Det synes at have været Muligheden af andre beboede Verdener
end vor Jord, der især er falden i Øjnene, fordi Gud
kun kunde blive Menneske én Gang. I Middelalderen stod
af samme Grund Troen paa Antipoder som et Kætteri.

Efter at have fulgt Bruno paa hans Livsgang og lært
nogle af hans vigtigste Tanker at kende, vil vi nu forsøge
en Karakteristik af hans Personlighed.

Ifølge Forhørene var han en lille slank Mand med spin%
% --- p. 57 ---
\setcounter{footnote}{0}%
\opage{57}kelt Legeme, blegt Ansigt, eftertænksomt og livligt, men
dog sørgmodigt Udtryk. --- Af hans Skrifter, med deres direkte
og indirekte Selvkarakteristik, faar vi Indtrykket af
en Karakter fuld af Modsætninger. Stærke sanselige Lidenskaber
boede ved Siden af høj intellektuel Stræben; udpræget
Selvfølelse, der ofte kunde give sig bombastisk Udtryk,
udelukkede ikke selvforglemmende Begejstring og Opgaaen
i de store Idéer, de uendelige Horizonter, der havde aabnet
sig for ham. Dette sidste Træk belyser han paa flere Steder
paa billedlig Maade. Ligesom Aktæon frejdig gik paa Jagt
med sine Hunde, men inden Aften selv blev Hundenes
Bytte, saaledes kan det gaa Tænkeren, naar han glæder sig
over sine Tanker; de kan tilsidst faa saadan Magt over
ham, at ethvert Selvhensyn forsvinder. En Dreng løser i
Kaadhed en Baad ved Strandbredden, men Strømmen kan
føre ham ud paa det store Hav, og han maa erfare, at der
er større Magter til end dem, han kan tumle i en Leg. Gennem
andre Billeder anskueliggøres de modstridende Følelser
hos den, der stræber efter Erkendelse: snart føler han
sig ringe og ussel overfor den høje Genstand for hans Erkendelse,
idet han føler sin Tankes Svaghed; snart løfter
han igen sit Hoved i heroisk Stræben, der fører ham ud
over alle hindrende Skranker. Den heroiske Sjæl vil hellere
segne i Kampen, eller dog føle Smerten ved egen Ufuldkommenhed,
end han vil slaa af paa det ideale Maal. Bruno har
idethele Sans for sjælelige Tilstande, der betinges ved
stærke Modsætninger. Her staar han --- ligesom paa sin Vis
Shakespeare --- som en moderne Psykolog, der ikke, som de
antike Tænkere dvæler ved enkelte klare Karaktertræk eller
er bange for, at sjælelig Harmoni ikke skulde kunne bevares
trods stærke, stridende Modsætninger. Det var ikke
blot den ydre Verden, der blev større i Renæssancen; ogsaa
den indre Verden blev rigere og dybere.

Naar Bruno kun i store Træk gav Udkast til det Verdensbillede,
der i de følgende Aarhundreder blev gennemført i
% --- p. 58 ---
\setcounter{footnote}{0}%
\opage{58}det enkelte, hænger det sammen med, at hans Fantasi er
nok saa stærkt udviklet som hans Tænkning. Og dog er
han klar over, at Erfaring og Forsøg maa træde til, for at
de dristige Tanker stadfæstes. Han vil netop aabne Øjnene
for Nødvendigheden af nye Undersøgelser. Han er ogsaa
klar over, at selv den dristigste og alvorligste Forsken ikke
vil kunne udtømme Tilværelsen.

\piecerule

% --- p. 59 ---
\setcounter{footnote}{0}%

\opage{59}\essayhead{SPINOZA}

\essaysource{24. NOVEMBER 1632 --- 21. FEBRUAR 1677}

\romanhead{I}
% CHECK roman numeral read

I de Forelæsninger, Henrik Steffens holdt ved Universitetet
i København 1802, og som fik saa stor Betydning for
dansk Aandsliv, gav han en begejstret Skildring af Spinozas
Liv og Karakter og udtalte derefter: „Aldrig kan jeg tænke
mig denne herlige Mand i hans rolige Ensomhed, med et
Lys i hans Hjerte, som intet Tilfælde kunde bortjage, med
en indvortes Fred, som ingen Forfølgelse kunde forstyrre,
uden at en Straale af hans Lys vækker det skønneste Haab
hos mig.“ --- Hvad Steffens her har udtalt, vil alle, hvorledes
de end stiller sig til Spinozas Lære, blive enige i. Den
inderlige Fordybelse i Tankearbejdet, --- den Trofasthed,
med hvilken den vundne Sandhed fremstilles uden Hensyn
til Mængdens og de Mægtiges Meninger, --- den Renhed i
Sindet, der satte det aandelige Gode, fri Erkendelse er, over
de ydre Goder, som saa let drager mange til sig, --- alt
dette vil vække Beundring overalt, hvor forudfattede Meninger
ikke raader ved Bedømmelsen af Menneskers indre
Liv. Men det er ikke nok til at begrunde den Begejstring,
med hvilken saa mange fra alle Sider i Aandens Verden
fejrer enhver Mindedag, der fører Tanken tilbage til ham.
Den 21. Februar 1677 døde han i sin tarvelige Bolig i Haag,
og naar Mindet om ham i den Anledning fremdrages rundt
om i hele Verden, er det ikke blot paa Grund af hans rene,
ideale og uegennyttige Sindelag. Naar det for nogle Aar
siden stiftede internationale Selskab Societas Spinozana
køber det Hus, han boede i ved sin Død, for at gøre det til
et Mindesmærke om ham ved Siden af den Billedstøtte, der
% --- p. 60 ---
\setcounter{footnote}{0}%
\opage{60}for 50 Aar siden rejstes ligeoverfor det, saa er det, fordi
der af hans Ensomhed er fremgaaet et af den menneskelige
Tankes ypperste Værker, hans „Ethica“, som i kort og kraftig
Form fremsætter Idéer, i hvilke den nye Filosofis og
den nye Naturvidenskabs Grundtanker var indvævede.

Historisk staar Spinoza som Fuldender af Renaissancens
Tænkning. Han forener Giordano Bruno's Naturbegejstring
og store Horizont med Fordring om den strenge Begrundelse
af hver Enkelthed, som Galilei, Kepler og Descartes
havde hævdet. Den Enhed i Naturen, der for Renaissancen
blev Et med Guddommen, havde ifølge Spinoza sin Begrundelse
gennem den indre Sammenhæng, Videnskaben
søger at paavise mellem alle de enkelte Ting, den undersøger.
Det er ingen Forklaring af et Fænomen eller en Begivenhed
--- hævder Spinoza --- at finde deres Aarsag i Gud,
ti Gud antages at være Aarsag til alle Ting, og ingen Forklaring
af den enkelte Ting kan derfor umiddelbart udledes
af ham. Den enkelte Ting eller Begivenhed maa derimod
paavises som Led i en Række, indenfor hvilken ethvert
Led forudsætter de foregaaende Led. Forklaringen af en
Enkelthed kommer da til at ligge i dens Forhold til andre
Enkeltheder. Hvor der kan dannes en saadan Række, kundgør
sig en Lov. Og gennem Lovenes indbyrdes Sammenhæng
naas der til Ideen om en Guddom, der virker i Tingenes
Indre gennem den Lovmæssighed, der betinger deres
Væsen og deres Existens. Hvad der hos Renaissancens tidligere
Tænkere var en begejstret Skuen, det blev for Spinoza
gennem Studiet af den nye Naturvidenskab, som var
det syttende Aarhundredes store Gave til Menneskeheden,
en Tanke, der kunde lægges til Grund ved videnskabeligt
Arbejde. Spinoza udtaler ligefrem: „Jo mere vi forstaar
de enkelte Ting, des mere forstaar vi Gud“, --- en Sætning,
der var Goethe's, denne store Spinozadyrkers Yndlingssætning.
Hvad Spinoza udtaler som Led i hele sit System,
minder, tildels i Ordlyden, om Sætninger hos Kepler og
Galilei. --- Kepler havde (i sit berømte Skrift om Planeten
Mars) fordret, at de Aarsager, man beraabte sig paa i Vi%
% --- p. 61 ---
\setcounter{footnote}{0}%
\opage{61}denskaben, selv skulde paavises i Erfaringen, ikke blot være
Genstande for Tro eller Spekulation. Han berigtigede derved
sine egne tidligere mystiske Forklaringer. Det Udtryk,
Kepler bruger om en i Erfaringen paaviselig Aarsag (vera
causa), forekommer ogsaa hos Spinoza. Det var paa den
Tid et teknisk Udtryk, og bruges endnu af Newton. Det
udtaler de naturlige Aarsagers Princip. --- Galilei hævdede
(i sine Dialoger om astronomiske Verdenssystemer) overfor
skolastiske Modstandere, at det ikke kunde nytte at beraabe
sig paa Gud som Aarsag til et enkelt bestemt Fænomen,
hvis Tydning Astronomerne var uenige om, ti Gud
kunde i sin Almagt lige saa let lade Solen bevæge sig om
Jorden som omvendt. Og i Analogi dermed lærer Spinoza,
at man i Videnskaben ikke kan paaberaabe sig et Væsen,
der antages at ligge til Grund for Alt, som Aarsag til en
enkelt bestemt Ting. Hver enkelt Ting skal forklares ved
dens Sammenhæng med en anden Ting, denne igen ved
dens Sammenhæng med en tredie Ting, og saaledes i det
uendelige. (Ethica I, 28).

Hvad selve Aarsagsforholdet angaar, hævdede Spinoza, i
Tilslutning til Descartes, at der ikke kan være mere i Virkningen,
end der var i Aarsagen. Hvis der var mere, vilde
der jo være noget, man ikke havde fundet Aarsagen til!
Aarsagsforholdet staar for Spinoza som et Identitets- eller
Ækvivalensforhold mellem forskellige Fænomener. Denne
Opfattelse var idethele karakteristisk for det syttende Aarhundredes
Tænkere. Den udvikles nærmere af Huyghens
og Leibniz, og den blev i det nittende Aarhundrede af Betydning
for Opdagelsen af Loven om Energiens Bestaaen.
Robert Mayer, den første i Rækken af de Mænd, der hver
ad sin Vej opdagede denne Lov, skal stadig have ført Sætningen
causa æqvat effectum i Munden; den stod i hans
Øjne som selvfølgelig, og den ledte ham ved hans afgørende
Iagttagelser og Forsøg. Ogsaa de andre Naturforskere,
der selvstændig opdagede Energiens Bestaaen, gik ud fra
% PRINTED AS IS (read at the image, odd spots o1)
Tanken om et Ækvalensforhold mellem Aarsag og Virkning
som selvfølgeligt. Men Selvfølgeligheden skriver sig
% --- p. 62 ---
\setcounter{footnote}{0}%
\opage{62}først fra det syttende Aarhundrede, og Spinoza er den, der
tydeligst har udtalt den.

Sammensmeltningen af Renaissancens Mystik og den nye
Naturvidenskab blev for Spinoza Grundlag for en Lære om
en evig Bestaaen, en absolut Substans, som han kaldte
Gud eller Naturen, der ligger til Grund for Alt og aabenbarer
sig baade i den indre Kraft, med hvilken hvert enkelt
Væsen søger at hævde sin Bestaaen, og i det Baand, der
knytter alle Væsener i Verden sammen til et stort Hele.
Vor Erkendelse begynder ganske vist med Opsamling af
enkelte Iagttagelser, men gennem Sammenligning mellem
disse og gennem Bestemmelse af deres indbyrdes Sammenhæng
kan den indre Grund til Alt findes, og Højdepunktet
for al Erkendelse var for Spinoza en Skuen af det Enkelte
i dets rationelle Sammenhæng med det Hele, --- en afsluttende
Intuition. Det enkelte i dets Isolation, uden Sammenhæng
med Naturhelheden var for ham kun en Abstraktion.
Men den afsluttende Intuition naas kun gennem strengt
Arbejde med Enkelthederne, og Spinoza erklærer selv, at
det kun er meget lidt, han er naaet til at forstaa ad denne
Vej\footnote[1]{I \textit{Den nyere Filosofis Historie} har jeg givet en udførligere Karakteristik af Spinoza's Filosofi og i \textit{Spinoza's Ethica. Analyse og Karakteristik} (Videnskabernes Selskabs Skrifter. 1918) har jeg gennemgaaet hans Hovedværk.}.

\romanhead{II}
% CHECK roman numeral read

Men hvorledes kom Spinoza til denne Filosofi? --- Dette
Spørgsmaal fører til at fremdrage hans Biografi og hans
personlige Udviklingsgang.

Spinoza blev født i Amsterdam den 24. November 1632.
Hans Forældre var portugisiske Jøder, der hørte til de saakaldte
Marannos, som under det spanske Herredømme var
blevne tvungne til at opgive deres jødiske Kultur og deres
hebraiske Sprog. Hans Forældre var som saa mange andre
Marannos flygtede til Amsterdam, hvor man først oversaa,
siden taalte deres Religionsøvelse. Spinozas Fornavn har
% --- p. 63 ---
\setcounter{footnote}{0}%
\opage{63}været Bento; med dette portugisiske Navn underskrev han
sig paa juridiske Dokumenter. Som Student ved Talmudskolen
i Amsterdam kaldtes han Baruch, og i Breve og i et
af sine Skrifter kaldte han sig Benedictus (ligesom Bento
svarende til det hebraiske Baruch). Han studerede med Iver
de hebraiske Religionsbøger og Talmud, og man ventede
sig en fremragende Talmudlærer af ham. Men efterhaanden
vakte han en pinlig Opmærksomhed hos sine Trosfæller
ved sin historiske og filosofiske Kritik, og tilsidst blev han,
24 Aar gammel, højtidelig udstødt af Synagogen. Nu først
gjorde han ret Bekendtskab med Vestens nye Videnskab
og Filosofi, der, som vi har set, har sat dybe Spor i hans
egen Filosofi. Da ogsaa den protestantiske Gejstlighed forargedes
over hans Kritik af det gamle Testamente, fik man
Øvrigheden til at udvise ham af Amsterdam. Siden levede
han forskellige Steder paa Landet og ernærede sig ved at
slibe optiske Glas, hvori han opnaaede stor Færdighed. Tilsidst
opslog han sin Bolig i Haag, og her døde han af Brystsygdom
den 21. Februar 1677. Hans rene og harmoniske
Livsførelse under tarvelige Forhold, hans Uegennyttighed og
ikke mindst hans rolige Død var for Tidens Teologi et stort
Problem. Kunde en Mand med hans Anskuelser leve et ordentligt
Liv og dø i Ro? Der var mange Rygter i Omløb
om hans udsvævende Liv og om hans Død i Fortvivlelse.
Men en brav Præst, der kom til at bo i Spinozas Lejlighed
og hos de samme Værtsfolk, har i en Biografi af Filosofen
sat sig for at gendrive disse Rygter. Det er et helt lille hollandsk
Genrebillede, man faar ved hans Skildring, der grundede
sig paa Beretninger af Spinoza's Vært og Værtinde,
som holdt meget af deres Logerende. Spinoza kom ofte ned
til dem og samtalede med dem ved en Pibe Tobak. Endnu
om Morgenen paa den Dag, der skulde blive hans sidste,
var han nede hos dem og samtalede med dem, og de var
meget forundrede, da de kom hjem fra Kirke om Eftermiddagen
og fandt ham død. Han var død i Ro, i Nærværelse
af en af sine Venner, der var Læge. Den gode Præsts Vidnesbyrd
er saa meget des vigtigere, som han nærede den
% --- p. 64 ---
\setcounter{footnote}{0}%
\opage{64}største Afsky for Spinoza's Lære og var meget forarget
over, at der kom Regninger til Boet, der omtalte ham som
„salig Hr. Spinoza“.

Lad os vende tilbage til Tiden lige efter Udstødelsen af
Synagogen. Vi véd meget lidt om de Aar, der forløb fra
dette Tidspunkt, indtil hans Brevvexling med forskellige
Venner begyndte (1656---1660). Det var dog dels den Tid,
da hans filosofiske Ideer modnedes, dels en Gæringsperiode,
hvad hans personlige Maade at tage Livet paa angaar. Ad
begge Veje forberedtes den vidunderlige Samklang mellem
Lære og Liv, der stiller ham i første Række af dem, der har
arbejdet for fri Forsken og ideal Livsførelse.

Spinoza har i denne Periode foretaget Overgangen fra
Teologi til Filosofi. Som rabbinsk Teolog havde han i
Gudsbegrebet, saaledes som den hebraiske Religion opfattede
det, fundet et Princip til Forstaaelse. Gud indeholdt
for ham Forklaring til alle Ting. Den Forandring, der nu
foregik med ham, kan udtrykkes ved, at medens Gud forhen
for ham var Forstaaelsens Princip, var det nu Forstaaelsens
Princip, der for ham blev Gud. Overgangen
skete ved, at han spurgte, hvad det egentlig er at forstaa
noget. Svaret fandtes i den nye Videnskab og Filosofi, som
han nu gjorde Bekendtskab med. Det gjaldt at finde saadanne
Aarsager til Tingene, som selv kundgjorde sig i Erfaringen.
Det var Keplers, Galilei's og Descartes's Principper,
han optog i sin Filosofi. I Forudsætningerne for saadan
Aarsagserkendelse er Princippet for al Forstaaelse givet,
og dette Princip bestemte væsentlig hans Gudsbegreb. Han
har sikkert følt Modsætningerne mellem de to Gudsbegreber,
mellem hvilke Overgangen skete, meget skarpt. Men
uagtet han i et særligt Skrift (Tractatus theologico-politicus)
satte sig for at indskærpe Forskellen mellem Teologi
og Filosofi, mente han sig dog berettiget til at beholde det
samme Ord, idet han nu bruger Ordet Gud enstydig med
Ordene Substans og Natur. Han mente sig at have et Begreb
om Gud, selv om han ikke kunde have noget tilsvarende
Billede. Hans Gudsbegreb var baade Princippet for
% --- p. 65 ---
\setcounter{footnote}{0}%
\opage{65}Tilværelsens Bestaaen og for dens Forstaaelighed. Det
fyldte hans Sind med den højeste Glæde, at Tilværelsen i
saa høj Grad var tilgængelig for videnskabelig Forstaaelse.
Aandslivets Højdepunkt var for ham, hvad han kaldte den
intellektuelle Kærlighed til Gud (amor intellectualis dei),
som fremkaldtes ved Tanken om, at sand Erkendelse var
mulig i denne Verden, der har sin Grund i Gud som den, der
er Et med Naturen. Det at filosofere, var for ham en religiøs
Gerning.

Men de samme Overgangsaar, i hvilke disse Ideer modnedes
i ham, har tillige været optaget af indre Kamp af en
anden Art end de intellektuelle, Kampe, der dog ogsaa
førte ham til i Tænkearbejdet at se det højeste aandelige
Gode. Der har efter hans egen Bekendelse været en Tid,
hvor dels Trangen til sanselig Nydelse, dels Trangen til
Ære og Rigdom har rørt sig stærkt hos ham. Han har følt
sig dragen til forskellige Sider. Men han erfarede under
disse Kampe, at der kom Fred i Sindet, naar han blot
kunde samle sig i alvorlig Eftertanke (modo possem penitus
deliberare!). I sit eget personlige Liv har han saaledes
erfaret Forstaaelsens Betydning, og Erkendelsesglæden har
han lært at kende baade gennem Tankernes og gennem
Drifternes Brydning.

Man har ment, at hans Opdragelse i sit Folks Religion
har været afgørende for ham i hele hans Liv, og at det var
Efterklangen af hans religiøse Barndomstid, der senere
virkede til at give hans Livsførelse et Præg af Renhed og
Ophøjethed. Men der kan ingen Tvivl være om, at det Liv
i Tankens Verden, hvilket han gennem indre Kampe naaede
op til, har havt en afgørende Indflydelse for hans Liv og
vist sin Bærekraft i dette. Naar en Søfarende kommer fra
den sydlige Halvkugle op over Ækvator, er det nye Stjerner,
der leder ham, selv om han (maaske med Pietet) mindes,
at andre Stjerner forhen bestemte hans Fart. --- Jeg
kan her ikke lade være at tænke paa min Lærer Professor
Hans Brøchner (som iøvrig har forfattet en grundig Monografi
over Spinoza 1857). Han staar ogsaa som et smukt
% --- p. 66 ---
\setcounter{footnote}{0}%
\opage{66}Vidnesbyrd om, at der kan leves et ædelt Menneskeliv i
Lyset af de nye Stjerner, Tankearbejdet har ført til at opdage.

\romanhead{III}
% CHECK roman numeral read

For en Prøvelse ud fra vor Tids Erkendelsesteori staar
Spinoza's Filosofi nærmest som en i dogmatisk Form given
Fremstilling af videnskabelig Erkendelses Ideal, ikke som
en afsluttende Løsning af Problemer. Men den, der vil fordybe
sig i hans Værk, vil under den faste, dogmatiske Formulering
mærke en arbejdende Tanke. Derfor kunde Lessing,
en af de første, der forstod Spinoza's Betydning, erklære,
at hvis han skulde kalde sig efter nogen, vilde han
kalde sig efter Spinoza, uagtet han selv havde udtalt den
berømte Sætning, at ikke Besiddelsen af Sandheden, men
den stadige Stræben efter Sandhed var det højeste, vi Mennesker
kan naa. Spinoza har selv sagt, at der staar mange
Gaader tilbage, selv om man gaar ind paa hans Filosofi. Det
var jo meget faa Emner, han mente at have forstaaet paa
fuldkommen Vis. Og han kunde ikke tænke sig, at Tilværelsens
Fylde kunde være udtømt ved de eneste to Former
for Existens, vor Erfaring viser os, den aandelige og den
materielle, hvis indbyrdes Forhold frembyder saa mange
Vanskeligheder. Der maa, mente han, være utallige andre
Former for Tilværelse, selv om den Erkendelse, vi formaar
at vinde af de Tilværelsesformer, der er os tilgængelige, er
nok til at vi kan komme i inderlig Forbindelse med Tilværelsens
inderste Grund.

Det var Tankelivet, som for Spinoza var det højeste, det,
der gav alle andre Sider af Livet Værdi. Fra sin Ungdom
af havde han jo erfaret, hvilken Betydning Tankelivet kan
faa, naar det gaar i Dybden. Og han staar som et stort Forbillede
ved sin inderlige Hengivelse til intellektuelt Arbejde.
Der er andre Sider af Livet, der i vore Øjne ikke
kommer til deres Ret hos ham. Saaledes Kunst, Poesi og
det frie Følelsesliv, der i sine dybeste og inderligste Former
kan kaldes Religion. 

% --- p. 67 ---
\setcounter{footnote}{0}%
\opage{67}Alligevel havde han Evne og Lyst til at iagttage Menneskelivet,
som det viste sig i ham og omkring ham. Ligesom
en hel Række af Renæssancetidens Tænkere fandt han i
Selvopholdelsestrangen, i Trangen til at fastholde og uddybe
sit Væsen, Grundlaget for alt menneskeligt Sjæleliv.
I alle Væsener virker jo efter Spinozas Tro det i sig bestaaende
og ved sig selv bestemte uendelige Væsen, han
kalder Gud eller Naturen, og det er en Straale fra dette
guddommelige Væsens Fylde, der rører sig i de endelige
Væseners Trang og fører til Selvhævdelse. Lyst og Ulyst,
Glæde og Sorg, Haab og Frygt, Kærlighed og Had er forskellige
Ytringsformer for den Selvopholdelsestrang, der
rører sig i os alle. Og Spinoza giver en ikke blot for hans
Tid mesterlig Skildring af disse Følelsers Udvikling. Mennesket
kan blive omtumlet af Lidenskaberne ligesom et
Skib paa et oprørt Hav, hvis det ikke søger en klar Forstaaelse
af dem og deres Aarsager. Gennem saadan Selvforstaaelse
kan der udvikle sig en Aandsstyrke, en Tapperhed
(fortitudo) overfor Livet, der dels fremtræder som
Livsmod (animositas), dels som det Højsind (generositas),
der tager andres Byrder op foruden sine egne. I disse forskellige
Former ytrer den sande Frihed sig, der ikke er en
mystisk Evne, men en Karakteregenskab, som forudsætter,
at Mennesket har indre Betingelser for at bestaa og udfolde
sig. Naar Selvforstaaelse bliver til videnskabelig Erkendelse,
lægges Grunden til den sande Frihed saa meget
fastere.

Spinoza havde en for hans Tid sjelden Forstaaelse af
andre religiøse Standpunkter end hans eget. Gennem sit
historiske Studium af sit Folks Religion blev han Forgænger
for moderne religionshistorisk Forsken. Han saa de
positive Religioners Betydning deri, at de i deres Gudsforestilling
fremholdt et Livsideal (exemplar vitæ), i Lydighed
mod hvilket der kunde føres et godt og retfærdigt Liv.
Det afgørende overfor religiøse Forestillinger var for Spinoza,
der herved staar som Forgænger for Mænd som Carlyle
og Kierkegaard, at den Enkelte med fuld og redelig
% --- p. 68 ---
\setcounter{footnote}{0}%
\opage{68}Tilslutning (sine ulla hæsitatione, sed integro animi consensu)
følger det Livsideal, han anerkender. Det var efter
hans Biografs Udsagn i denne Aand, Spinoza samtalede
med sine Værtsfolk, naar han drøftede med dem, hvad
deres Præst havde sagt i sine Prædikener. --- I et særligt
Skrift (den ovenfor nævnte Tractatus Theologico-politicus)
søger han at vise, at naar Religionen opfattes paa denne
Maade, vil Staten uden Fare kunne give fuld Religions- og
Tænkefrihed. Her som paa andre Punkter har han udtalt
Ideer, der peger langt ud over hans Tid.

\romanhead{IV}
% CHECK roman numeral read

Trods Alt er Spinoza et Barn af sin Tid. Hans Tænkning
omfatter i koncentreret Form Problemer, som nu maa behandles
hvert for sig, for at der saa bagefter kan prøves,
hvorvidt og hvorledes Løsningerne af de enkelte Problemer
kan bringes i Overensstemmelse med hverandre. Det
er Arbejdets Deling paa det intellektuelle Omraade, og paa
Aandslivets Omraade idethele, der gør, at vi nu maa filosofere
anderledes end Spinoza. Men det vil bestandig føre til
Berigelse og Befæstelse at gaa tilbage til den stille Tænker i
det syttende Aarhundrede, den Tid, den moderne Videnskab
blev grundlagt. Og det styrker Troen paa Tankearbejdets
Bærekraft og Betydning, naar man dvæler ved Billedet
af ham i hans simple Bolig, arbejdende paa sit monumentale
Værk under Misforstaaelse, Ligegyldighed og
fanatisk Forfølgelse. Den hellige Ild brændte i hans Indre
og lod ham glemme Kulden og Mørket udenfor. Den Aand,
der besjælede ham, maa ikke forsvinde fra Aandens Verden,
saa meget end baade Opgaverne og Midlerne til deres
Løsning skifter.

Hvad der bestaar i Tankens som i Naturens Verden, det
er ikke de Former, som til given Tid stod til Raadighed,
men det er de Kræfter, som laa til Grund ved disse Formers
Opstaaen, og som for en Tid, men ogsaa kun for en Tid,
% --- p. 69 ---
\setcounter{footnote}{0}%
\opage{69}lagde sig til Hvile i dem. Til dem søger vi tilbage for at
styrke vort Haab til Fremtiden. Selv om vi nærer en fast
Tro paa, at der stedse paany vil opstaa aandelige Værdier,
kan vi ikke paa Forhaand afgøre, af hvad Art disse Værdier
vil blive. Men vi kan faa en Anelse derom ved at gøre
os fortrolige med den Aand, i hvilken Fortidens Store virkede.

\piecerule

% --- p. 70 ---
\setcounter{footnote}{0}%

\opage{70}\essayhead[Pascal og Kierkegaard]{PASCAL OG KIERKEGAARD\footnote[1]{I Anledning af Trehundredaaret efter Pascal's Fødsel (19. Juni 1923) udgav Revue de Métaphysique et de Morale et Pascalnummer, til hvilket jeg efter Opfordring af Redaktionen skrev denne Sammenligning mellem de to religiøse Tænkere, som samtidig udkom paa dansk (i „Tilskueren“).}}

Med to Aarhundreders Mellemrum mødtes to Tænkere,
som begge med hele deres Sjæls Inderlighed klamrede sig
til Kristendommen, i den Maade, paa hvilken de gjorde op
med deres Religion, saaledes som denne havde udviklet
sig gennem Tiderne. De hævdede begge, at der var en uforsonlig
Modsætning mellem den Aand, i hvilken Kristendommen
traadte ind i Verden, og de nyere Tiders Kristendom.
Modsætningen var opstaaet ved, at man havde forsøgt
at træde i intimt Forhold til „denne Verden“, det vil
sige, til den paa rent menneskelig Grund udviklede Kultur.
Opgøret blev af dem begge foretaget med en saadan Lidenskab,
og gennemført med en saadan Konsekvens, at der
vilde have været Stof nok til en stor Revolution i Aandens
Verden. Men der knytter sig den Tragedie til dem begge,
at de brød sammen, da de havde sagt deres sidste Ord. Den
aandelige Kamp — i deres Indre og med deres Omgivelser
--- havde fortæret dem. De gik bort som unge, Pascal 39
Aar gammel (1662), Kierkegaard 42 Aar gammel (1855). De
efterlod sig en stor Gaade, skarpere formuleret end nogensinde
før, men uløst som før. Dog, en Heltegerning i Aandens
Verden øves aldrig forgæves. De har med Fynd og
Klem stillet det Spørgsmaal, om de kristne Kirkers religiøse
Overleveringer stadig kan lægges til Grund for Tanke og
Liv af oprigtige og redelige Sandhedssøgere.

Fra Kirkernes Side er der ikke givet Svar paa dette
% --- p. 71 ---
\setcounter{footnote}{0}%
\opage{71}Spørgsmaal. Man mener at kunne fortsætte sin Praxis uden
at tage Hensyn til det. Man støtter sig til, hvad der jo ogsaa
er en Kendsgerning, at der hidtil ikke er opstaaet nogen
anden Organisation, der formaar i større Kredse at vække
til alvorlig Livsførelse og tillige at bringe Trøst og Lindring
under Livets haarde Kaar. Men Spørgsmaalet bringes
herved ikke ud af Verden. Det fortælles om Scipio Afrikanus,
at da han blev anklaget i Folkeforsamlingen for
Uorden i sine Regnskaber, sagde han: „Paa denne Dag
slog jeg Hannibal ved Zama; --- lad os gaa til Kapitolium
og takke Guderne!“ Folket fulgte ham --- og glemte Regnskabsopgørelsen.
Den Tvivl staar stadig ved Magt, --- og er
efter de to Tænkeres Tid bleven indgaaende begrundet ved
kritiske Undersøgelser, --- om Kirkerne har Ret til at mene
sig i Kontinuitet med den dybe Aandsbevægelse, der for
1900 Aar siden brød frem i en Krog af det store Verdensrige,
og hvis ledende Tanke var, at Livet nu maatte føres
i spændt Forventning om et snart indtrædende overnaturligt
Frembrud af et nyt Rige, en Forventning, i Forhold til
hvilken Livet i Familie og Stat, i Kunst og Videnskab blev
af forsvindende Værdi, forsaavidt det ikke var en positiv
Hindring.

Det er dog ikke blot for dette Opgørs Skyld, saa betydningsfuldt
det end var, at vi mindes de to Mænd. Paa deres
Vej til Opgøret har de udviklet Tanker af stor Betydning,
og --- som saa ofte paa Tænkningen Omraade --- er Vejen
ikke mindre betydningsfuld end Maalet. De har undervejs
kastet Lys over det menneskelige Aandslivs Natur og Vilkaar.
De har gjort en Mængde betydningsfulde psykologiske
Iagttagelser. De har ydet Bidrag til, hvad vi nu vilde
kalde Erkendelsesteori, idet de har prøvet Tankens Kraft
paa selve Tankens Grænse. De har vist, at selv hvor Følelse
og Lidenskab fører Ordet, er det dog Tanker, der som
Ledestjerner bestemmer den sjælelige Bevægelses Retning.
De er Tankens Riddere, og deres Kamp har Interesse, hvad
enten man mener, at den har ført til Sejr eller til Nederlag.
--- Deres Veje faldt ikke sammen. Det er to forskellige
% --- p. 72 ---
\setcounter{footnote}{0}%
\opage{72}Personligheder, vi her har for os, hver med sine ejendommelige
indre og ydre Forudsætninger, og de historiske Forhold,
under hvilke deres Tanker tilspidsedes, var højst forskellige,
skønt Analogier ikke mangler.

Allerede i deres Stil er de forskellige. Det maa her vel
erindres, at Pascal's „Pensées“, hans vigtigste Værk, ikke er
et afsluttet Arbejde, men en Række Udkast, som han har
nedskrevet i Mellemrummene mellem sine Sygdomsanfald,
ofte afbrudt af disse. De er kortfattede, sammentrængte,
vibrerende af en Sindsbevægelse, der snart lokkede Tankerne
frem, snart fremkaldtes af dem. Der er en enestaaende
Kraft og Klarhed over disse Tanker, og disse Egenskaber
vilde næppe være gaaede tabt, hvis Bogen var bleven
udarbejdet. Kierkegaards Stil har ikke denne slaaende
Kraft, denne epigrammatiske Karakter. Forfatterskabet var
for ham, i hans første Periode i hvert Fald, en Kur: saalænge
han skrev, fik Tungsindet ingen Magt over ham. Han
tog sig derfor god Tid. Og desuden elskede han Sproget
for dets egen Skyld; han legede med det og tænkte ikke
altid paa, at Læseren maaske ikke havde saa god Tid som
han, men længtes efter at kunne overse Tankeløbet og se,
hvor det fører hen. Episoder af lyrisk og naturbeskrivende
Art indflettes, saavel som Exempler fra det daglige Liv.
Den Lidenskab, der ogsaa ligger bag ved hans Produktion,
mærkes mere som en stadig Skælven end som øjeblikkelig
Sindsbevægelse. Først i hans Forfatterskabs allersidste Periode
bliver denne Skælven til et Jordskælv. ---

Hvis jeg i dette Forsøg paa en Sammenligning skulde
komme til at dvæle forholdsvis meget ved den danske Tænker,
er det, fordi jeg skriver for franske Læsere. At ville
give Bidrag til Forstaaelse af Pascal i et fransk Tidsskrift
vilde være at bringe Ugler til Atén, saa rig og grundig er
den franske Pascalforsken. Jeg ønsker her særlig at udtale
min Beundring for Ste Beuve's „Port Royal“, som jeg nu
efter næsten 60 Aars Mellemrum har fornyet Bekendtskabet
med, og som, trods det Arbejde, der er gjort i Mellemtiden,
fremtræder med samme Friskhed og Grundighed som for%
% --- p. 73 ---
\setcounter{footnote}{0}%
\opage{73}hen. Desuden maa jeg nævne den pragtfulde Udgave af
Pascal's samlede Værker, som Léon Brunschwicg har besørget
under Samarbejde med Pierre Boutroux og Félix
Gazier, og som i sine Indledninger og Noter bringer mange
uundværlige Oplysninger. Navnlig giver Brunschwicg i Indledningen
til „Pensées“ en overordentlig interessant Redegørelse
for den Tankegang, Pascal kan antages at have
villet gennemføre i sin Apologi. Jeg citerer i det følgende
„Pensées“ efter den af Brunschwicg foreslaaede Ordning
af Fragmenterne. --- Kierkegaard's Værker udkommer for
Tiden i anden samlede danske Udgave. De foreligger ogsaa
paa Tysk (hos Diederichs, Jena). Værkerne udgør 14 Bind.
Men desuden er der begyndt en fuldstændig Udgave af
Kierkegaard's Dagbogsoptegnelser og Papirer (ved P. A.
Heiberg og V. Kuhr). Hidtil er 11 Bind udkomne, og der er
Materiale til 7 Bind endnu. Det er altsaa en meget omfangsrig
Produktion, der her er Tale om. Paa fransk foreligger
der en smuk Afhandling af Delacroix i Revue de
Métaphysique et de Morale (1900), og en Tale af mig paa
Københavns Universitet paa Hundredaarsdagen efter Kierkegaard's Fødsel
(5. Maj 1913) er trykt i samme Tidsskrift
(1913). Jeg tillader mig tillige at henvise til min Monografi
om Kierkegaard som Filosof (1892), der er oversat paa
Tysk.

Om nogen Indflydelse af Pascal paa Kierkegaard kan
der næppe være Tale. I Fortegnelsen over Kierkegaards
Bibliotek findes Pascal's Værker (i tysk Oversættelse). Men
det er først paa et sent Stadium af Kierkegaards Udvikling,
at Pascal findes omtalt i Dagbøgerne (efter hvad Hr. Arkivar
Heiberg godhedsfuldt har meddelt mig). Det var i Efteraaret
og Vinteren 1850, og da havde han allerede udgivet
de Skrifter („Sygdommen til Døden“ 1849 og „Indøvelse
i Kristendom“ 1850), hvis Indhold mest kunde minde
om Pascal. Han har saa ogsaa læst forskellige Værker
om Pascal, navnlig Reuchlin's Værk om Port-Royal, og er
bleven opmærksom paa Forskellen mellem ældre Udgaver
af „Pensées“ og den af Faugère besørgede Udgave, i hvil%
% --- p. 74 ---
\setcounter{footnote}{0}%
\opage{74}ken de mest revolutionære Udtalelser først forelaa. (De
tyske Oversættelser, der fandtes i Kierkegaard's Bibliotek,
var foretagne efter de ældre Udgaver).

--- Idet jeg nu forsøger en Sammenligning mellem de to
Tænkere, dvæler jeg successivt ved Temperament og Karakter,
ved intellektuelle Forudsætninger, ved Forholdet til
Kristendommen og ved Spørgsmaalet om, hvilke Veje der
stod aabne, da de to Mænd havde sagt deres sidste Ord.

\sectionhead{1. \textit{Temperament og Karakter.}}

Pascal's Karakter beskrives af hans to Søstre (Jaqueline
i hendes Breve, Mme Périer i hendes Biografi af Broderen)
paa overensstemmende Maade. Han havde et gærende eller
kogende Sind (humeur bouillant), og hans livlige Aand
medførte en saadan Utaalmodighed, at det undertiden var
vanskeligt at gøre ham tilpas; hans ældste Søster beundrede,
at han trods sit fyrige Sind var bleven bevaret for Ungdomsudskejelser.
Det er sikkert karakteristisk for ham,
naar han i sine Pensées (fr. 139) mener at have opdaget, at
Menneskenes hele Ulykke kommer fra en eneste Ting, den,
at de ikke kan holde sig i Ro. I den mærkelige „Discours
sur les passions de l'amour“ (som han sandsynligvis er
Forfatter af) prises Elskovslidenskabens Skønhed, netop
naar den ytrer sig paa en overvældende Maade, og Lidenskab
idethele beskrives især derved, at alle Tanker styrte
afsted i en eneste Retning. Der er her et Varsel om den
Troens Lidenskab, der greb ham efter hans definitive Omvendelse.
Hans Karakters Grundpræg bevarede sig, saa forskellige
Udviklingsstadier han end gennemløb.

Skønt der maatte en dobbelt Omvendelse til, før han kom
til sit sidste, afgørende Stade, har han dog aldrig været irreligiøs.
Gennem Opdragelsen var en alvorlig religiøs
Grundtone raadende i hans Sjæl; men ved Siden af den
rørte sig en tidlig vakt intellektuel Lidenskab, der førte
ham ind paa epokegørende matematiske og fysiske Undersøgelser,
og en Tid deltog han desuden i det verdslige Liv,
% --- p. 75 ---
\setcounter{footnote}{0}%
\opage{75}et aandfuldt og literært interesseret Samliv med begavede
Verdensmennesker. I denne verdslige Periode har han havt
den Oplevelse, der synes at ligge til Grund for Afhandlingen
om Elskovslidenskaben. Ifølge Strowski (Pascal et son
temps. II p. 278) synes Genstanden for hans Lidenskab at
have været en smuk og fornem Dame, der yndede eventyrlige
og urolige Stemninger, men ikke har givet ham Anledning
til noget, han maatte angre. Dog maatte det senere
staa for ham som en stor Synd, at hans Sind overhovedet
havde været optaget af andet end Gud. For Pascal var der
her en Modsætning som mellem Liv og Død. Han hørte til
de Karakterer, der kun gennem et afgørende Brud kunde
naa til deres endelige Stade; det fulgte af hans „humeur
bouillant“. Katastrofen kom den Nat (23. Nov. 1654), da
det i en ekstatisk Tilstand gik op for ham, at den levende
Gud kræver alt, saa at en ubetinget Forsagelse maatte træde
i Kraft. Hans foregaaende Liv stod nu for ham som et Liv
i Tomhed, men han jublede over, at han nu havde fundet
Livets Kilde.

Paa Baggrund af, hvad han saaledes havde oplevet,
maatte jesuitiske Skriftefædres Forsøg paa at lempe de
store Idealers Krav efter verdslige Menneskers Behov forarge
ham i højeste Grad. Det var umuligt at forbinde Gud
og Verden (joindre dieu au monde) (fr. 935). Verden vil
nok have Religion, men det skal være en blid og mild Religion
(fr. 956). I sine „Lettres écrites à un provincial“ og
i de sidste „Tanker“, han har nedskrevet, taler han i stedse
stærkere Udtryk om den fordærvelige Vej, Kirken er slaaet
ind paa. Kirken er ikke længer Kirke, hvis den ikke har
Gud med sig, naar den løser eller binder. (Si elle absout ou
si elle lie sans dieu, ce n'est plus l'Église. fr. 870.)

Under den Kamp, Pascal herved kom ind i, og medens
han var i fuld Gang med Udgivelsen af „Provincialbrevene“,
helbrededes en ung Slægtning af ham for en Øjensygdom,
Lægerne havde erklæret for uhelbredelig, derved at det
syge Øje berørtes af en Torn af Jesu Tornekrone, der hørte
til Relikvierne i Port Royal's Kirke. Dette Mirakel bekræf%
% --- p. 76 ---
\setcounter{footnote}{0}%
\opage{76}tede ham i hans Overbevisning om at være paa den rette
Side. Miraklernes Tid er ingenlunde forbi, som Modstanderne
mente. Og naar de kirkelige Myndigheder fornegter
Sandheden, maa Gud umiddelbart gribe ind (fr. 832; 839).
Tillidsfuldt appellerer derfor Pascal fra Rom, der fordømmer
hans Skrifter, til Jesu Christi Domstol (fr. 920).

I sit eget Liv gennemførte Pascal den absolute Forsagelse.
Vel gik han ikke i Kloster, men han levede sine sidste
Aar i den strengeste Askese, ene optagen af Kærlighedsgærninger,
men samtidig ængstelig afværgende, at nogen
af hans Nærmeste knyttede deres Kærlighed til hans Person
og derved berøvede Gud den fulde Hengivelse, der tilkom
ham. Sin langvarige Sygdom følte han som en Lettelse,
skønt den hindrede ham i at fuldføre det Forsvarsskrift for
Kristendommen, hvortil „Pensées“ indeholder Udkast. Aandelig
Sygdom er jo de Kristnes naturlige Tilstand, og det
er lettere at bevare sig i den, naar man tillige er legemlig
syg!

Det er et noget forskelligt Karakterbillede, der fremtræder
for os i Kierkegaard. Der er større Kontinuitet i hans
Udvikling, skønt den frembyder afgørende Overgange. Han
kan ikke regnes til „de to gange Fødte“. Hverken Ekstase
eller Mirakel spiller en Rolle i hans Liv.

Ligesom Pascal var han bleven opdraget i alvorlig Kristendom.
Men i sine Ungdomsaar var han optagen af Poesi
og Filosofi. Det var Romantikens Tid, og Ideen om en Harmoni
mellem Kunst, Videnskab og Religion beherskede
Sindene. Nylig var Romantikens Tankegang bleven systematiseret
af Hegel, og Kierkegaard har en Tid været optagen
af hans Ideer. Naar det var saaledes, at alle Modsætninger
tilsidst mødes i Harmoni, kunde man foreløbig
lade Stemninger og Fantasier frit udfolde sig; ogsaa de
hører jo med til den rige, mangfoldige Verden, som tilsidst
kunde finde sin Forklaring i de kristelige Grundtanker.
Kierkegaard har dog tidlig mærket, at her var en Dobbelthed
tilstede. Han har i sine Dagbøger fra denne Tid brugt
Goethe's Linier om Gretchen som en Slags Motto: „Halb
% --- p. 77 ---
\setcounter{footnote}{0}%
\opage{77}Kinderspiel, halb Gott im Herzen!“ Og der var endog en
Tid, da han blev dragen ind i et vildt Liv, som efterlod
angerfulde Minder, Minder, der maatte nære det Tungsind,
der laa dybt i hans Natur, og som undertiden kunde overvælde
ham saaledes, at han følte sig paa Grænsen af Sindssygdom.

Han saa det nu som sin Opgave at forme en harmonisk
Helhed ud af de Modsætninger, der rørte sig hos ham, ---
Legen og Alvoren, Letsindigheden og Tungsindigheden.
Han opdagede her ad sin egen personlige Erfarings Vej,
at det i Praxis ikke var saa let at naa til en „højere Enhed“,
som det var i Hegels System. Dette indre Arbejde gjorde
ham ensom. Han følte sig isoleret. Hans Tungsind gjorde
det (mente han) umuligt for ham at indtræde i Livsforhold,
der hvilede paa aaben og glad Hengivelse. Han hævede sin
Forlovelse med en ung, livsglad Kvinde, og han følte sig
ude af Stand til at paatage sig en Embedsstilling. Der var
for ham en grel Modsætning mellem andre Menneskers lette
Maade at tage Livet paa og hans egen Kamp med mørke
Magter. Hvis han kunde have levet i Middelalderen, vilde
han, mente han, være gaaet i Kloster. Det indre Arbejde
paa Udarbejdelsen af sit personlige Liv har han smukt
skildret paa følgende Maade: „Jeg sidder og lytter til Tonerne
i mit Indre, Musikens glade Vink og Orglets dybe
Alvor. At sammenarbejde dem, det er en Opgave, ikke for
en Organist, men for et Menneske, der i Mangel af større
Fordringer til Livet indskrænker sig til den simple, at ville
forstaa sig selv.“ --- Nu bor altsaa „Gott“ og „Kinderspiel“
ikke længere fredelig Side om Side i Hjertet. Modsætningen
mellem dem stiller en Opgave, der kan lægge Beslag paa
hele Livets Energi. Hvad der her fremtræder, gentager sig
hele Kierkegaards Liv igennem. Han har ikke bevæget sig
fra en Yderlighed til en anden. Modsætningerne har stedse
været der, men de skærpes, efterhaanden som hans Udvikling
skrider frem, og Arbejdet paa at overvinde dem
bliver derfor større og større. Af sin egen Erfaring kunde
han hente Maalestokken for det personlige Livs Højde.

Naar det gjaldt om at hæmme Tungsindet og lindre den
% --- p. 78 ---
\setcounter{footnote}{0}%
\opage{78}indre Spænding, var der et Middel, der stod til hans Tjeneste.
Han behøvede blot at følge sin Trang til literær Produktion.
Det var ham en Befrielse at hengive sig til poetisk
og filosofisk Produktion. Som Prinsessen i „Tusind og én
Nat“ frelste jeg, siger han, mit Liv ved at fortælle, det vil
sige, ved at skrive. I den første Periode af hans Forfattervirksomhed
(1843---1846) blev nogle af hans mest læste
Skrifter til, saaledes „Enten --- Eller“ og „Stadier paa Livets
Vej“. De angik, hvad Kampen i hans Indre stod om,
--- at blive Personlighed, at vinde et fast Midtpunkt, om
hvilke alle Livets Elementer kunde ordnes. Og derved kunde
disse Skrifter tillige tjene til at paavirke den slappe Tidsalder,
der var Romantikens og den spekulative Filosofis
Epigontid og var tilbøjelig til at forflygtige Livets Modsætninger
og derved sænke det aandelige Livs Niveau. Kierkegaard
turde ikke sige, at denne Paavirkning havde været
den oprindelige Hensigt med hans Forfattervirksomhed.
Han var saa ærlig at vedgaa for sig selv, at det havde været
en Spænding i hans Indre, der trængte til Afledning. Medens
Pascal i sine „Lettres à un provincial“ med fuld Bevidsthed
stillede sig den Opgave at kæmpe mod en Sænkning
af det etiske Niveau, kom Kierkegaard mere uvilkaarlig
ind paa et Arbejde i denne Retning. Men han var heller
ikke endnu, som Pascal, paa Højden af sin religiøse Udvikling.

En ny Periode i Kierkegaards Forfattervirksomhed (1847
—1855) fremkaldtes netop ved, at hans indre Liv antog en
stedse mere udpræget religiøs Karakter. „Jeg føler nu,“
skriver han i sin Dagbog (1847), „en Trang til i en dybere
Forstand at komme til mig selv ved at komme nærmere til
Gud\dots{} Der rører sig noget i mig, der tyder paa en Metamorfose\dots{}
Jeg skal derfor forholde mig stille.“ Man ser
her Modsætningen til Pascal's Ekstase: en halvt ubevidst
Udfoldelse af en ny Livsretning, overfor hvilken han staar
afventende. Aaret efter skriver han saa: „Nu er jeg ved
Troen i dybeste Forstand“. Kristendommen fik nu en Virkelighedsbetydning
for ham, som den hidtil ikke havde
% --- p. 79 ---
\setcounter{footnote}{0}%
\opage{79}havt. Efter at have afsluttet den første Periode af sin Forfattervirksomhed
troede han en Stund, at hans literære
Virksomhed nu var afsluttet, og han tænkte paa at trække
sig tilbage til en ensom Egn paa Landet. Men den stærkere
vakte religiøse Bevidsthed vakte saa en skarpere Kritik af
den bestaaende Kristenhed. I sin første Forfatterperiode
havde han taget Sigte paa Slappelsen i den individuelle
Maade at tage Livet paa; nu fandt han i Kirken en Gaaen
paa Akkord, en Lempelse af Kristendommen efter verdslige
Tilbøjeligheder, medens de oprindelige Idealer idethøjeste
fik en Plads i Baggrunden. Kristendommen var i
hans Øjne faktisk bleven afskaffet ved at blive gjort til
Mildhedens og Trøstens Religion. „Og nu lever der i den
bestaaende Kristenhed en forkælet, stolt og dog fejg, trodsig
og dog blødagtig Slægt, der lejlighedsvis hører disse
milde Trøstegrunde foredragne, men som knap ved, om
den vil gøre Brug af dem, selv naar Livet smiler skønnest,
og som i Nødens Stund, naar det ses, at saa milde er de dog
egentlig ikke, forarges.“ Denne Udtalelse findes i „Indøvelse
i Kristendom“, der tilligemed „Sygdommen til Døden“ er
karakteristisk for den anden Periode af Kierkegaards Forfatterskab.

Fra den officielle Kirkes Side tog man ikke Hensyn til
Paastanden om den skarpe Modsætning mellem den oprindelige
Kristendoms strenge Fordring om at gaa op i det
ene Fornødne og den moderne Kristendoms i Ly af Forsoningsdogmet
byggede Idyl. Da begyndte Kierkegaard sin
lidenskabelige Kamp mod den bestaaende Kirke, den voldsomste
aandelige Kamp, der er ført indenfor dansk Aandsliv.
% PRINTED AS IS (read at the image, odd spots o1)
Han sidste Ord var: „Det nye Testamentes Kristendom
er ikke til.“ Midt under Striden blev han syg og døde.

\sectionhead{2. \textit{Intellektuelle Forudsætninger.}}

Pascal dyrkede Matematik og Fysik og har et smukt
Navn indenfor disse Videnskaber. Han hentede fra disse
Studier sin Opfattelse af, hvad videnskabelig Erkendelse er.
% --- p. 80 ---
\setcounter{footnote}{0}%
\opage{80}En Systematik som Descartes', med Forsøg paa at føre Naturerkendelsen
tilbage til visse simple Sætninger og saaledes
naa et rationelt Grundlag for den menneskelige Fornufts
Arbejde, ansaa han for umulig. Man har træffende
kaldt ham experimental Positivist. Allerede i sin „verdslige“
Periode indsaa han, at den strengt matematiske Tænkning
(esprit de géométrie) ikke gennem sine Deduktioner kan
føre til fuld Forstaaelse af, hvad der foregaar i det menneskelige
Sjæleliv. Dertil er de sjælelige Fænomener for
mangfoldige, for sammensatte og for skiftende. Det maa
paa det sjælelige Livs Omraade gælde om at føre de mangfoldige
Fænomener tilbage til visse Typer, leve sig ind i enhver
af disse for sig, sætte sig ind i deres Motiver og Interesser.
Hertil egner matematisk Tænkning sig ikke. En fin
Sans for Typer og Nuancer (esprit de finesse) sætter ene
istand til at finde de væsentlige Elementer og forfølge dem
i deres Forbindelser. Mme Périer roser netop sin Broder for
hans Evne til at sætte sig ind i Forudsætningerne hos dem,
han henvender sig til. Han kendte, mener hun, „alle Drivfjedre
i det menneskelige Hjerte“. Denne Sans var bleven
udviklet hos ham i hans Omgang med Verdensmennesker,
og den kom ham til Gode senere, da det gjaldt om at vinde
Sjæle for den sande Religion. Spekulative Beviser fører
ikke til Maalet, indser han; hvad enten de er holdbare eller
ikke, er deres Virkning indskrænket til saadanne Øjeblikke
i vort Liv, da vi er rent intellektuelt interesserede; siden
glemmes de let. Hvad det gælder om, er at tage Menneskene
dér, hvor de staar, og saa lede dem til at komme videre ved
deres egen Evne; man overbevises nu engang bedst ved
Grunde, man selv finder. Man skal lede dem til at se, hvilken
underlig Blanding af Storhed og Usselhed deres Natur
indeholder. Tanken kan føre dem op i højere Regioner,
men Drifter og Lidenskaber drager dem atter ned i Smudset.
Mennesket skal føres til at se, hvilket ubegribeligt
Uhyre (monstre incompréhensible) det er. I den sande Religion
vil Mennesket finde Forklaring paa Gaaderne i sin
Natur og tillige Midler og Veje til at ophæve Spliden i sit
% --- p. 81 ---
\setcounter{footnote}{0}%
\opage{81}Indre. En Betragtning, Pascal særlig lægger Vægt paa, er
den, at Mennesket maa føle sig ensomt i den uendelige Verden,
fra hvis Dyb ingen Stemme høres, der kunde yde Forstaaelse
og vække Haab. Pascal selv var af Uro, Forundring
og Angst overfor de store Gaader ført til at gribe som eneste
Udvej den, som Kristendommen aabner. Kristendommen
indeholder ganske vist Mysterier ved Siden af de Aabenbaringer,
den bringer. Men en fuldstændig Afsløring af Tilværelsens
Hemmeligheder vilde udelukke det aandelige Arbejde,
den Kamp og det Mod til at vove, til at sætte Alt paa
ét Kort, der nu engang er Betingelsen for at være værdig
til det Højeste.

Der er ifølge Pascal to Hovedtyper, hvad de naturlige
Menneskers Forhold til de store Gaader angaar. Den ene
Type holder sig til Tanken og Pligten, til det Store og Faste
i Mennesket, og betragter alt Andet som ligegyldigt. Den
klassiske Repræsentant for denne Type er Stoikeren Epiktet.
Den anden Type tager Livets Bølgegang, som den kommer,
nyder de enkelte Øjeblikke og holder sig iøvrig til Skik og
Brug, idet den én Gang for alle er klar over alle Tings Forgængelighed.
Den finder en tryg Hvile i den Indsigt, at man
i Grunden Intet véd. Montaigne er den klassiske Repræsentant
for denne Type.

Studiet af disse to Typer var for Pascal en Forøvelse til
det kraftige Tag, med hvilket han i sine „Pensées“ griber
og ryster den Vantro for at bringe ham til det afgørende
Punkt, hvor der kun viser sig en eneste Udvej fra Fortvivlelse,
--- den, han selv personlig havde fundet. Det var ham
en Skuffelse, at der var saa liden Interesse for Studiet af
den menneskelige Natur, dens Typer og deres praktiske
Betydning. „Da jeg,“ siger han, „begyndte Studiet af Mennesket,
indsaa jeg, at de abstrakte Videnskaber ikke passede
for Mennesket, og at jeg fjernede mig mere fra min
sande Stilling i Verden ved at fordybe mig i dem end andre
Mennesker ved at være uvidende om dem.“ Saa havde han
ventet at finde mange Fæller i dette nye Studium, men opdagede,
at der paa dette Omraade var endnu større Lige%
% --- p. 82 ---
\setcounter{footnote}{0}%
\opage{82}gyldighed hos Folk end overfor Geometrien (fr. 144). Selv
kastede han sig nu med Lidenskab over dette Studium, medens
hans tidligere videnskabelige Interesser mere og mere
forsvandt. I et Brev til den berømte Matematiker Fermat
skriver han, henimod Slutningen af sit Liv (10. August
1660), at han stadig betragter Geometri som det højeste Omraade,
paa hvilket aandelige Kræfter kan øves, men at han
nu anser den for ganske unyttig, især da han er kommen
ind paa helt andre Studier, der er saa forskellige fra hans
tidligere Studier, at han næsten ingen Erindring mere har
om disse.

Stadig hævder han dog, at Menneskets Storhed og Adel
beror paa Tanken, og at det er den, vi maa holde os til. Vi
fylder ikke i Rummet som de store Masser omkring os, og
vor Kraft er forsvindende i Forhold til deres. Vi er som
Siv, men vi er tænkende Siv, og kan de ydre Masser end
knuse os, naar de dog ikke op til den Værdighed, Tanken
giver. (fr. 346---348). Men de Tanker, Pascal tillægger saa
stor Betydning, var dog mere og mere kun de Tanker, der
tjener til religiøs Selvforstaaelse. Ligesom han i sin „verdslige“
Periode havde beskrevet Elskovslidenskaben som en
Koncentration, en Henstyrten (précipitation) af alle Tanker
i en eneste Retning, saaledes var nu Religion bleven
hans Lidenskab, og kun de Tanker, der førte i denne Retning,
vilde han nu kendes ved.

Kierkegaard møder med helt andre intellektuelle Forudsætninger
end Pascal. Opvoxet mellem Romantikens Epigoner
betragtede han en Tid Hegel's Filosofi som den højeste
Form for Tænkning, medens for Pascal Geometrien var
Idealet. I et Tankesystem, indenfor hvilket hvert Led med
Nødvendighed førte over til et andet, skulde ethvert Emne
finde sin Plads. Poesi, Religion og Videnskab stod som forskellige
Former eller Trin af en og samme guddommelige
Sandhed. Det var hans Ungdomsdrøm. I Tilslutning til sin
Lærer, Professor Poul Møller (den mest danske af danske
Forfattere), arbejdede han sig ud af denne Drøm. Hans Be%
% --- p. 83 ---
\setcounter{footnote}{0}%
\opage{83}tydning som Filosof beror paa, at han var en af de Første,
der kritiserede spekulativ Filosofi ud fra et erkendelsesteoretisk
Synspunkt. (Uvidenskabelig Efterskrift. 1846).
Han viser, at selv om der kunde gives en „absolut“ Viden,
maatte den Enkelte dog altid begynde den Tankegang, der
skulde føre ham op paa dette høje Stade, ud fra et bestemt
Punkt indenfor Erfaringen, og hvorfor mere fra et Punkt
end fra et andet? Og kan saa dette faktiske Begyndelsespunkt
selv udledes af den „absolute“ Viden, som skal være
Resultatet? Vil der desuden ikke altid være visse Forudsætninger,
som selv den „absolute“ Viden hviler paa, og vil man
nogensinde kunne være vis paa, at man er naaet tilbage til
de sidste Forudsætninger? --- Kierkegaard kom til at se, at
Erkendelsesproblemet er det mest centrale af de filosofiske
Problemer. Han lægger særlig Vægt paa, at vi lever i Tiden.
Vi lever forlænds, men vi forstaar baglænds. En Afslutning
kan derfor ikke naas, og Kierkegaard erklærer sig enig med
Lessing i den Sætning, at ud over den stadige Stræben efter
Sandhed kommer Menneskene ikke.

Ligesom Pascal fastholder Kierkegaard, at det personlige
Liv behøver Tanker, men at disse Tanker angaar det, der
er vort Tilhold under Livets Gang. Gudstanken er en Nødhjælp,
ikke et Emne for rent intellektuel Interesse. Han
skelner skarpt mellem „væsentlig“ og „uvæsentlig“ Erkendelse,
og til den uvæsentlige regner han Naturvidenskab,
Filologi, Historie og ren Filosofi, --- alle Tanker, der er
uden Betydning for, hvorledes vi stiller os til Livets og Dødens
store Spørgsmaal. Der gives ingen rent intellektuel
Maade at tage Livet paa, naar man da ikke vil gøre sig selv
til en Abstraktion. Hvad de forskellige Maader at tage Livet
paa angaar, anvender han en lignende Fremgangsmaade
som Pascal, idet han forsøger at karakterisere visse Typer.
Vi kan her holde os til dem, han har fremstillet i sine mest
populære Skrifter („Enten—Eller“ og „Stadier paa Livets
Vej“). Den æstetiske Type gør Øjeblikkets Nydelse til det
Højeste. Kunsten er her den at bevare en vis Elasticitet, saa
at man kan gribe hvert nyt Øjeblik. Det gælder at bryde af
% --- p. 84 ---
\setcounter{footnote}{0}%
\opage{84}i rette Tid for ikke at blive holdt fast i et varigt Livsforhold.
Derfor er en vis Vexeldrift at anbefale. Æstetikeren er som
en Tangent til Livets Cirkel (eller rettere: Cirkler). Kun paa
et enkelt Punkt finder en momentan Berøring Sted; derefter
søges et andet Berøringspunkt. Derfor: flygtig Elskov, men
ikke Ægteskab; Kammeratskab, men ikke Venskab; flygtig
Beskæftigelse, men ikke Borgerskab; intellektuel Nydelse,
men ikke Videnskab! Medens Æstetikeren er en Tangent,
beskriver Etikeren en Cirkel om et Midtpunkt og bestemmes
ved Forholdet til dette. Det er Trofasthedens Type.
Etisk leves der i faste Samlivsforhold og i stadigt Arbejde.
Ud over begge Typer hæver sig det religiøse Liv, hvis Stræben
og hele Retning bestemmes ved et Maal, der ligger ud
over alle Øjeblikke og over al Tid, inkommensurabelt med
alle andre menneskelige Formaal, ja endog paa sit Højdepunkt,
i Kristendommen, i den skarpeste Modsætning til,
hvad der rent menneskeligt kan tænkes som Formaal eller
Opgave. Jo mere udpræget en Karakter det religiøse Liv
faar, des mindre lever Mennesket i sit naturlige Element,
men bliver som en Fisk paa Landjorden. Og tilsidst staar
Evigheden ikke (som hos Platon og Spinoza) som den store,
altid nærværende Baggrund for Tiden i Livet, men bryder
ind midt i Tiden, idet Guddommen selv træder frem som
Menneske paa et bestemt Tidspunkt og paa et bestemt Sted
i Rummet. Da bliver denne Aabenbaring, og Forholdet til
den afgørende for al Vurdering af Livet og dets Forhold.

Det er en sammenlignende Livsfilosofi, Kierkegaard her
har givet. Den kan minde om Pascal's Karakteristik af Montaigne
og Epiktet og af den kristelige Livstype i Modsætning
til dem begge. Og ligesom Pascal erklærede Sygdom
for de Kristnes naturlige Tilstand, saaledes kommer Kierkegaard
til det Resultat, at den højeste Livstype, den kristelige,
medfører Spænding og Lidelse, fordi de Modsætninger,
mellem hvilke Livet skal føres, staar i den mest skærende
Modsigelse til hinanden. Og denne Livstype er netop
derfor den højeste, fordi den stiller de største Fordringer
til den syntetiske Energi, alt personligt Liv kræver.

% --- p. 85 ---
\setcounter{footnote}{0}%
\opage{85}Det er ifølge Kierkegaard ved et Spring, man kommer
fra den ene Livstype til den anden, og det største Spring er
det, der fører en i Fortvivlelse stedt Sjæl over til den kristelige
Type, Men end ikke paa det mest afgørende Punkt faar
Springet hos Kierkegaard Ekstasens Karakter. Og en vis
Kontinuitet udelukker han ikke, da Maalestokken (den personlige
Synteses Energi) skal være fælles for hele Rækken
af Typer. I Grunden har Pascal jo den samme Maalestok,
den der hentes fra Forholdet mellem den Livsfylde, der
staar til Raadighed, og den Kraft og Koncentration, hvormed
det fyldige Indhold ordnes. Det er nu ogsaa den eneste
Maalestok, der kan anvendes overfor de personlige og historiske
Forskelligheder i Menneskenes Maade at tage Livet
paa, og den maa anvendes ogsaa af dem, der finder
Livets Højdepunkter andre Steder end de to Tænkere, vi
her staar overfor. --- En Indvending maa vi allerede her gøre
mod Pascal's og Kierkegaards Anvendelse af Maalestokken.
De lægger udelukkende Vægt paa den indre, subjektive
Modstand, der er at overvinde, altsaa paa det indre Arbejde,
hvorved der naas Enhed og Koncentration i den Enkeltes
Sind. Men der maa ogsaa spørges om dette indre Arbejdes
Betydning for større Kredse af menneskeligt Liv, for det
Samfund, i hvilket de lever, maaske for hele Menneskeheden.
Dette vilde hverken Pascal eller Kierkegaard gaa ind
paa. De hævder tværtimod, at der er et svælgende Dyb mellem
den i deres Øjne højeste Livstype og det menneskelige
Kulturarbejde.

\sectionhead{3. \textit{Kristendommens Problem.}}

Pascal optræder indenfor Katolicismen, Kierkegaard indenfor
Protestantismen. Dette maa naturligvis paa mange
Maader faa Indflydelse paa deres Udviklingsgang saavel
som paa deres definitive Standpunkt. Men det er af stor
religionshistorisk Interesse at se, at de ved Slutningen af
deres Stræben efter at trænge ind til Kærnen i deres Religion
% PRINTED AS IS (read at the image, odd spots o1)
kommer til at staa overfor et og sammme Problem,
% --- p. 86 ---
\setcounter{footnote}{0}%
\opage{86}et Problem, der kan formuleres saaledes: Kan den oprindelige
Kristendoms Livsopfattelse og Livsførelse fastholdes
sammen med Tilegnelse af og Arbejde for den aandelige
og materielle Kultur, som Kristendommen ikke har frembragt,
men som --- før og efter Kristendommens Opstaaen
--- har udviklet sig efter sine egne Love? Det er jo ogsaa
et Problem, der er fælles for Katolicismen og Protestantismen.
Lad os, før vi gaar videre, i Korthed se, hvorledes det
er opstaaet.

De første kristelige Slægtled levede i Forventning om
Kristi nære Genkomst, der skulde bringe Afslutning af hele
Verdensløbet gennem en Verdensdom. I denne Forventning
bestod den oprindelige kristne Tro, og baade Livsførelse
og Livsopfattelse bestemtes helt igennem ved den. Overfor
Udsigten til den nære Katastrofe, der betød Gennembruddet
af en ny Himmel og en ny Jord, maatte alle rent menneskelige
og sociale Opgaver tabe deres Betydning. Derfor er
der i det nye Testamente ikke Tale om Kunst, Videnskab
og Statsliv, og selv Familielivet anerkendes kun med et vist
Forbehold. Det nye Testamentes Etik maatte blive, hvad
Albert Schweitzer træffende har kaldt en „Interimsetik“.
Det gjaldt om at samle al Opmærksomhed mod et eneste
Punkt i en nær Fremtid. Guds Riges Komme, hvorom der
bedes i „Fader vor“, betød for de første Slægtled noget aldeles
Bestemt, ikke blot en vag Bøn om aandelige Goder.
Og hos de Troende var Forventningen forbunden med en
Begejstring, der ofte steg til Ekstase, og for hvilken ethvert
Offer blev let. Alt i det sædvanlige Menneskeliv blev smaat
overfor det Store, der ventedes.

I Modsætning til et Standpunkt som det urkristelige fremtræder
baade Katolicismen og Protestantismen som Organisationer
paa langt Sigt, byggede paa Udsigten til en lang
Udviklingsgang for det religiøse Liv som for Menneskelivet
idethele. Man blev mindre og mindre belavet paa hurtigt
Opbrud. Selv om man søgte at fastholde den oprindelige
Forventning, henlagdes dens Opfyldelse dog til en stedse
fjernere Fremtid, og Forventningen mistede efterhaanden
% --- p. 87 ---
\setcounter{footnote}{0}%
\opage{87}sin praktiske altbestemmende Magt. Og baade den fra den
klassiske Oldtid overleverede Kultur og den nye Kultur,
der udviklede sig uafhængig af Kristendommen, søgte man
(efter tilbørlig Modstand, ofte af bitter Art) efterhaanden
at tilegne sig. Men trods Alt hævder begge Kirker at være
i Kontinuitet med Urkristendommen og søger, hver paa
sin Maade, at begrunde Paastanden om at fastholde de
gamle Idealer i de nye Tider.

Den mest storartede Løsning er Katolicismens. Den beror
paa, at der skelnes mellem Grader i Fuldkommenhed.
Munken og Nonnen opgiver Deltagelsen i Verdenslivet for
helt og udelt at leve for det Hinsidige og i det ene Fornødne.
Gennem dem vedligeholdes Sammenhængen med
Kristendommens Urtid. Lægfolket fører saa det menneskelige
Liv videre under Ledelse af Skriftefædrene, gennem
hvilke Kirken øver en kontrollerende Indflydelse paa de
Enkelte. I Kirkens Prælater saa allerede Augustinus (De
civitate dei XX, 9) en delvis Opfyldelse af Apokalypsens
profetiske Skildring af dem, der sidder paa Dommersædet
i Tusindaarsriget.

Pascal var klar over, at der til enhver Tid i Kirken maatte
raade en praktisk Anerkendelse af den urkristelige Fjernhed
overfor alt Kulturliv. Han levede i sin sidste Tid som
Munk, selv om han ikke gik i Kloster. Tiden indtil Jesu
Genkomst maatte være en Spændingens og Lidelsens Tid.
Scenen i Getsemane var i hans Øjne Udtryk for sande Kristnes
stadige Situation. Der maa leves i Samtidighed med
den lidende Frelser. Den eneste Gang, Jesus i sin Kvide
søgte Hjælp og Trøst hos Mennesker, var i Getsemane. Da
svigtede Apostlene ham. Men de Kristne har stadig den
Pligt, som da blev svigtet. Jesu Dødsangst varer ved til Verdens
Ende, og de Troende maa holde sig vaagne indtil da.
(Jésus sera en agonie jusqu'à la fin du monde; il ne faut
pas dormir pendant ce temps-là. fr. 583). Det er altsaa ikke,
som hyppigst i Kirken, den sejrende, men den lidende Jesus,
Blikket skal rettes imod. Getsemanes Have er Mod%
% --- p. 88 ---
\setcounter{footnote}{0}%
\opage{88}stykket til Edens Have. Det faldne Menneske har at opholde
sig i Getsemane og kan kun derved have Udsigt til,
at Eden skal oplukkes for ham.

I jesuitiske Skriftefædres Praxis saa Pascal nu en forbryderisk
Tillempelse af kristeligt Liv efter Verdenslivet,
og han rettede en skarp Protest derimod i sine „Lettres à
un provincial“. Trods Fejltagelser i det Enkelte har han
Ret fra et strengt kristeligt Standpunkt. Det er en forskellig
Etik, der forkyndes i det nye Testamente og i den jesuitiske
Kasuistik, og selv fra et rent menneskeligt Synspunkt
kunde der være Betænkeligheder ved denne Kasuistiks Resultater.
Men foruden den Kamp, Pascal i Provincialbrevene
førte mod den milde Religion, der efterhaanden havde afløst
Getsemanestemningen, har han i en særlig Afhandling
anstillet en Sammenligning mellem de første Tiders Kristendom
og vor Tids. Naar Afstanden her er saa stor, finder
han Aarsagen dertil i, at nu gaar Daaben forud for Undervisningen
i, hvad Kristendom er, medens oprindelig Ingen
gik over til Kristendommen uden at være grundig underrettet
om dens Lære og uden at have givet Afkald paa
Verden. Dengang kunde man kende Forskel paa Verden
og Kirken, men nu har Kirken gennem Barnedaaben gjort
det muligt, at verdsligt Liv kan brede sig indenfor den. Det
var i en god Mening, føjer Pascal til, at Kirken indførte
Barnedaaben, nemlig af Frygt for at udelukke de Børn,
der døde før Daaben, fra Frelseren. Men hvad der her er
gjort af Omsorg for Børnene, er blevet til Fordærv for de
Voxne! ---

Det var, som vi allerede har set ovenfor, ikke blot paa
dette Punkt, at Pascal's Brud med Kirken skete. Han førtes
til at forkaste Kirkens Autoritet, baade hvad Læresætninger
og hvad Fastsættelse af Kendsgerninger angaar. Sammen
med sine jansenistiske Venner hævdede han, at de Sætninger,
Paven havde fundet i Jansenius' Værk „Augustinus“
og fordømt som kætterske, slet ikke fandtes i den nævnte
Bog, og han negtede Paven Ret til at fastslaa noget som
Kendsgerning, enhver selv kunde og maatte overbevise sig
% --- p. 89 ---
\setcounter{footnote}{0}%
\opage{89}om. Allerede tidligere havde han (i en Afhandling om det
tomme Rum) bestemt afvist Autoritetstroens Indblanding
i Erfaringsvidenskabens Problemer. Og i det 18. Provincialbrev
udtaler han, at hvis Galilei's Lære om Jordens Bevægelse
er falsk, saa er det ikke Pavens Dekret, der begrunder
det, og at hvis der virkelig kan føres Bevis for, at Jorden
bevæger sig, saa gaar den rundt baade med den, der
negter Bevægelsen, og med den, der antager den. Pascal
var ikke Kopernikaner, men han rørte her ved et for en
Katolik paa den Tid meget farligt Punkt, det, der havde
bragt Bruno paa Baalet og tvunget Galilei til Mened. En
mod kirkelig Autoritet oprørsk Aand danner en underjordisk
Forbindelse mellem Pascal som Naturforsker og som
religiøs Tænker. Han udtaler endog, at Paven hader og frygter
de Forskere, der ikke har aflagt Ed til ham (som Munke
eller Præster) (fr. 873). Denne Udtalelse fremkommer i Anledning
af, at Paven fordømte Læren om den guddommelige
Naades Almagt og om den menneskelige Villies fuldkomne
Afhængighed af Guds Villie, denne augustinske
Lære, som Jansenius havde fornyet i skarp Modsætning til
skolastisk og jesuitisk Teologi. Saa er Paven selv altsaa
ikke ortodox, og han kan ligesaalidt i Trossager som angaaende
Kendsgerninger, være ufejlbar Autoritet. Men Paven
har jo ogsaa saa meget Andet at tage vare, at han ikke
har Tid at sætte sig ind i Trossager (fr. 882)! Og naar han
fordømmer Sandheden, saa maa man raabe saa meget des
højere, og der maa appelleres fra Rom til Himmelen! Vi
skal adlyde Gud mere end Menneskene (fr. 920).

Denne sidste Udtalelse fandt Pascal's jansenistiske Venner
det bedst at undertrykke, da de udgav hans „Pensées“.
Det vilde have været en dyb Sorg for Pascal at mærke, at de
her ikke gav Sandheden Æren. Kun Døden havde hindret
ham i selv at udgive sine sidste Tanker. --- Paa sit Dødsleje
modtog han Sakramentet af sin Sognepræsts Haand.
Deri saa hans Modstandere en Tilbagekaldelse; men det
betød kun, at han ikke ansaa sit Forhold til den sande
Kirke for brudt. ---

% --- p. 90 ---
\setcounter{footnote}{0}%
\opage{90}Kierkegaard havde udviklet sig indenfor den protestantiske
Kirke, fra først af med stor Ærbødighed for dens Førere,
de nulevende som de tidligere. Men i Aarene 1847—
50 skærpedes hans Blik for Modsætningen mellem oprindelig
og moderne Kristendom, og i to af hans dybest gaaende
Skrifter („Sygdommen til Døden“ 1849 og „Indøvelse i Kristendom“
1850) forberedes et stort Opgør med Kirken. Hovedtanken
i de nævnte Skrifter var, at man, især i Protestantismen,
havde glemt Forbilledet for Forsoneren. Man
tog Forsoningen som given Kendsgerning og indrettede sig
hyggelig i Verden, som den er. Kierkegaard henstiller nu
til Kirkens Repræsentanter at gøre den Indrømmelse, at
der var en stor Afstand mellem moderne Kristendom og
det nye Testamentes Kristendom, og at der var sket et Affald
i Forhold til Kristendommens Heltetid. Men da man
fra den officielle Kirkes Side tav hertil og endog ved given
Lejlighed udtrykkelig proklamerede Paastanden om at
være i Kontinuitet med „Apostlenes Dage“, udstedte han en
voldsom Protest herimod, især i sit Agitationsskrift „Øjeblikket“
(1855). Midt under Striden faldt han sammen og
døde. Han døde uden Nadverens Sakramente, da han ikke
vilde tage det af en Præsts Haand. Han er her et Skridt
længere ude end Pascal; men han har ikke forkastet Sakramentet
overhovedet, og her hænger han for sit Vedkommende
endnu sammen med den gamle Kirke.

„Sygdommen til Døden“ bestaar i den Fortvivlelse, som
uundgaaelig fremkaldes ved, at der i Kristendommen er sat
Mennesket den høje Opgave at leve et Aandens Liv i den
skarpeste Modsætning til det Liv, naturlige Følelser og Drifter
vilde føre til. Hvor denne Opgave ikke er kommen til
Bevidsthed, føres der maaske et Liv i barnlig Glæde over
Livet. Men Angsten vil atter og atter melde sig og gøre
Ende paa Trygheden. Mennesket kan da i Svaghed, i Indesluttethed
eller i Trods forsøge paa at dæmpe Angsten og
unddrage sig fra at lytte til den store Fordring; men Fortvivlelsen
vil stedse være skjult tilstede. Kun ved at Fortvivlelsen
bryder aabent frem og udfolder sig, er der Mu%
% --- p. 91 ---
\setcounter{footnote}{0}%
\opage{91}lighed for at komme ud over den. Da først er Betingelsen
tilstede for, at Naaden kan træde i Kraft. Men indenfor Kristenheden
har man, snart i Tankeløshed, snart i romantisk
Sværmeri, snart i Hykleri, slaaet af paa den store Modsætning
mellem det Aandelige og det Dyriske, søgt at udjævne
den istedetfor at arbejde paa at overvinde den.

„Indøvelse i Kristendom“ indskærper, at saalænge Jesus
ikke er kommen i sin Herlighed, har den Kristne at tage
den fornedrede Jesus til Forbillede. Det gælder om at se
bort fra, hvad man gennem Dogmatiken mener at have
faaet at vide om, at Jesus var Gud. Man skal gøre sig samtidig
med Jesus fra Nazareth, det Menneske, der sagde sig
at være Guds Søn og derfor blev bespottet og henrettet.
Istedetfor har man som Skoledrenge lusket sig til et Facit
og snydt sig fra at foretage Regningen paa egen Haand. ---
Parallelen med den af Pascal fordrede Getsemanestemning
ligger lige for Haanden.

Det blev mere og mere Kierkegaards Overbevisning, at
Menneskene havde begaaet et stort Forræderi overfor Kristendommen.
Man tog Mildheden og Naaden og glemte de
nødvendige Forudsætninger. Og han hævder, at især i Protestantismen
ligger dette Forræderi tydelig for Dagen. Katolicismen
vilde aldrig kunne synke saa dybt i Aandløshed
som Protestantismen. Og dog er Katolicismen allerede gaaet
paa Akkord. Meningen i Kristendommen brast, da det betragtedes
som noget Overordentligt at følge Jesus efter for
Alvor ved at unddrage sig Verdenslivet. Det er en Meningsløshed,
at der skulde gives to Slags Kristne. Men hvad der
saaledes var begyndt, fortsattes, da Luther brød ud af Klostret
og fik travlt med at lave en ny Kirke, der paa gemytlig
Maade udjævnede Modsætningen mellem Kristendommen
og Verden. „Luther, Luther, Luther, du har et stort Ansvar!“
Og Luther havde dog gennemgaaet Fortvivlelsens
Skole. Men senere Slægter stoler saa paa Forfædrenes Bedrifter
og har ingen Trang til selv at øve Bedrifter. Protestantismen
er en plebejisk Retning. Den nivellerer Forskellen
mellem det Høje og det Lave i Aandens Verden. Og
% --- p. 92 ---
\setcounter{footnote}{0}%
\opage{92}Familielivet har for en stor Del Skylden derfor. Gennem
Barnedaaben blev jo alle i de saakaldte kristne Lande gjort
til Kristne. „Og det“ --- udraaber Kierkegaard (Øjeblikket
VII) --- „vover man at byde Gud under Navn af kristelig
Daab! Daaben --- det var ved denne hellige Handling, at
Verdens Frelser indviedes til sit Livs Gerning, og efter ham
Disciplene, Mænd, der forlængst var komne til Skelsaar og
Alder, og som nu, døde for dette Liv, lovede at ville leve
som ofrede i denne Løgnens og Ondskabens Verden.“

Kierkegaards Polemik kulminerede i Opstillingen af én
eneste Tesis (i Modsætning til Luthers 95): „det nye Testamentes
Kristendom er ikke til.“ Og dertil føjedes en Advarsel
mod „dem, der gaar i lange Klæder“ og mod Deltagelse
i den offentlige Gudsdyrkelse, i hvilken Gud holdtes
for Nar!

Pascal's og Kierkegaards Optræden har stor religionshistorisk
Interesse. Fra et udpræget sociologisk Synspunkt
er det blevet hævdet, at Religion væsentlig er et socialt Fænomen.
Og faktisk interesserer Religionshistorien sig væsentlig
for saadanne religiøse Fænomener, der hviler paa
traditionel Kultus. Hvorledes kan da enkelte Personligheder,
hvis Stræben og Tænken fører dem til et Brud med
deres religiøse Samfund (uden at de gaar ind i et andet
Samfund) faa religionshistorisk Betydning? Man kunde besvare
Spørgsmaalet ved at henføre dem hver til sin „Sekt“,
for Pascal's Vedkommende til Jansenismen, for Kierkegaards
Vedkommende maaske (saaledes som det er sket i
et udmærket fransk sociologisk Værk) til Pietismen. Dette
vilde dog strande paa, at Pascal netop i sin sidste, afgørende
Vending under stor Smerte skilte sig fra sine tidligere
Meningsfæller, og at Kierkegaard kun har Satire og Haan
overfor Pietisterne og udtrykkelig skiller sin Sag fra deres.
Snarere kunde man sige, at de kun bryder med den i deres
Tidsalder raadende kirkelige Overlevering for at hævde en
ældre, i deres Øjne ene berettiget Overlevering. Men med
deres Appel til det i deres Øjne Oprindelige stod de ene i
% --- p. 93 ---
\setcounter{footnote}{0}%
\opage{93}deres Tid, og det var i den individuelle Samvittigheds Navn,
at de protesterede. Der var ingen social Instans, til hvilken
der kunde appelleres. Pascal fandt Bekræftelse for sin Lære
i et Mirakel, men dette Mirakel var ikke kirkelig sanktioneret.
Og Kierkegaard erklærede udtrykkelig, at den religiøse
Type, der for ham var den eneste gyldige, ikke mere
existerede. Kun i deres Forhold til Sakramentet paa Dødslejet
er der et Baand tilbage, som dog ikke viste sig tilstrækkelig
stærkt til at bringe Forbindelsen med Kirken tilveje
igen. --- Men det er en uholdbar Teori, der lægger saa stor
Vægt paa Overlevering og Kultus, at man glemmer, hvad
der foregaar i de enkelte Sjæle, som hylder disse sociale
Magter. Det maa blive en væsentlig Opgave for Religionsforskningen
at efterspore de religiøse Personligheders Udvikling
og Holdning. Og da ikke mindst saadanne individuelle
Udviklingsprocesser, som fører ud til vedkommende
Religions Grænse netop ved Forsøget paa at gennemleve
denne Religion i Aand og Sandhed. Og saadanne Forsøg
foreligger netop i Pascal's og i Kierkegaard's religiøse Udvikling.

\sectionhead{4. \textit{Mulige Veje.}}

Hvad de to Mænd vilde have gjort, hvis de var komne
til at fortsætte deres Liv, derom kan Intet vides. Men det
kan have sin Interesse at undersøge, hvilke Muligheder
der kan fremstille sig for saadanne Mennesker, som med
fuld Overbevisning og Tilslutning og støttede paa deres
egne Livserfaringer har fulgt de to aandelige Stridsmænd
lige til det Punkt, da de havde tænkt deres sidste Tanke.

Der er tre Muligheder, der her fremtræder.

Man kunde prøve paa at føre sit Liv med Genoptagelse
af Urkristendommens Forudsætninger, i den store Forventning
om det nære Gennembrud og med Opgivelse af
al Deltagelse i Kulturarbejdet og i det sædvanlige Menneskelivs
Interesser. Der er maaske --- hvem véd? --- dem,
der gaar denne Vej. Det var den Vej, som i Middelalderen
Franciscus af Assisi paa sin Vis fulgte.

% --- p. 94 ---
\setcounter{footnote}{0}%
\opage{94}Eller man kunde gøre den af Kierkegaard forlangte
Indrømmelse om Afstanden mellem oprindelig og moderne
Kristendom. Men man maatte da gaa et Skridt videre og
anerkende, at moderne Kristendom egentlig er en anden
Religion end den oprindelige Kristendom. Da undgaar man
den store Modsigelse at mene sig at leve efter Idealer,
man faktisk fornegter. Man maa da opfatte sin Religion
saaledes, at saa væsentlig en Faktor den ogsaa er for Livet,
maa Livsførelsen og Livsopfattelsen tildels ledes af Idealer,
der har en helt anden Oprindelse og fører i en helt anden
Retning end den Religion, man hylder. Det kan være vanskeligt
at faa disse forskellige Ledestjerner til at høre til
samme Verdenssystem. Men den Redelighed, de to store
religiøse Tænkere krævede, er ikke paa Forhaand krænket.

Kan man ikke gaa nogen af disse Veje, er der endnu den
Mulighed tilbage, at man forsøger at genoptage den græske
Stræben efter at grunde en etisk Livsførelse paa menneskelig
Natur og menneskelig Livserfaring, idet man undersøger,
hvilke Idealer og hvilke Formaal der konsekvent
maa stille sig paa dette Grundlag. Det er faktisk den Vej,
Bestræbelserne for en filosofisk Etik i de sidste Aarhundreder
er gaaet. Jeg har engang (i min Etik\textsuperscript{5} X, 3) udtrykt dette
saaledes: „Det platoniske Begreb om Retfærdighed som personlig
Harmoni kan vi optage til at udtrykke den indre
Forbindelse, den højere Enhed af Selvhævdelse og Hengivelse.
Det antike Harmonibegreb maa udvides og uddybes
ved Hjælp af de etiske Erfaringer, der ligger udtrykte i
Stoicismens, Kristendommens, Renaissancens og den moderne
Humanitets Karaktervurderinger. Men en saadan
Ændring af Begrebets Indhold er meget vel mulig uden at
sprænge dets Ramme. Vor Etik er den græske Etik, som
den senere etiske Udvikling gennem sine mangehaande
Svingninger har kunnet ændre og berigtige, men som aldrig
vil kunne fortrænges, saalænge der skal bestaa Noget, der
med Rette kan kaldes menneskelig Etik.“ --- Paa dette
Grundlag kan der ikke arbejdes, uden at der tillige udvikler
sig en Tro paa, at et saadant Arbejde stedse er muligt, ---
at Tilværelsen altsaa har Plads for det, ja, maaske først
% CHECK p. 94: superscript, no note: an edition number?
% --- p. 95 ---
\setcounter{footnote}{0}%
\opage{95}kommer til fuld Udvikling gennem det. Den nærmere Udformning
af en saadan Tro kan kun ske gennem poetisk
Symbolik. En stor og dyb Poesi hører derfor selv til de
Værdier, paa hvis Bestaaen og Udfoldelse der maa haabes.

Overfor de to andre antydede Standpunkter bliver Forholdet
det, at man --- istedetfor at forsage al naturlig Udvikling
af Menneskelivet, al selvstændig Kultur, alt Liv i
Familie og Samfund, i Kunst og i Videnskab, og istedetfor
at foretage en udvortes Blanding af østerlandsk Religion og
vesterlandsk Kultur --- stiller sig bestemt paa Menneskenaturens
og den vesterlandske Kulturs Grundlag og trøstig
arbejder videre, saa langt man kan naa. Her kan være Problemer
nok, og der kan komme store Skuffelser, saa at
Mange taber Modet. Der kan være Strid om Vejene i det
Enkelte, og om Begrundelsen af de Tanker, der leder os;
men Hovedlinierne vil bestandig paany kunne skelnes. ---
Og der kan meget vel drages Næring og høstes Belæring
fra andre Livsretninger. Den Pathos, hvormed den Troende
venter det altid nære Gudsrige, kan faa et mere eller mindre
udpræget Sidestykke i enhver Stræben efter høje Maal.
Uden Spænding gennemføres intet stort Arbejde. Sammenligner
vi Kristendommens Betydning for europæisk Menneskeliv
med Buddhismens Betydning for asiatisk Menneskeliv,
kan den Sætning forsvares, som jeg har fremsat i
min Religionsfilosofi, at Buddha har gjort Asien mild, men
Jesus har lært Europa et stort Excelsior. Men fremfor Alt
vil den Menneskekærlighedens Aand, som Kristendommen
har forkyndt og (ikke altid med fuld Konsekvens) stillet i
Spidsen for Dyderne, ogsaa finde sin Plads i en rent menneskelig
Etik. Allerede Aristoteles har udtalt, at Retfærdighedens
højeste Form har Kærlighedens Præg, og omvendt
gælder, at Kærligheden kun er en virkelig Dyd, naar den
forener sig med Visdom og derved bliver til Retfærdighed.
Paa det humane Standpunkt staar man overhovedet i et
lignende Forhold til Kristendommen som til andre Aandsretninger;
man tilegner sig deraf, hvad der stemmer med
stadig fortsat Stræben. De store Verdensreligioner har nu
engang, ligesom Kulturen, deres Udspring i menneskelig
% --- p. 96 ---
\setcounter{footnote}{0}%
\opage{96}Natur og menneskelige Livsforhold, og hvad der af Ædelt,
Stort og Skønt har udviklet sig indenfor Religionerne, tilhører
derfor tilsidst Menneskeslægten som Helhed, ikke en
enkelt Sekt, selv om denne har nok saa stor en Udbredelse.
---

Hvilken af de tre Veje man saa gaar, er der én Ting,
man kan lære af de to mægtige Skikkelser, vi her har dvælet
ved. De har begge forlagt Tyngdepunktet i Aandslivet
fra de ydre Autoriteter og fra Læreindholdet til Personligheden,
dens Erfaringer og dens Krav. Pascal motiverer Paastanden
om, at Kristendommen er den sande Religion, derved,
at først i den fyldestgøres Menneskenaturens dybeste
Trang, idet den paa én Gang forklarer og forliger Livets
største Modsætninger. Og for Kierkegaard var Gudstanken
et Nødanker under Livets Storme, som den, der stedse færdes
i smult Vande, ikke forstaar Nødvendigheden af. Det
kommer da, efter begges Tankegang, an paa, at der gøres
de rette, dybtgaaende Erfaringer under Livets Gang. Men
først og fremmest gælder det om ikke at bedrage sig selv
og Andre ved at mene sig at hylde en Tro, der ikke er fremgaaet
af egen Livserfaring og egen Trang. „Pour les religions“
--- siger Pascal (fr. 590) --- „il faut être sincère; vrais
paiens, vrais juifs, vrais chrétiens.“ Og i sin sidste Kamp
mod den bestaaende Kirke siger Kierkegaard: „Jeg tør ikke
kalde mig en Kristen; men Redelighed vil jeg, og for den
vil jeg vove.“

Men der er et Punkt, hvor de, der gaar den tredie Vej,
maa være afgjort uenige med Pascal og Kierkegaard. De er
begge klare over, at Livet, ogsaa det religiøse Liv, ikke kan
undvære Tanker, som fastholder og leder Livets Retning.
Men de skelner saa absolut mellem Livstanker og videnskabelige
Tanker, eller, som Kierkegaard udtrykker det,
mellem „væsentlig“ og „uvæsentlig“ Erkendelse. Til den
„uvæsentlige Erkendelse“ regnes al Videnskab. Og dog er
der videnskabelige Tanker, som kan blive af den største Betydning
for personlig Tro. Grænsen kan ikke én Gang for
alle afstikkes. Stille og umærkelig for de Fleste øver saaledes
% --- p. 97 ---
\setcounter{footnote}{0}%
\opage{97}den storartede Udvikling af Videnskaben, baade Naturvidenskaben
og Aandsvidenskaben, i de sidste Aarhundreder
sin Indflydelse paa de religiøse Forestillinger. Det er idethele
en falsk Opfattelse, at Personlighed og Videnskab
skulde staa i Modsætning til hinanden som hvilende paa
helt forskellige Interesser. Videnskab skyldes personligt Arbejde
og udspringer af personlig Trang, og Erkendelsesglæden
hører til Menneskets ædleste Følelser. Selv de Former,
i hvilke det videnskabelige Tankeliv bevæger sig, uden altid
at være sig det bevidst, fører ved nærmere Undersøgelse
tilbage til Personligheden, idet de er betingede ved den menneskelige
Aands Beskaffenhed. Paa den anden Side udvikles
Personligheden gennem det videnskabelige Arbejde, vinder
i Klarhed over Livet og dets Kaar og lærer at stille sig Opgaver,
der fører langt ud over de egoistiske Interesser. Og
som det gaar med Videnskaben, gaar det med al Kultur.
Mennesket er baade større og mindre end den til hver Tid
givne Kultur: større, fordi al Kultur tilsidst maa godtgøre
sin Værdi ved at bidrage til det personlige Livs Udvikling,
--- mindre, fordi Kulturen stiller Opgaver, der kræver Underordning
af individuelle Interesser under det, der tjener
til varig Fremgang for Menneskeheden idethele.

Det er to ensomme Skikkelser, vi har dvælet ved. Efterfølgere
har de ikke havt, og Skole har de ikke dannet. Men
Analogierne mellem de Tanker, der leder dem hver for sig
og mellem deres Optræden overfor det Bestaaende, de hver
for sig stod overfor, indeholder Belæring for alle Tider.
De har fra deres ensomme Stader faaet Øje paa store, afgørende
Problemer paa Religionshistoriens og Religionspsykologiens
Omraade, og de har sat deres Liv ind paa at
slaa disse Problemer saaledes fast, at de ikke kan glemmes.
Og Glemsel vil saa meget des mindre blive de to Mænds
Lod, som de har fremsat deres Problemer, ikke blot med
Lidenskab, men med en Tankerigdom og Stemningsfylde,
der i og for sig lader deres Værker høre til Menneskehedens
kosteligste Skatte.

% --- p. 98 ---
\setcounter{footnote}{0}%

\opage{98}\essayhead{SIBBERN OG KIERKEGAARD}

Begge disse Tænkere har haft stor Indflydelse paa min
Udvikling, og selv efter at jeg mente at være kommen til
Skelsaar og Alder i filosofisk Henseende, er jeg ofte vendt
tilbage til dem. Og da jeg nu i Vinter [1925] af nogle unge
Studenter fik Opfordring til at holde nogle Forelæsninger
over danske Filosofer, blev det især disse to Mænd, jeg
dvælede ved. Det er de Fremstillinger, jeg ved den Lejlighed
gav, som jeg nu udarbejder til en lille Afhandling.

Kierkegaard fødtes samme Aar, som Sibbern blev Professor
i Filosofi ved Københavns Universitet (1813). Aldersforskellen
mellem dem var altsaa stor, og Karakterforskellen
ikke mindre. Og dog har de staaet i personligt Forhold
til hinanden. Kierkegaard holdt sig en Tid til Poul Møller,
hvem han elskede og beundrede, men efter hans tidlige Død
har han en Tid staaet Sibbern ret nær. Det vides ikke blot
af et Brev, han fra Berlin skrev til Sibbern (1841). Men en
af Sibberns Døtre har i et Brev til mig fortalt, at det hører
til hendes Barndomserindringer, at hun gik med sin Haand i
sin Faders, medens Kierkegaard gik paa den anden Side, ivrig
talende og af og til standsende for ret at faa sagt, hvad
der laa ham paa Hjerte. Hun erindrer ogsaa, at Kierkegaard
ofte besøgte hendes Fader og holdt af at sidde i Dagligstuen
i Mørkningen, naar Ilden lyste fra Ovnen. Om Indholdet af
deres Samtaler berettes intet, men Forskellen mellem de to
Mænds Karakterer og Synsmaader maa atter og atter have
gjort sig gældende. Sibbern kendte af egen Erfaring Tider
af indre Kamp og Gæring; Skildringen af saadanne Tider
% --- p. 99 ---
\setcounter{footnote}{0}%
\opage{99}hører til det bedste i hans Psykologi. Men han bevarede
trods alt et aabent, naivt Blik paa Livet. Som Psykolog
havde han Forstaaelse og Sympati for alle sjælelige Fænomener,
og det var en Yndlingssætning hos ham, at ligesom
der for Botanikeren intet Ukrudt gives, saaledes var der
intet sjæleligt, som Psykologen ikke kunde have Interesse
af at studere. Ikke blot formedelst denne intellektuelle Interesse,
men ogsaa ved sin faste Tro paa Livet fik han bestandig
Øje paa en Harmoni bag alle Modsætninger. Kierkegaards
Tungsind og aldrig hvilende Reflexion lod ham vandre
langs Afgrunde, fra hvilke kun krampagtig Selvhævdelse
reddede ham. Hans skarpe Blik for Livets Modsætninger
førte ham til Tanken om et absolut Paradox, og Evnen
til at lide blev for ham Villiens højeste Form. Men
trods den store Modsætning mellem Sibberns aabne Blik
for alt menneskeligt og Kierkegaards lidenskabelige Begrænsning
til den snevre Vej, der for ham førte til det ene
fornødne, mødtes de dog som Forkyndere af Personlighedsprincipet,
af den Tanke, at den Sandhed, som har
Værdi, maa udspringe af Menneskets egen Erfaring og Eftertanke.
Dette Princip udtaltes mærkeligt nok et og samme
Aar (1846) af den ældre og den yngre Tænker. De vidste
neppe, at Personlighedsprincipet med stor Klarhed allerede
nogle Aar forud (1842) var forkyndt af den svenske Tænker
Erik Gustaf Geijer. Alle tre nordiske Tænkere brød ved
dette Princip med Romantiken, med dens Liv i Fantasiens
Verden saavel som med dens Afhængighed af det historisk
overleverede. Det er af stor Betydning at mærke sig, at det
ikke blot var streng naturvidenskabelig og historisk Forsken,
som førte bort fra Romantiken og den med den forbundne Konservatisme, men ogsaa den individuelle Personligheds
Krav paa fri, selvstændig Udvikling.

Sibbern (1785---1872) dukkede i sin Ungdom ned i den
romantiske Strømning, der føltes som Vederkvægelse efter
Rationalismens Tørhed. Men han dukkede op igen, og hans
hele Liv var en Stræben efter fri Forsken over Livet med
% --- p. 100 ---
\setcounter{footnote}{0}%
\opage{100}Bevarelse af den Sans for Fortidens Storhed, der havde udmærket
Romantiken.

Sibbern var først og fremmest beskrivende Psykolog. Jeg
vil her kun dvæle ved de Emner, han bedst beskriver. Dertil
hører, som allerede nævnt, Gæringsperioder. Han var
klar over, at jo dybere Livet aabner sig i et Gemyt, desto
snarere kan det komme til at gennemgaa en kaotisk Elementernes
Blanding og Røring i sit Indre. Og saadan Gæring
fører let til en sygelig Optagethed af sig selv, til Grublen
og Angst. Han mente, at saadan Hildethed og Nervøsitet
særlig var karakteristisk for „vor Tid, i hvilken dels den
større Kamp for fysisk og borgerlig Subsistens har virket
til Selvreflexion, dels Individerne ikke mere, saaledes som
forhen, føler sig, hverken i religiøs eller i borgerlig Henseende,
baarne ved en fast Tingenes Orden i Livet“. Denne
Udtalelse er fra 1828; der er vist mange, som mener, at den
passer den Dag idag. I Modsætning til Gæring og Selvreflexion
stiller saa Sibbern Sympati og Beundring, Glæden
ved Livet, som det udfolder sig omkring os og Optagethed
af det Store og det Ophøjede, som fører den enkelte
ud over sig selv.

Psykologien var for Sibbern en Del af en almindelig
Livslære, til hvilken ogsaa Fysiologien giver sine Bidrag.
Psykologi og Fysiologi hørte for ham nøje sammen. En og
samme Grundvirken træder frem for os under to Former.
Ordet Sjæl betyder ikke et for sig bestaaende Væsen, men
en Leven og Virken, der stadig har et tilsvarende indenfor
de organiske Processer. Sibbern var hverken Spiritualist
eller Materialist. Den af Spinoza og Leibniz hævdede Enhedsopfattelse,
som han slutter sig til, var traadt ham i
Møde hos hans Forgænger Treschow.

For Sibbern var i Følge denne Opfattelse Sjælelivet underlagt
de samme store Naturlove som alt andet, navnlig
Udviklingens Lov. Tidlig blev det for ham et Hovedpunkt,
at al Udvikling foregaar sporadisk, ud fra forskellige Udgangspunkter,
og at Helheden paa alle Omraader bliver
til ved Sammensmeltning af Processer, der en Tid kunde
% --- p. 101 ---
\setcounter{footnote}{0}%
\opage{101}synes ikke at have noget med hverandre at gøre. Han belyser
det ved Exempler baade fra den uorganiske og den
organiske Verden, og paa Sjælelivets Omraade finder han
her en Forklaring paa, hvorfor Gæring, Tvivl og Disharmoni
saa let optræder; de skyldes den Maade, hvorpaa
Trang og Evne saa ofte først optræder som forskellige,
maaske modsatte, saa at det bliver Opgaven at harmonisere
dem. Paa det sociale Omraade optræder det sporadiske
i Menneskers forskellige Trang og Opfattelse og i Modsætninger
mellem Folkene.

Sibberns Verdensanskuelse er betinget ved denne Opfattelse.
Der foregaar en stor Udvikling fra Kaos til Harmoni.
Verden, og særlig Menneskene, er maaske ikke kommet ret
langt frem ad denne Vej. Hvad Menneskene angaar, er de
maaske kun i Fostertilstanden, eller i det Højeste i Lømmelalderen.
Og trøstig ser Sibbern da hen til Udviklingens
fortsatte Løb. Han var, som hans nære Ven A. S. Ørsted
sagde om ham, en uforbederlig Optimist. ---

I sin Opfattelse af Filosofiens Stilling og Betydning
stod Sibbern i Modsætning til den Hegelske Filosofi, der i
Danmark væsentlig hyldedes af Æstetikere som Heiberg
og Teologer som Martensen. Et en Gang for alle afsluttet
System ansaa han for umuligt. Selv om det skulde lykkes
at sammenfatte Tidens Kundskab til et hele, saa vilde, medens
dette systematiske Arbejde stod paa, Erfaringen være
skreden frem, og der vilde faktisk være kommen en ny Tid
med nyt Indhold. Desuden finder Sibbern ogsaa her Anvendelse
for sin Lære om sporadisk Udvikling. Enhver Erkendelse
udfoldes ud fra et enkelt Punkt i Tilværelsen. Og
saaledes staar ogsaa enhver Filosofi paa et bestemt Stade
og kan kun fremstille, hvorledes Tilværelsen tager sig ud
fra dette enkelte Punkt. Vi forstaar da, hvorfor der er Stridigheder
mellem Filosofer. Det er et Tegn paa, at vor Erkendelse
er i sin Vorden.

Ogsaa hvad angaar Filosofiens Forhold til andre aandelige
Omraader har Sibbern en anden Opfattelse end Hegel.
Kunst og Religion er ikke ufuldkomne Former for Filosofi,
% --- p. 102 ---
\setcounter{footnote}{0}%
\opage{102}men selvstændige Former for aandeligt Liv, hvis Betydning
det er Filosofens Opgave at klargøre.

I sin Ungdom var Sibbern stærkt betagen af positiv Kristendom.
Han priste den for den Aandsfylde, den aabnede
Adgang til, og for den Hvile og Vederkvægelse, den kunde
bringe. Han fastholdt dog Filosofiens Selvstændighed overfor
Kristendommen som overfor alle andre Former for
Aandsliv. Efterhaanden førtes han til at skelne skarpt mellem
de udformede Dogmer og de Livserfaringer, paa hvilke
de oprindelig var byggede. Og hans Idé om det sporadiske
i al aandelig Udvikling maatte ogsaa føre ham til den Overbevisning,
at en og samme Tro ikke svarer til alle Trin
eller til alle individuelle Livserfaringer og Livsformer.
Hans egen Livserfaring førte ham bort fra Kristendommen.
Alt, hvad der har aandelig Værdi, begynder indefra, som
subjektivt. Det er Aandens Vej. Senere kan saa Formuleringer
og Organiseringer finde Sted og en tilsyneladende
Objektivitet naas. Men naar man dogmatiserer, rykker
man let Sandheden bort fra dens levende Rod og gør en
Mumie af den. Og den Omstændighed, at store Samfund
slutter sig sammen om visse Dogmer, gør disse ikke derfor
mere objektive, hvis man ved „objektiv“ forstaar det
samme som ved „gyldig“.

Disse Tanker fremsatte Sibbern i et Universitetsprogram
1846, et af de vigtigste Dokumenter for det religiøse Problems
Historie i Danmark. I sin Utopi (forfattet i Aarene
1858---1872), som jeg strax skal vende tilbage til, skildrer
han Fremtidens Mennesker, som de efter hans Haab vil
være til Sinds om tre Hundred Aar. De vil ikke mere kalde
sig Kristne, fordi Kirkens Historie i deres Øjne er fuld af
Anmasselse og Herskesyge. Dog finder de meget stort og
skønt i Kristendommen og skelner derfor mellem Kristelighed
og Kristenskab. Men de læser i Naturens og Livets
store Bog, ikke blot i Bibelen. De lever ud af en Kraft, der
rører sig i deres Indre, og de har den Tro, at deres egen
Fremadstræben og den guddommelige Virken gaar i ét.

Det sociale Spørgsmaal, som Sibbern i nævnte Utopi
% --- p. 103 ---
\setcounter{footnote}{0}%
\opage{103}(hvis Titel er „Meddelelser af et Skrift fra Aar 2135“) behandler,
havde allerede før Revolutionen i 1848 vakt hans
Opmærksomhed. Den første Fortale til Utopien er dateret
1846. Og i et Universitetsprogram fra 1849 opfordrede han
Statsmændene til i Tide at sætte sig ind i Kommunismens
Idé, som er Udtryk for højere, fordelende Retfærdighed, i
Modsætning til Produktion og Konkurrence, der ikke indeholder
nogen Garanti for, at Udbyttet virkelig kommer
store Kredse tilgode. Han lovpriser det sociale Spørgsmaal
i Modsætning til det konstitutionelle (parlamentariske),
som mest optog Sindene. I det Fremtidssamfund, han skildrer
i sin Utopi, er der slet ikke Tale om Stats- og Kirkevæsen,
saa lidt som om Rets-, Handels- og Ejendomsvæsen.
Arbejdet er organiseret saaledes, at alle tager Del baade i
aandelig og i materiel Virken. Udbyttet er fælles. Af offentlige
Forraadshuse tager enhver, hvad han behøver.
Forholdene skildres i det enkelte med mange naive Træk.
Naiviteten viser sig i det store ved, at der ikke spørges om,
hvilken Sikkerhed der er for, at den ideale Organisation
opretholdes, lige saa lidt som om, hvorledes den er bleven
til. Det for alle ideale Fremtidsskildringer vanskelige
Spørgsmaal om, hvorledes Overgangen fra den nuværende
økonomiske og aandelige Kultur til den ideale Tilstand kan
foregaa, undgaar den naive Utopist ved at lade Menneskene
(med Undtagelse af en lille udvalgt Flok) falde i en lang
Søvn, under hvilken de gamle Indretninger falder sammen.
Ved Opvaagningen er man saa ikke til Sinds at oprette
dem igen. Heller ikke spørges der om, hvorledes Driften
til Arbejde og Initiativet til nye Retninger kan bevares
under de nye Forhold.

Sin største Interesse har Sibberns Utopi derved, at medens
Romantiken var vendt mod Fortiden, i hvilken den
fandt umiddelbare Aabenbaringer af det evige, viser den
gamle Romantikers Sind sig nu, med skarp Kritik af Fortidens
Levninger, vendt fremad mod nye Muligheder. Og
som andre Utopier (Fourier's, St. Simon's, Owen's) fra
samme Tid vidner Sibberns om en varm Følelse for dem,
% --- p. 104 ---
\setcounter{footnote}{0}%
\opage{104}der lider under det nærværende Samfundslivs Ufuldkommenheder.
---

Som jeg engang før har udtalt, har næppe hos nogen
dansk Forfatter psykologisk Sans og Menneskelighedsfølelse
været forbundne med saa stor Aandsfrihed som hos
Frederik Christian Sibbern.

Gennem Poul Møller staar Kierkegaard i Forbindelse
med Sibbern. Denne overbragte ham en Hilsen fra Poul
Møllers Dødsleje, og Kierkegaard holdt sig, som vi har set,
en Tid til ham. Dette personlige Forhold har Interesse ved
at vidne om Sammenhængen i dansk Filosofi.

I sine Papirer (udg. af Heiberg og Kuhr. Femte Bind
p. 20) siger Kierkegaard: „Den danske Filosofi, hvis der
engang skal kunne blive Tale om en saadan, vil deri være
forskellig fra den tyske, at den slet ikke begynder med
Intet, eller uden al Forudsætning, eller forklarer Alt ved at
mediere, da den tværtimod begynder med den Sætning, at
der er mange Ting mellem Himmel og Jord, som ingen Filosofi
har forklaret“. Denne Sætning vil, mener han, ved
at være optagen i Filosofien, afgive det behørige Korrektiv
og tillige kaste en humoristisk-opbyggelig Bevægethed over
det hele. --- Hvad Kierkegaard her venter af en dansk Filosofi,
der kunde komme, passer efter min Opfattelse godt
paa den danske Filosofi, som den har udviklet sig fra Holbergs
Tid af. I hele Rækken af danske Filosofer før Kierkegaard
--- Holberg, Sneedorff, Steffens, Treschow, A. S. Ørsted,
Sibbern --- finder vi\footnote[1]{I min Bog \textit{Danske Filosofer} (1909) har jeg nærmere søgt at vise dette, og i \textit{Udvalgte Stykker af dansk filosofisk Literatur} (1910) vil man finde Vidnesbyrd derom. } som Hovedejendommelighed
en Sans for det individuelle og i nøje Sammenhæng dermed
en Vægt paa Erfaringen. Man kan jo ikke forstaa Individualiteten
uden ved Erfaringens Hjælp. Der kan have været
Uenighed om, i hvilket Omfang det problematiske optræder
i Verden. Saaledes mente Holberg, at med Udvidelse
af Erfaringen vilde vi staa overfor endnu flere Problemer,
% --- p. 105 ---
\setcounter{footnote}{0}%
\opage{105}end vi hidtil har at arbejde med, medens den mere optimistiske
Sneedorff mente, at mange af vore nuværende
Problemer vil falde bort ved udvidet Erfaring. Men gennemgaaende
% PRINTED AS IS (read at the image, odd spots o1)
passer den givne Karakteristik. Og Kirkegaard
selv betegner et paradoxalt Højdepunkt af den raadende
Tendens, hvad Fremhævelsen af det individuelle angaar;
men denne fremskredne Stilling, hvad det ene Spørgsmaal
angaar, købes hos ham med Indsnævring af den brede
Erfaringsbasis, som dansk Tænkning tidligere byggede
paa.

Jeg vil her kun tale om Kierkegaard som Filosof, og
Hovedskrifterne er her „Filosofiske Smuler“ og „Uvidenskabelig
Efterskrift“. Dette sidste Skrift, der udkom 1846,
afsluttede den rige Produktion i æstetisk, religiøs og filosofisk
Retning, som Kierkegaard øvede i den første Halvdel
af Aarene 40, og som han senere, hvad der ses af hans
Papirer fra 1849, gjorde Forsøg paa at udlede af en bestemt
forudfattet Plan, et Forsøg, hans store Sandhedskærlighed
dog førte ham til at opgive, idet han blev klar over,
at det var den uvilkaarlige Trang og Evne til at producere,
som var den drivende Kraft. Tilfredsstillelsen af den stærke
Produktionstrang var for ham tillige et Middel til at beherske
det dybe Tungsind, som inderst inde raadede i ham.
„Uvidenskabelig Efterskrift“ indtager en central Stilling i
Kierkegaards Produktion, for saa vidt som det giver et Tilbageblik
paa hans forudgaaende Værker og yder Oplysninger
til ret Forstaaelse af dem. I filosofisk Henseende
har det den store Betydning at klargøre hans Stilling som
Erkendelsesteoretiker og som Livsfilosof, de to Synspunkter,
fra hvilke vi her vil betragte ham. Han havde en Stund
tænkt sig, at dette Skrift skulde danne Afslutningen paa
hans Produktion i det hele. Men en Uddybelse af hans religiøse
Liv førte i de følgende Aar til en ny og vældig Udfoldelse
af religiøst Forfatterskab, som igen førte til, at han
kom til at føle sig i stærkere og stærkere Modsætning til
den sædvanlige Opfattelse af Kristendommen og til den bestaaende
Kirke.

% --- p. 106 ---
\setcounter{footnote}{0}%
\opage{106}--- For Kierkegaard var Tænkning karakteriseret ved, at
den stadig søgte og fandt Vanskeligheder, hvor den jævne
Forstand, og ofte ogsaa den spekulative Systematisering,
fandt Vejen jevn og farbar. Det er Vanskeligheder, Gaader,
der i deres højeste Grad bliver til Paradoxer, som holder
Tanken i Aande. Og da Livet stedse frembyder nye Gaader
og Opgaver, hvis Løsning kræver aandeligt Arbejde,
kan intet System fuldendes. Livet er ikke afsluttet, men
udfolder sig stedse videre, og Tanken, der forklarer Livet,
kommer stedse bagefter. Vi lever forlænds, men vi forstaar
baglænds.

Allerede Sibbern havde paa lignende Maade stillet sig i
Modsætning til spekulativ Filosofi. Men hos Kierkegaard
udformes Tankegangen med større Eftertryk og Konsekvens.
Han greb med begejstret Tilslutning Lessings Ord
om den evige Stræben efter Sandhed, som det højeste,
Mennesker kan naa, og han fandt i dette Ord et Udtryk
for den Selverkendelse, ethvert ærligt Menneske maa komme
til, og derfor et Udtryk for etisk Livsanskuelse. Men
ogsaa rent teoretisk tager Kierkegaard Spørgsmaalet op.
Han forlanger en Revision af Erkendelsesproblemet i
Kant's Aand og staar derved som Forgænger for det store
Arbejde, der hen imod Aarhundredets Slutning blev gjort
indenfor Erkendelsesteorien.

Han spørger, om Videnskabens Grundtanker, de saakaldte
Kategorier, der stadig lægger sig for Dagen i det
specielle videnskabelige Arbejde, kan findes gennem abstrakt
Spekulation, eller om de skyldes Erfaring. Det var et
Spørgsmaal, han ønskede at behandle nærmere. Dette Ønske
fik han ikke opfyldt, da hans filosofiske Produktion
afbrydes formedelst den ovenfor nævnte Vending i hans
aandelige Liv. Men i Grunden har han givet Svaret paa sit
Spørgsmaal ved at hævde, at Forstaaelse altid kommer
efter Oplevelse.

Jeg tør maaske her indskyde, at jeg har forsøgt at drøfte
dette Spørgsmaal i min Bog „Den menneskelige Tanke“,
idet jeg viser, at Kategorierne findes ved Analyse af det
% --- p. 107 ---
\setcounter{footnote}{0}%
\opage{107}videnskabelige Arbejde, saaledes som det udvikler sig gennem
Tiderne. Jeg lægger særlig Vægt paa, at hvad der er
selvfølgeligt for en Tids Forskere, ikke altid har staaet saaledes
for tidligere Tiders Forskere. En Fremstilling af videnskabelige
Grundtanker, en Kategorilære, kan derfor
kun betyde et Stadium i menneskelig Tænknings Udvikling.
Kategorier kan opstaa, og de kan dø. --- Ogsaa min
Religionsfilosofi bygger paa denne Opfattelse af Forholdet
mellem Oplevelse og Forstaaelse.\footnote[1]{Naar Kierkegaard mente, at der for Filosofen vil gøre sig „en humoristisk-opbyggelig Bevægethed“ gældende i Sindet, naar han bliver sig Erkendelsens Begrænsning bevidst, saa har den Erfaring, jeg under mit Tankearbejde har gjort, for mit Vedkommende bekræftet dette, og jeg har baade i \textit{Den menneskelige Tanke} og senere i \textit{Den store Humor} (se dette Skrift p. 117---123) søgt at give Udtryk for en Stemning af nævnte Art.}

Da de Forudsætninger, hvorpaa al Begrundelse hviler,
ikke selv kan begrundes, kommer man i Følge Kierkegaard
kun ved et Spring ind i videnskabelig Tankegang.
Der gør sig her, hvad Kierkegaard kalder en Begyndelsens
Dialektik, gældende. Jeg maa begynde min Eftertanke paa
et enkelt Punkt, med en enkelt Antagelse, der først senere
hen, gennem fortsat Erfaring og Eftertanke, kan komme
til at staa som begrundet, altsaa foreløbig staar som vilkaarlig.

Paa dette Punkt har Kristian Kroman udviklet Kierkegaards
Tanke videre, idet han i sin Bog „Vor Naturerkendelse“
med stor Skarphed hævdede, at Videnskabens første
Principer fastslaas ved en vilkaarlig Akt i Postulatets
Former. Om min egen Stilling til dette Spørgsmaal vil jeg
senere komme til at udtale mig. ---

Naar Kierkegaard betegnes ikke blot som Erkendelsesteoretiker,
men som Livsfilosof, betyder det, at han har
fremdraget Synspunkter, der er af Betydning for Menneskers
Maade at tage Livet paa, og at disse Synspunkter ere
fremgaaede af dyb psykologisk Indsigt og af egen dyrekøbt
Erfaring. I vor Tid grasserer der meget, især i Tyskland,
% --- p. 108 ---
\setcounter{footnote}{0}%
\opage{108}som kalder sig Livsfilosofi, men som kun er Forsøg paa i
en Haandevending at gøre spekulative og dogmatiske Antagelser
gældende, uagtet de har mistet deres teoretiske
Grundlag. Det er Forsøg paa at løse Problemerne uden at
stille dem for Alvor. Men netop den Alvor og Energi, med
hvilken Kierkegaard har skildret Livets Forhold til Tanken,
og særlig til de Tanker, Livet selv ikke kan undvære,
har nu ogsaa i Udlandet mere og mere henledet Opmærksomheden
paa ham og paa hans Synsmaader, som jeg nu
skal forsøge at karakterisere.

Som ovenfor udtalt danner Kierkegaard Afslutningen paa
den karakteristiske Tendens i ældre dansk Filosofi til at
hævde det individuelle og lægge Vægt paa Erfaringer, men
en Afslutning, der er karakteriseret ved, at Individualiteten
(eller, som Kierkegaard udtrykker sig, Subjektiviteten)
pointeres til det yderste, medens Erfaringselementet samtidig
begrænses i sin Betydning.

Som det i Følge Kierkegaard var ved et Spring, vi kom
ind i Tankens Verden, er det ogsaa ved et Spring, vi kommer
ind i den Livsanskuelse, vi havner i. Det er personlig
Trang, Følelsen af aandelig Nød, til sidst Fortvivlelsen, der
fører til at gribe det Nødanker paa Livets stormfulde Sø,
som enhver Livsanskuelse er. Her tilspidser den Lessingske
Sætning om den evige Stræben sig til den Sætning, at Personligheden
(Subjektiviteten) \textit{er} Sandheden. Det vil sige,
at den Sandhed, Mennesket kan finde paa Livets Vej, ikke
har anden Grund end netop Personlighedens egen Trang.
Ingen Videnskab kan hjælpe os her. Enhver Livsanskuelse
behøver Tanker, og for saa vidt Erkendelse. Men Kierkegaard
gør saa her en skarp Forskel mellem væsentlig og
uvæsentlig Erkendelse, en Forskel, der bestemmes ved
Modsætningen mellem de Tanker, der er os nødvendige
for at kunne holde os oppe under Livets Kamp, og saadanne
Tanker, som kun intellektuel Interesse fører til at
fordybe sig i. Til uvæsentlig Erkendelse regner han al objektiv
Videnskab, teoretisk Filosofi, Historie og Filologi,
saavel som Naturvidenskab. Han anede, at Naturvidenska%
% --- p. 109 ---
\setcounter{footnote}{0}%
\opage{109}ben ved sine nøje sammenhængende Synspunkter og Metoder
kunde blive en lignende Fare for hans Synsmaade,
som den spekulative Filosofi og Teologi, han længe havde
bekæmpet. Ud fra den Betragtning, at Livet er noget individuelt,
ikke tilgængeligt for abstrakte Teorier, løsnede han
det fra enhver Sammenhæng med Videnskaben.

Det Spørgsmaal fremstiller sig da, om enhver Livsanskuelse
er ligesaa god som enhver anden, eller om der gives
nogen Maalestok, efter hvilken her en Vurdering kan ske.
Hvad er højt og hvad er lavt, --- hvad er stort og hvad er
smaat paa dette Omraade?

Allerede tidligere var i dansk Aandsliv Spørgsmaalet om
en saadan Maalestok kommet frem. Det var i Striden mellem
Baggesen og Grundtvig, i sidste Grund en Strid mellem
Filosofi og Teologi. Grundtvig var enig med sin Modstander
i at føre Kampen paa rent menneskelig Grund. „Ej
spørger her vi stort om Konfessioner!“ udraaber han. Men
naar Striden mellem de to aandelig bevægede Mænd tabte
sig i Sandet, var det, fordi de ingen fælles menneskelig
Maalestok kunde finde.\footnote[1]{Se om denne Strid \textit{Danske Filosofer,} p. 66---72.} Kierkegaard har forsøgt at hævde
en saadan Maalestok, og hans Betydning som Livsfilosof
beror ikke mindst herpaa.

Kierkegaards Maalestok er hentet fra selve Begrebet Personlighed.
Ved at bestemme dette Begreb vil man altid
komme tilbage til, at det forudsætter Energi til at sammenfatte
en større eller mindre Mangfoldighed af Tilskyndelser
og Oplevelser, saaledes at en samlet Optræden overfor
Omverdenen kan blive mulig. Jo fyldigere og mere
omfattende det Indhold, der skal sammenfattes, er, des
større Energi forudsættes der, naar den inderlige Koncentration
ikke skal falde bort. Det er i Grunden det fra Leibniz
og Kant stammende Syntesebegreb, som her ligger til
Grund, og Livsfilosofien vil her altid vise sig at staa i Gæld
til den teoretiske Filosofi.

Ifølge denne Maalestok inddeler saa Kierkegaard de forskellige
Livsanskuelser eller, som han med et uheldigt Ord
% --- p. 110 ---
\setcounter{footnote}{0}%
\opage{110}benævner dem, de forskellige „Stadier“. --- Koncentrationen
kan begrænses til det enkelte Øjeblik og det, der kan udfylde
dette (det „æstetiske“ Stadium). --- Man kan være
tilsyneladende optagen af den enkelte Situation og dog i
sit Indre staa fri overfor den. Det er den i Alvor skjulte
Spøg, der karakteriserer Ironien. --- Man kan i inderlig
Alvor fordybe sig i sin Pligt (det etiske Stadium). --- Man
kan betragte Livets Indhold og Opgaver som ude af Stand
til at opfylde Sindets hele Tragten (Humor). --- Man kan
endelig se Livets egentlige Opgave ligge ud over al Erfaring,
Tiden som forsvindende i Forhold til Evigheden. Det
er det religiøse Stadium, der i Kristendommen fremtræder
i paradoxal Form, fordi der det evige selv fremtræder i
Tiden og i en menneskelig Skikkelse, man kan møde paa
sin Vej.

Som man ifølge Kierkegaard ved et Spring kommer over
fra Livet i Tankens Verden til den Tro eller Livsanskuelse,
efter hvilken man fører sit Liv, saaledes er det ogsaa ved
et Spring, man kommer over fra den ene Livsanskuelse til
den anden. Det er under alle Former Kontinuiteten, Kierkegaard
vil til Livs --- mærkværdig nok netop ud fra den
indre, sluttede Sammenhæng, der for ham med Rette karakteriserer
Personligheden!

Baade som Erkendelsesteoretiker og som Livsfilosof er
Kierkegaards store Betydning betinget ved den paradoxale
Skarphed i Problemstillingen. Det gjorde ham til en streng
Lærer for den Generation, der udviklede sig under hans
Indflydelse. Og Sporene heraf vil kunne mærkes, selv hvor
Udviklingen har ført langt bort fra hans Stade.

Fra mit nuværende Stade vil jeg bemærke følgende.

Ethvert Spring maa for Tænkningen staa som en Opgave,
ikke som et sidste Ord. Hvorfor springer man netop
paa dette bestemte Sted og i denne bestemte Retning? Og
betyder selve Springet ikke en Kontinuitet, selv om det betyder
et Fald! --- Filosofien, som har den Opgave at undersøge
Forholdet mellem de forskellige Videnskaber, maa
ogsaa drøfte Springene, de bratte Overgange mellem disse.
% --- p. 111 ---
\setcounter{footnote}{0}%
\opage{111}Den finder disse Overgange betinget ved Optagelse af nye
Forudsætninger, men saa bliver Nødvendigheden af disse
nye Forudsætninger at undersøge. Og det samme gælder
% PRINTED AS IS (read at the image, odd spots o1)
ved Overgangen fra Tanke til Liv, en Overgang, som nøvrig
stedse igen forudsætter en forudgaaende Overgang fra
Liv til Tanke. Her er altsaa en stadig Vexelvirkning.

Ved Overgangen fra Tanke til Liv kan der ikke, som
Kierkegaard vilde, en Gang for alle skelnes mellem væsentlig
og uvæsentlig Erkendelse. Kierkegaards Standpunkt
betegner her en Indsnævring i Forhold til det brede
Grundlag, paa hvilket tidligere danske Tænkere, især Mænd
som Holberg og Sibbern, opbyggede deres Livsanskuelser.
Verdensbilledets Udvidelse var for Holberg, den store Udvikling
i Naturen fra Kaos til Kosmos var for Sibbern meget
væsentlige Elementer baade i deres Tro og i deres Viden.
Hvilken Betydning har da ogsaa Videnskaben, naar
den ikke, maaske ad mange Omveje, til sidst faar Indflydelse
paa menneskelige Livs- og Verdensanskuelser? Og
saadan Indflydelse øves da ogsaa stadig og i al Stilhed,
selv om kun et i Filosofiens Histories Skole uddannet Blik
opdager den. ---

Jeg har nu for sidste Gang talt om de to danske Tænkere,
der har været mest for mig. Og det er mig en Glæde, at
Opfordringen til endnu en Gang at sysle med dem kom til
mig fra en Kreds af unge Studenter.

% --- p. 112 ---
\setcounter{footnote}{0}%
\opage{112}\essayhead{FILOSOFI \\ OG KIRKELIG RATIONALISME}

Dr. Michael Neiiendams interessante Afhandling om
Christian Bartholin har klart belyst en kirkelig Retning,
som fordetmeste af Eftertiden blev haanet eller dog lidet
forstaaet, men som nu underkastes en rent historisk Belysning.
Hovedpersonens Karakter og Tankegang, der er
saa karakteristisk for hans Tid, det 18. Aarhundredes sidste
Halvdel, træder klart og tydeligt frem.

For Filosofien har den kirkelige Rationalisme i sig selv
meget liden Interesse. Den hviler paa en Populærfilosofi
uden klare Principer, og den bygger gerne paa Smuler, der
var faldne fra de egentlig filosofiske Forskeres Bord. Derimod
er det af Interesse at se, hvad Afhandlingen ogsaa
viser, at strengere filosofisk Tænken medvirkede til Opløsningen
af det uholdbare Forsøg paa at løse det religiøse
Problem, som ligger til Grund for kirkelig Rationalisme.
Men i det Galleri af Personligheder, som Afhandlingen
fremdrager, og som, saavidt jeg kan se, yder gode Bidrag
til Forstaaelse af det 18. Aarhundredes aandelige Forhold
i Danmark, forekommer to Personligheder, jeg ikke finder
grundig og retfærdig behandlede. Der er fra et filosofisk
% PRINTED AS IS (read at the image, odd spots o1)
Synspunkt væsenlige Indvendinger at gøre mod Forfatterens
Fremstilling.

1. Jeg ser ikke, at det havde været nødvendigt for Forfatteren
at omtale \textit{Lessing.} Der har, som han selv siger,
ikke paa hin Tid fundet nogen Indflydelse af Lessing paa
dansk Teologi Sted, afset fra, at hans polemiske Skrifter
lejlighedsvis benyttes. En saadan Indflydelse øvedes derimod
senere gennem Kierkegaards forstaaende og begej%
% --- p. 113 ---
\setcounter{footnote}{0}%
\opage{113}strede Omtale af Lessing, og hvis Forfatteren havde lagt
Mærke til den, vilde han have udtalt sig anderledes om
Lessing, end han gør. Han lader ham samvirke med Oplysningstiden
og styrke dennes kritiske Sans. Dette er i
hvert Tilfælde ikke den fulde Sandhed. Lessing hører, ligesom
fra andre Sider Rousseau, Hume og Kant, til de Personligheder,
der --- alt imedens de i mange Henseender er
gennemtrængte af deres Aarhundredes Aand --- peger ud
over deres Samtid til rigere og dybere Opfattelse af Problemerne.
Hvad der er det mest karakteristiske for Lessing,
er hans Sans for ægte og oprindeligt Aandsindhold,
hvor det saa findes. Denne Sans førte ham til i sin dramatiske
Kritik („Hamburgische Dramaturgie“) at fremdrage
Shakespeares vældige Poesi i Modsætning til fransk-klassisk
Tragedie; --- den førte ham (i „Laokoon“) til at drage
en Grænse mellem Billedkunst og Poesi; --- og den førte
ham til (overfor Rationalisme som overfor Ortodoxi) at
henvise til det religiøse Livs inderlige og umiddelbare Bevægelser
i Modsætning til den Reflexion og den Lærdom,
som religiøse Dokumenter kan give Anledning til. Men
Lessing indskærpede tillige, at naar man overhovedet vilde
anvende Tanker i sin Livsanskuelse, skulde man tænke
dem til Ende. Der var to forskellige Strømninger, som
rørte sig hos ham lige fra hans Ungdom af, den ene næret
af hans historiske Sans, den anden af hans filosofiske
Tænkning. Det lykkedes ham ikke at bringe disse to Strømninger
til at forene sig. Men hvem er det i Grunden lykkedes?
--- Ved at gøre dem begge gældende overfor det religiøse
Problem har han hævdet Personlighedens og den
selvstændige Tænknings Ret overfor de uklare Mæglingsforsøg,
som ikke blot den kirkelige Rationalisme, men
ogsaa selve Ortodoxien stedse gør sig skyldige i, og han
faar derved en Betydning, der gaar langt ud over hans Tid.

De Mæglingsforsøg, den kirkelige Rationalisme foretog,
var Lessing mere forhadte end den gamle Ortodoxi, fordi
de (som han skrev til sin Broder) „under Paaskud af at
ville gøre os til fornuftige Kristne gør os til ufornuftige
% --- p. 114 ---
\setcounter{footnote}{0}%
\opage{114}Filosofer“. Men ogsaa Ortodoxiens (dengang) moderne Repræsentanter
(som Schumann og Goetze) maatte for ham
høre til de uklare Mæglere. Den frie Religiøsitet, han haabede
paa, vilde --- mente han --- snarere udfolde sig, naar
man lod den gamle Tro aabent og ærlig føre Ordet. Han
vilde ikke have det oprindelige religiøse Indhold sammenbundet
med spekulative Idéer, saaledes som det tidlig var
sket i Kirken. Positiv Religion er noget Konkret og Historisk
og derpaa kan ingen evige Sandheder bygges. Positiv
Religion er helt igennem bygget paa, hvad der antages at
være historisk Sandhed, og kan hverken føre til eller udledes
af almene Idéer. Religionens egentlige Bevislighed er
inderlig Tilegnelse af et historisk Indhold. --- Hvis Lessing
havde været mere Psykolog, end han var, havde han kunnet
tale tydeligere. Dog udtaler han sig tydelig nok mod
Intellektualismen paa det religiøse Omraade.

Det var ganske vist Lessings Tro, at Religionen maatte
komme til at gennemgaa en Udvikling, hvorved den mere
og mere kom ud over det rent historiske og derfor mere
eller mindre vilkaarlige Grundlag. Men denne Udvikling
maatte ske langsomt. Han raadede fra at slukke sine Lys,
før Solen var staaet op. Og denne Varsomhed havde netop
sin Rod i hans Sans for personlig Tilegnelse, for hvad man
senere har kaldt personlig Sandhed. Hans historiske Sans
lærte ham desuden, at hverken de saakaldte Aabenbaringssandheder
eller de saakaldte Fornuftsandheder var de
samme til alle Tider.

Overfor den ortodoxe Hauptpastor Götze hævdede Lessing,
at Kristendommen ikke stammer fra skriftlig Overlevering,
men fra de Begivenheder, disse Beretninger handler
om, og det var ogsaa overfor ham, at Lessing fik Anledning
til at udtale den berømte Sætning, at den stadige Stræben
efter Sandhed er for et Menneske at foretrække for den
stadige Besiddelse af Sandheden. Senere fik han over Jacobi's
sentimentale Halvteologi Lejlighed til at forsvare den
store og stille Tankebevægelse, der ligger bag ved Spinoza's
krystalliserede Systematik. Begge Steder har hans Geniali%
% --- p. 115 ---
\setcounter{footnote}{0}%
\opage{115}tet i Forening med hans Kongenialitet fremdraget betydningsfulde
Synspunkter. Det er karakteristisk for Lessing,
at det er to Sammenstød, der har fremkaldt hans mest betydningsfulde
Udtalelser. Han var den store Søger, hørte
ikke til nogen Konfession eller til noget System. Og det er
i hvert Tilfælde ikke rigtigt, som det er sket i Dr. Neiiendams
ellers saa interessante Bog, at fremstille ham blot
som Repræsentant for Oplysningstiden.\footnote[1]{Se nærmere om Lessing \textit{Den nyere Filosofis Historie}\textsuperscript{3}, III p. 16---25.}

2. Den anden Skikkelse, som er fortegnet i den nys
nævnte Bog, er den danske Forfatter \textit{Johannes Boye} (1756
--- 1830), hvis Skrift („Statens Ven“) dér kaldes „den vigtigste
Kilde til Forstaaelse af den hjemlige Oplysning“. Men
Boye kommer ikke til sin Ret ved at betegnes som Oplysningsmand.
Han er en Renæssancenatur, hævder antik
Klarhed og Kraft og paaviser, hvorledes etiske Idéer udvikler
sig under Kampen for Tilværelsen som den eneste
Mulighed til at gøre Ende paa fysisk Magts Eneherredømme.
Skønt han bekæmper Kant's Begrundelse af Moral ud
fra ren Fornuft, er hans mandige Livsopfattelse, ifølge
hvilken aandelig Udvikling gaar gennem Strid og Arbejde,
ganske i Kant's Aand. „Naturen,“ siger han, „optugter sine
Børn strengt, men det sker for at frembringe Helte.“ --- Boye
lægger stor Vægt paa de historisk-psykologiske Forudsætninger
for Tilblivelsen af en rent menneskelig Moral. Men
det har Kant selv gjort paa et betydningsfuldt Stadie af
den Udvikling, hans etiske Idéer er undergaaet. Menneskets
Anlæg som Fornuftvæsen kan efter Kant kun komme
til Udvikling i Løbet af Slægtens Gang, der gennem
Striden mellem Lidenskaber fører til Idéen om et Samfund,
indenfor hvilket den Enkeltes Frihed gennem Lovbundethed
bliver forenelig med Andres lige Frihed. Uden
Strid vilde de menneskelige Evner ikke udvikles. Træerne
i Skoven, siger Kant, søger at berøve hverandre Lys og
Luft, men netop derved nødes de til at søge Lys og Luft i
højere Regioner. (Smlgn. nærmere om dette Punkt Den
% --- p. 116 ---
\setcounter{footnote}{0}%
% CHECK p. 116: notes paired with the marks in order (the layer misread a number)
\opage{116}\textit{nyere Filosofis Historie}\textsuperscript{3}. III, p. 74---78). Hverken Kant selv
eller hans fremragende danske Fortolker Anders Sandøe
Ørsted var klar over Forholdet mellem psykologisk-historisk
Undersøgelse og etisk Principteori --- og dette Forhold
er et Problem den Dag idag. Ørsted gik ud fra den rene
Fornuftmoral som fuldt færdig, Boye fra den stadige Kamp
for Tilværelsen. Og det var ingenlunde et Nederlag, Boye
led ved at fremdrage denne Side af Problemet.\footnote[1]{Se nærmere om den mærkelige Tænker, Boye var, \textit{Danske Filosofer,} p. 33---35. --- Omtrent samtidig med mit Forsøg paa Rehabilitation af Boye (i nævnte Bog) gav \textit{Vilhelm Andersen} en meget interessant Karakteristik af ham i „Typer og Tider“. }

% --- p. 117 ---
\setcounter{footnote}{0}%
\opage{117}\essayhead{EN MODERNE BUDDHIST}

\essaysource{Arthur Pfungst: Gesammelte Werke. I. 1926.}

Gennem min afdøde Kollega, den berømte Paliforsker
Viggo Fausbøll lærte jeg den mærkelige Mand at kende,
hvis samlede Værker, der skal udgøre tre Bind, nu begynder
at udkomme. Pfungst's Bekendtskab med Fausbøll
stammede fra hans store Interesse for Buddhismen, som
for ham ikke blot havde historisk Betydning, men i væsentlige
Træk bestemte hans Livsanskuelse og gav ham
Udtryk for, hvad han ad sin egen Vej var kommen til.
Fausbøll udgav som bekendt det buddhistiske Skrift Dhammapada
(Den sande Læres Ord) efter de Manuskripter,
Rask i sin Tid havde bragt hertil fra Ceylon, og senere Jatakabogen
efter samme Kilder. Der var en stor Modsætning
mellem de to Mænd, hvis Venskab var bygget paa den fælles
Interesse for de gamle indiske Visdomsbøger. Fausbøll
var en stille Lærd, optagen af sit videnskabelige Arbejde,
der gav ham et smukt Navn i Indologernes Kreds og endog
gjorde det kendt og skattet blandt de indiske Buddhister,
skønt han ikke hyldede deres Lære. Han beundrede Buddhismen,
ikke mindst fordi der ingensinde under dens
Herredømme var øvet saadanne Afskyeligheder og Grusomheder
som under Kristendommen. Men han var ikke Buddhist,
da han ikke kunde opgive Troen paa Sjælens Udødelighed
og paa en personlig Gud. Denne sin Anskuelse udtalte
han aabent i en lille Selvbiografi, han i Aaret 1900
skrev i det buddhistiske Blad „Journal of Maha-Bodhi Society“.
Pfungst beundrede Aabenhjertigheden af denne Udtalelse,
da det nævnte Blad var et Propagandablad for Buddhismen,
og da de Lærdomme, Fausbøll holdt fast ved, var
% --- p. 118 ---
\setcounter{footnote}{0}%
\opage{118}Kætterier fra et buddhistisk Standpunkt. At Bladet uden
videre optrykte Fausbøll's Udtalelser, var jo netop et Tegn
paa den store Tolerance, som Fausbøll priste hos Buddhismen.
Arthur Pfungst var en Verdensmand, rig, Arving til
en stor Fabrik i Frankfurt am Main, stærkt optagen af alle
Bevægelser indenfor de forskellige Kulturomraader, politisk
og religiøst radikal, Digter og Forretningsmand paa én Gang.
Hans Forretningsrejser førte ham rundt i Verden, og han
kom ofte her til Norden, hvor han aldrig glemte at besøge
Fausbøll. Han var livlig og munter og kun af Udtalelser,
der fremkom lejlighedsvis, lærte man, at der bag hans
Munterhed og hans spillende Vid laa den Overbevisning, at
Livet ikke var værdt at leve. Buddha's Lære var Pfungst's
egen Religion, og han har virket for den ikke blot som
Digter, men ogsaa i en Række af Afhandlinger, som han
1904 samlede i Bogen „Aus der indischen Kulturwelt“. Man
kunde spørge, hvorledes en saadan Dobbelthed, en saadan
Modsætning mellem det Indre og det Ydre var mulig. Men
efter hvad jeg har kendt til ham, tror jeg, at en af hans
nærmeste Venner, Dr. Angermann, der har skrevet den
„biografiske Indledning“ til det første Bind, har Ret, naar
han udtaler:

„For Pessimisten staar den ubehagelige Side af Tingene
fast som det mest Selvfølgelige i Verden, saa at der fra den
Side ingen Overraskelser kommer for ham, medens der fra
den anden Side kan komme gode Overraskelser og Øjet
kan skærpes for, hvad glædeligt og spøgefuldt Livet kan
bringe. Forkastelsen af denne Verden staar fast for Pessimisten,
men des mere véd han, at han indenfor dens Rammer
ikke tør gøre Livet vanskeligere for sig selv og andre,
end det i Forvejen er. Med andre Ord, Pessimistens Humor
er den dybeste og fineste; den er efter sit inderste Væsen
Vismandens rene og gode Sindelag over for Andre og sig
selv. Derfor manglede hos Pfungst ogsaa alt Skarpt og Saarende,
og netop det priste Alle, der kendte ham, som den
skønneste Egenskab ved hans Humor.“ ---

Det er vel ikke nødvendigt at bemærke, at Ordet Humor
% --- p. 119 ---
\setcounter{footnote}{0}%
\opage{119}i denne Udtalelse bruges i anden Betydning end den populære,
om Evne til at faa Folk til at le. Det vil heller ikke
være nødvendigt at minde om, at der gives flere forskellige
Former eller Arter af Humor foruden den ejendommelige
Art, Pfungst's Venner fandt hos ham. Og naar vi om lidt
skal høre nærmere om Pfungst's ejendommelige Forhold
til Buddhismen, vil det vise sig, at hans Humor i Virkeligheden
har haft andre Elementer i sit Grundlag end den
rene Pessimisme, Elementer, der bringer den i nært Slægtskab
med den af Thackeray og George Eliot beskrevne Humor,
som skyldes en Forening af komisk Sans og Sympati
(wit and love).

Det har hos Pfungst neppe været Studiet af Buddhismen,
der gjorde ham til Pessimist. Heller ikke har Læsningen af
Schopenhauer været afgørende i denne Henseende. Heller
ikke har det været egne Skuffelser og Sorger, der har ført
ham til hans Totalfølelse ved Livet. Han har, saavidt det
vides, ikke kunnet sige med Edgar i Kong Lear, at han er
kommen under Skæbnens Slag og gennem egne Sorger har
lært at have Medfølelse med Andre. Og der har næppe hos
ham, som hos Schopenhauer, tidlig rørt sig ubændig sanselig
Trang og en ligesaa ubændig Hidsighed, som kunde lade
hans Fantasi opfatte Tilværelsen som pinefuld Brynde. Hvis
% PRINTED AS IS (read at the image, odd spots o1)
man ud af hans store episke Digt „Lakaris“, som jeg siden
skal omtale, vilde slutte, at en stærk Livstrang hos dets Forfatter
er blevet kuet gennem Skuffelser og Sorg, saa har
hans Nærmeste i hvert Tilfælde ikke vidst noget derom.
Hans Biograf mener, at hans Livsanskuelse er bleven
grundlagt i de saa ofte skæbnesvangre Overgangsaar, i
hvilke der ofte er stor Modtagelighed for Livets tragiske
Sider. Den Verdenssmerte, som herved let opstaar, og som
for de Fleste kun er en Overgang, har hos Pfungst faaet
blivende Betydning, og den fik tidlig Udtryk i hans Digte.
Han var voxet op i et rigt og lykkeligt Hjem og har havt
en lys Barndom. Moderens Ømhed og Livslyst har kunnet
bøde paa den energiske Faders pessimistiske Tendenser.
Det maa være gaaet Arthur Pfungst paa lignende Maade
% --- p. 120 ---
\setcounter{footnote}{0}%
\opage{120}som det efter Legenden gik Buddha selv. Buddha levede jo
i Pragt og Herlighed, indtil han en skønne Dag fik Øje paa
den Elendighed, der kan være i Verden, paa Sygdom, Alderdom
og Fattigdom. Det Blik glemte han aldrig. Ogsaa
Pfungst var overbevist om, at han havde kastet et Blik ind
i Livet og dets Forhold, et Blik, der maatte være bestemmende
for, i hvilket Lys alt, hvad Livet bragte, maatte staa.
Men dette Blik affødte hverken Bitterhed eller Passivitet.
Han skal selv have udtrykt sit Standpunkt saaledes, at han
var Pessimist i sin Tanke, Optimist i sin Gerning. Han trak
sig ikke tilbage fra Livet, men førte en Kamp for at lindre
menneskelig Lidelse og navnlig for at modarbejde Undertrykkelse
og Tynge fra Statens eller Kirkens Side overfor
fri Stræben paa de forskellige Omraader. Som Frankfurter
følte han sig i mange Henseender i Modsætning til det
preussiske Væsen, og han talte med stor Glæde om, hvorledes
det var lykkedes ham at afværge et Forsøg fra den
preussiske Regering paa at faa Myndighed over et Legat,
der var stiftet til Bedste for et frireligiøst Samfund. Jeg fik
gennem ham engang optagen en Artikel i „Frankfurter Zeitung“
om Forholdene i Nordslesvig, og han anmeldte selv i
dette Blad mit Foredrag om „Fichte's Reden an die deutsche
Nation og de danske Nordslesvigere“. I sit Tidsskrift „Das
freie Wort“ kæmpede han mod Autoritetstro og Klerikalisme
saavel som mod politisk Tvang. Han følte sig vel som
en Fremmed overfor Menneskers naturlige Tro paa, at Alt
vil blive godt, og at det er til for deres Skyld. Et af hans
Digte har til Overskrift: „Ich bin ein Fremdling auf der
Welt gewesen“. Men naar det gjaldt at bekæmpe Lidelse og
Uret, følte han sig ikke som Fremmed. Selv om han ikke
havde den naive Optimists Tro paa, at der kunde vindes
Sejr én Gang for alle, saa svækkedes derved ikke hans Arbejds-
og Kamplyst. Kampen var for ham skønnere end
Sejren. Et af hans Digte har netop den Overskrift: „Der
Kampf ist schöner als der Sieg“.

Som saa mange orientalske Religioner har Buddhismen
en passiv, selvopgivende Karakter. Den prædiker Medliden%
% --- p. 121 ---
\setcounter{footnote}{0}%
\opage{121}hed, men den advarer tillige og især mod at knytte sit
Hjerte til, hvad der er endeligt og derfor tilsidst uden Realitet.
Det hedder i Dhammapada (Tuxens Oversættelse)
(Kap. 16): „Kærlighed avler Bekymring, Kærlighed avler
Frygt, Den, der har frigjort sig for Kærlighed, kender hverken
% PRINTED AS IS (read at the image, odd spots o1)
Bekymring eller Frygt.' Og fremdeles (Kap. 26): „Den,
der opgiver al Forbindelse med Mennesker og al Forbindelse
med Guder, --- den, som er uden al Forbindelse, ---
ham kalder jeg en Brahman.“ Men ligesom moderne Kristendom
kun blev mulig ved, at Forventningen om Kristi
snare Genkomst efterhaanden gav nogen Plads for andre
Forventninger og for Stræben efter Udvikling af menneskelig
Kultur, saaledes maa ogsaa en moderne, europæisk Buddhist
optage Elementer i sin praktiske Livsførelse, som ikke
fandtes i den oprindelige Buddhisme. Intellektuelle, kunstneriske,
sociale og politiske Opgaver faar positiv Værdi, en
Værdi, som Higen mod Nirvana eller Længsel mod et hinsidigt
Himmerige ikke kendte. De orientalske Verdensreligioner
opstod under Kulturforhold, der var meget forskellige
fra de Kulturforhold, under hvilke deres moderne Tilhængere,
især de vesteuropæiske, lever. Da Tagore var her,
fandt han, at vi her i Vesten har for mange Problemer, Opgaver
og Vanskeligheder. Kom over til os, sagde han; der
bliver Alt saa simpelt. Denne orientalske Tankegang var
Arthur Pfungst langt fra at dele. Han tog Kampen op, hvor
han stod. Men naar han saa priser Kampens Herlighed, er
han mere Optimist, end han som Buddhist har Lov til at
være. Ti selv om han ikke tror paa Sejr for sine humane
Bestræbelser, maa han dog tro paa, at der er Kræfter til at
holde Kampen gaaende, og allerede en saadan Tro eller et
saadant Haab bringer mere Lys ind i Verden, end der
ifølge Buddhismen kan være. For Buddha havde Intet i Erfaringens
Verden en saadan Realitet, at man derpaa kunde
bygge Haab om at naa frem til, hvad der har sand Værdi.
Hermann Oldenberg sammenstiller Buddhismens Negtelse
af enhver positiv Værdi i Verden med Bramaspekulationens
„nej, nej“ overfor ethvert Forsøg paa at sige, hvad Brama er.

% --- p. 122 ---
\setcounter{footnote}{0}%
\opage{122}Den moderne Buddhist staar overfor Nirvana, som den
moderne Kristne staar overfor den Forventning om Kristi
snare Genkomst, der var Grundlaget for Kristendommens
Etik. En Opfyldelse af denne Forventning midt under det
fremstræbende Kulturarbejde, vilde stanse dette, vise det i
dets Betydningsløhed. I en Brevvexling mellem to danske
Præster, jeg for nogle Aar siden anmeldte i „Tilskueren“
(senere i „Oplevelse og Tydning“), fremtræder den
for moderne Kristendom herved begrundede Vanskelighed
paa naiv Maade. Man tør ikke opgive den urkristelige Forventning,
men hvorledes vilde det, hvis denne Forventning
pludselig opfyldtes, gaa med den fremskridende Udfoldelse
af Folkelivet, som de to Venner begge haaber paa?

De to store Verdensreligioner bestaar i nyere Tid kun under
stadig Optagen af nye Elementer, og det store Problem
bliver da, hvorledes disse Elementer kan forenes med de
Elementer, der var de afgørende i de store Begyndelsestider.
I selve den Organisation af Kultus, Dogmatik og
Præsteskab, uden hvilken disse Religioner ikke kunde blive
„positive“ Religioner, ligger der Vidnesbyrd om, at der er
optaget mangt og meget, som ikke fandtes i det Oprindelige,
og som vilde blive forstyrret og skudt til Side, naar
dette Oprindelige igen pludselig blev sat i Kraft. Fra denne
Side har Dostojevski i en Fortælling, hvis Indhold er henlagt
til den middelalderlige Kirke, skildret den Bestyrtelse,
Kristi pludselig Genkomst vilde vække i den bestaaende
Kirke, der har indrettet sig paa langt Hold og aldeles ikke
er indrettet paa, at Kristus selv tager Styret.\footnote[1]{I „Oplevelse og Tydning“ (p. 205), hvor jeg allerede har omtalt denne Fortælling, henførte jeg den til Thor Lange som Forfatter. Men den stammer fra Dostojevski's vidunderlige Værk „Brødrene Karamazov“. (Første Del p. 227 i den franske Oversættelse). } Et tilsvarende
Problem vil kunne stilles for den kirkelige Buddhisme, der
nu bestaar i Østen som en af de ærværdigste religiøse Institutioner
i Verden. Saalænge Religion bygger paa Autoritet
og Overlevering vil et Problem af denne Art atter og
atter komme op. Kun en Livsanskuelse, der bygger paa
% --- p. 123 ---
\setcounter{footnote}{0}%
\opage{123}egen Erfaring om Livet, vil kunne drage Næring af Fortidens
store Tanker og store Skikkelser uden at hildes i
nogen Modsigelse. ---

Men jeg vender tilbage til Arthur Pfungst og til den foreliggende
Del af hans Værker. Denne første Del indeholder
hans Digte. Jeg betragter dem her fra Tankeindholdets
Side, navnlig med Hensyn til, om der i dem antydes Motiver
til den ejendommelige Maade at tage Livet paa, som
Forfatteren tidlig er kommen til. Digtene tyder paa, at selv
om Digterens Biografi ikke kan meddele bestemte Erfaringer,
der kunde have betinget hans Tro, saa har han dog
gaaet en streng Skole igennem og er naaet til sit Stade gennem
Selvbeherskelse og Resignation. Han udtaler i et Epigram,
at der er et Herredømme, i hvilket Herskertrangen
kan finde Tilfredsstillelse, selv om man ingen Undersaatter
har, nemlig Herredømmet over sig selv. Og i et Digt, hvis
Overskrift er „Zurück!“, bekender han, at han engang har
villet være som andre Mennesker, optagen af Livet og dets
brogede Mangfoldighed, glædet sig ved Solskin og ved Elskov,
---

\begin{quote}
jedoch in meinem Herzen rief es „Nein!“\\
mein Genius stand vor mir und sprach „zurück!“
\end{quote}

I det store Digt „Laskaris“, som optager den største Del
af den foreliggende første Del af hans Værker, skildres en
Livsgang gennem Haab, Brøde og Sorg, og som en Art Korsang
forkyndes det atter og atter, som Indledning til de
enkelte Sange, at Undergang og Opløsning truer Alt, hvad
man haaber, Alt hvad man glæder sig ved. Laskaris er en
ung Græker fra Cypern. Han bliver af den gamle Alkymist
Philaletes indviet i Kunsten at gøre Guld, men til Indvielsen
knytter den Gamle det Ord, at den Skat, hans Kunst kan
frembringe, og som alle Mennesker tragter efter, ingenlunde
har gjort ham selv lykkeligere. Derfor opfordrer han
Laskaris til kun at bruge den som Middel til at raade Bod
paa Menneskers Nød og Sorg. Men den unge Mand føler
% --- p. 124 ---
\setcounter{footnote}{0}%
\opage{124}Trang til Magt og Ære og agter ikke paa Philaletes' vise
Tale. Gennem omskiftende Skæbner --- ved August den Andens
prægtige Hof i Dresden og i Karl den Tolvtes Lejr ---
glimrer han ved sin Kunst og skaffer sig Ære og Rigdom.
Men han forraader sin opofrende Hustru for den blændende
Aurora von Königsmark, og senere volder han ved sin Ubesindighed
Hustrus og Barns Død. Tilsidst vil han vende tilbage
til sin Hjemstavn, men omkommer, da Skibet strander
paa Cyperns Kyst.

Det synes at have været Forfatterens Mening at indskærpe,
at kun den kan sige Nej til Livet, som virkelig har
prøvet det. Philaletes var fra Begyndelsen af klar over, at
hans vise Tale ikke vilde holde hans unge Ven fra at
styrte sig ind i Livet. Kun den, siger han, som \textit{har} levet,
kan mismodig sige, at Livets Goder ikke er Møjen værd! Det
er en Advarsel til unge Pessimister, men saaledes som den
udformes gennem Digtet, gælder den ogsaa unge Optimister,
som blindt stoler paa, at Alt maa blive godt. Ud over
Modsætningen mellem de to Stemninger, der vel stedse vil
kæmpe mod hinanden i Menneskers Sind, peger den Tanke,
at det fremfor Alt gælder om at kende Livet gennem egen
Erfaring og eget Arbejde. Men heraf følger naturlig den
Tanke, som Arthur Pfungst ikke vilde gaa ind paa, at Livserfaring
og Livsarbejde aldrig kan afsluttes saaledes, at en
absolut Vurdering af Livet i dets Helhed skulde blive mulig.
I et ufuldendt dramatisk Digt om samme Emne som
„Laskaris“ faar den unge Niceketas af Guderne Ret til Opfyldelse
af tre Ønsker, og hans tredie Ønske bliver da: „Lad
Verden gaa under!“ Men da svarer Guderne: „Du kender
ikke Livet i dets Fylde, ikke Glæden, ikke Kærligheden.
Vent saa idetmindste et Aar, inden du forlanger dit Ønske
opfyldt!“ Niceketas gaar ind herpaa, men erklærer, at hans
Ønske efter et Aars Forløb vil være det samme som nu.
Dramaet blev ikke fuldendt, saa det vides ikke, hvilke
Skæbner den unge Mand skulde gaa igennem, og hvilken
Dom han derefter vilde fælde. Men i Aandernes Sang, der
slutter Forspillet, hedder det: 
% --- p. 125 ---
\setcounter{footnote}{0}%
\begin{quote}
\opage{125}Nicht im Grimme darfst du scheiden,\\
Weh', dein Herz ist liebeleer!\\
Wer da zürnend flucht das Leben,\\
Bringt ja Rettung nimmermehr!\\
--- --- --- --- ---\\
Lern' wie man das Leiden wehrt\\
\textit{Lern' das Leben wieder segnen!}
\end{quote}

Da jeg i Sommeren 1909 besøgte Pfungst i Frankfurt og
tilbragte en Aften med ham, havde jeg Indtrykket af, at der
var kommet noget mere dæmpet og tilbageholdende over
ham. Han var ikke ganske den Samme, som naar han paa
sine Besøg heroppe sprudlede af Vid og Virkelyst, af Iver
for intellektuelle og praktiske Bestræbelser, og naar hans
Humor kunde ytre sig paa saa kostelig en Maade. Maaske
har han dengang allerede lidt af den Hjertesygdom, af hvilken
han tre Aar efter døde.

Gennem Dr. Angermanns smukke Biografi og gennem
fornyet Læsning af hans Digt kan de, der kun kendte ham
gennem enkelte Samtaler, faa et Helhedsbillede, de ikke vil
glemme. Hans Minde vil leve gennem hans Værker og gennem
en Stiftelse han har grundlagt til Fremme af intellektuelle
og filantropiske Formaal. Det vil særlig leve hos dem,
som mødte ham paa hans Vej, og som ved at fordybe sig
i det foreliggende Værk har vundet Forstaaelse af ham.

\piecerule

% --- p. 126 ---
\setcounter{footnote}{0}%

\opage{126}\essayhead{„BAKKEHUS OG SOLBJERG“}

\essaysource{Foredrag for nordiske Lærere og Lærerinder paa Hindsgavl \\ den 30. Juli 1923.}

\romanhead{I}
% CHECK roman numeral read

Skønt jeg ikke er Literaturhistoriker, har jeg dog maattet
love at tale her om min afdøde Ven Troels Lund's
„Bakkehus og Solbjerg“, og det blev sagt mig, at man netop
ønskede det, fordi Forfatteren var min Ven. Jeg tror nu
ikke, at Bogen behøver at ty til en Vens Indlæg for at ses
i sin Betydning. Og desuden kan man meget vel være gode
Venner og dog se meget forskellig paa mange Ting. Jeg vil
da heller ikke lægge Skjul paa, hvilke Indvendinger jeg
har at gøre, baade hvad dens Fremstilling angaar, og hvad
de Synspunkter angaar, der ligger til Grund for hans
Domme.

„Bakkehus og Solbjerg“ handler om den danske Poesis
Historie i den første Halvdel af forrige Aarhundrede. Det
er ikke et lærd literaturhistorisk Værk. Der vil sikkert øves
skarp Kritik overfor det af Literaturkritikere af Faget, og
Værket vil gennemgaa en lignende Skærsild, som hans store
kulturhistoriske Værk „Dagligt Liv i Norden“ i sin Tid
maatte gennemgaa fra kritiske Historikeres Side. Men hans
kulturhistoriske Arbejde vandt dog tilsidst Anerkendelse.
Det blev af vor ypperste kritiske Historieforsker sagt om
ham, at „hans levende Forestillingsevne, hans psykologiske
Skarpblik har tjent ham som Forsker“, og at „han har fundet
meget i Kilderne, som var der, men som Andre først
saa, da han drog det frem“, samt at „Værket vidner om
hans dybe Kærlighed til Folk og Fædreland, præget af at
% --- p. 127 ---
\setcounter{footnote}{0}%
\opage{127}være skrevet af en Mand, der var et fint Udtryk for noget
af det Bedste i den danske Folkeaand“.

En lignende Karakteristik kan der, tror jeg, gives af
hans literaturhistoriske Værk, som han arbejdede paa lige
til sin Død og maatte efterlade ufuldendt. Det er en af sit
Fædrelands Historie og Aandsliv dybt greben Mands Forsøg
paa at udtale, hvad han skylder den rige og store Poesi,
der udfoldede sig her i Danmark i den første Halvdel af
forrige Aarhundrede. Og han lægger ikke Skjul paa, at der
var en indre Trang hos ham, som netop i denne Poesi søgte
og fandt Tilfredsstillelse. Han var en harmonisk Natur.
Hvad Livet bragte ham af Modstand eller Sorg, lukkede
han ikke ude, men lod det fortone sig gennem fortsat Haaben
og Stræben, og han saa stedse Lys forude. Han erklærer
sig da ogsaa i det foreliggende Værk for en „Lysseer“.
Men han vil kun give et Indlæg i den store Debat om Livets
Værd fra sit Synspunkt. „Jeg véd jo egentlig kun,
hvordan det Sete har spejlet sig i mit Lyssind. Men, Herregud,
lad saa Andre lære mørkt. Jeg kan jo kun lære fra
mig, hvad Livet har lært mig.“

Den Periode af Danmarks Historie, hvis daglige Liv hans
kulturhistoriske Værk handlede om (det 16. Aarhundrede),
var en Opgangsperiode, og hans Lyssyn fik i mange Henseender
Bekræftelse ved at sysle med den. Men han saa for
lyst paa den. Den maa dog have baaret Spirer i sig til den
ulykkelige Tid, der i det 17. Aarhundrede brød ind over
Landet og bragte det Undergangen nær. Noget lignende
gælder, efter min Opfattelse, ogsaa om hans literaturhistoriske
Værk: Han har faaet Øje paa en Lysstrøm og følger
den, men spørger ikke, om Strømmens Kilder er tilstrækkelige
til at nære dens fremtidige Løb. I et af sine mærkeligste
Skrifter, som han kaldte „Livsbelysning“, har han
gjort Rede for sin almindelige Opfattelse af, hvorledes
Menneskenes Livssyn har udviklet sig. Dette betinges ved
de store Modsætninger, Erfaringen frembyder mellem Lys
og Mørke, Lyst og Smerte, Godt og Ondt, Liv og Død. I
Udtrykkene Lys og Mørke sammenfatter han alle disse
% --- p. 128 ---
\setcounter{footnote}{0}%
\opage{128}Modsætninger, og det betegner da for ham det største Fremskridt
i Menneskehedens store Udvikling, at man drog den
kække Slutning: ikke Mørket, men Lyset er det Sejrende!
Det er i Forasiens og Grækenlands Religioner, at dette store
Skridt gøres. Man begynder med at antage de lysende Himmellegemer
for Guder. Men derfra højnes Tanken til, at de
kun er Lysgudernes Boliger, og tilsidst opdukker Troen paa
en eneste usynlig Lysets Gud. Og forgæves danner længe
Djævletroen, en Rest af den paa Frygt hvilende Tro paa
mørke Magter, en Hæmning for Lystroen. Ogsaa denne
Hæmning, der var saa stærk i den kristne Middelalder,
overvindes. Og for dansk Aandslivs Vedkommende sker
dette Gennembrud, efter Forfatterens Mening, netop i den
store Poesi, der udfoldede sig i første Halvdel af det 19.
Aarhundrede. I „Bakkehus og Solbjerg“ har han villet godtgøre,
hvorledes den store verdenshistoriske Bevægelse fra
Mørke til Lys gennemløb en af sine Faser indenfor dansk
Aandslivs snevre Ramme.

Bogens Tilblivelse er karakteristisk for dens Forfatter.
Da han havde afsluttet sit store kulturhistoriske Værk,
følte han Trang til at leve i sine egne Erindringers Verden
og tænkte paa at skrive sin Selvbiografi. Denne Plan beskæftigede
ham i flere Aar, og han omtalte den ofte. Men
jo mere han fordybede sig i sine Barndoms- og Ungdomserindringer,
des mere kom han til at dvæle ved de Somre,
han med sin Familie tilbragte ude i „Bakkehuset“ paa
Valby Bakke, og ved de Erindringer, som derved igen vaktes
om en Periode i dansk Digtning, hvor denne Egn spillede
en ikke ringe Rolle. For hans Fantasi myldrede der en
Række af Digterskikkelser frem, som enten havde boet
paa Bakkehuset eller haft deres aandelige Tilhold der. Han
blev træt af at beskæftige sig med sig selv og valgte at
fordybe sig i disse Skikkelser, hvis Værker havde været saa
meget for hans aandelige Udvikling. Han gør dog Uret
mod sig selv, naar han siger, at han overfor de rige Billeder,
hans Erindring og Fantasi vakte hos ham, følte sig
% --- p. 129 ---
\setcounter{footnote}{0}%
\opage{129}som den Træpind, paa hvilken Kagekonerne binder de
Kirsebær, de sælger til Børn. Der var fra hans Ungdom af
et indre Tankeliv i ham, der søgte Udtryk. Jeg erindrer
med Glæde den Tid i mine unge Dage, da vi kom sammen
paa mit Kvistkammer paa Gammeltorv (skraas overfor hans
Faders Gaard) og studerede Platon sammen. Troels Lund
var meget interesseret i dette Studium, som for hans Vedkommende
førte til en aandfuld Doktorafhandling om Sokrates.
Men han var mere optagen af Kunstneren i Platon,
end af Tænkeren Platon. Det var især Platon's Evne til at
lade en Tanke udforme sig i Billede og beaandes af Stemning,
der tiltalte ham. Tænkning og Forskning fyldte hans
Liv fra Ungdommen af, men fik dog især deres Værdi for
ham som Grundlag for Udfoldelse af Stemning og Fantasi.
Og derpaa beroede den Værdi, kulturhistoriske og literaturhistoriske
Emner hele hans Liv igennem fik for ham. Ofte
fik, det kan ikke nægtes, og det gælder ogsaa hans literaturhistoriske
Værk, Digteren i ham Overhaand over Forskeren.
Men i dette Værk, som han udarbejdede i sin høje Alder,
fandt hans Sind Hvile ved Fordybelse i, hvad der stod
for ham som det Vigtigste i det Aandsliv, han var bleven
indviet i fra sin Barndom af, --- den store, af hans Folk
udsprungne Digtningens Verden. Den stod for ham som en
Lysverden, der ikke blot havde kunnet vejlede og nære
hans eget personlige Liv, men ogsaa, medens Verdenskrigen
og dens Følger kuede Landene, kunde varsle om, at Lysets
Magter ikke kunde overvindes i Sjælenes indre Liv.

Inden jeg gaar ind paa Bogens specielle Indhold, maa jeg
forklare dens Navn, som skyldes en for Troels Lund karakteristisk
Symbolik. Danmark har ingen Bjerge, men des
gladere er vi for Bakkerne, der muliggør vide Udsigter over
vort flade Land, hvor Lys og Skygge formedelst den megen
Blæst ofte skifter meget hurtig. Lige udenfor København
har vi Valby Bakke. Valby er forlængst en Forstad til København,
og Villaer, høje Huse og Fabriker bedækker det
% --- p. 130 ---
\setcounter{footnote}{0}%
\opage{130}forhen saa idylliske Sted, til hvilket saa mange Erindringer
fra det forrige Aarhundredes Literatur knytter sig. Oehlenschläger
levede som ung paa Frederiksberg Slot, hvor hans
Fader var Slotsforvalter. I Degneboligen ved Kirken skrev
han „Helge“, et af sine allerypperste Værker, og i det nære
Bakkehus, hos Rahbeks, havde han gjort Bekendtskab med
sin Brud, Kammas Søster. I sin Alderdom havde han Sommerbolig
paa Fasangaarden i Frederiksberg Have. I Bakkehuset
paa Bakkens Skraaning mod Øst boede Rahbeks, hvis
Hjem i lange Tider var et Midtpunkt for dansk Aandsliv.
Næsten alle navnkundige Digtere færdedes der, og efter
Ægteparrets Bortgang har baade Heiberg og Grundtvig haft
Sommerboliger der. At Troels Lunds Forældre boede der en
Række af Aar om Sommeren, bragte Stedet og dets Historie
saa levende frem for hans Fantasi, da han fordybede sig i
sine Livserindringer, at han istedetfor en Selvbiografi kom
til at give et stemningsfuldt Overblik over dansk Poesis
Historie i det nittende Aarhundrede.

Solbjerg var Navnet paa en Landsby, der i gamle Dage
laa paa Bakken, men maatte vige Pladsen, da Frederiksberg
Slot blev bygget. For Troels Lund betyder Solbjerg et
Bjerg eller en Bakke, hvor Solen dyrkedes. Og baade Bakkehus
og Solbjerg blev nu for ham symbolske Udtryk. Bakkehus
betyder et Hjem paa Vejen op mod den Lysdyrkelse,
der strækker sig gennem Menneskehedens Historie, og som
ifølge ham paa ejendommelig Maade arbejdede sig frem i
dansk Poesi i det 19. Aarhundrede, et lille Exempel paa den
Higen mod Lyset, som er Menneskeaandens inderste Væsen.
Troels Lund har holdt sig til dette Exempel, fordi det
var det, han bedst kendte. Andre Folk kendte han, siger
han, for lidet til for at kunne skrive om deres Hjemsteder
for Lysdyrkelse. Ikke blot her i Norden er der mange Solbjerge
og mange Bakkehuse, men ethvert for Lyset modtageligt
Sted er et lille Solbjerg.

% --- p. 131 ---
\setcounter{footnote}{0}%
\opage{131}\romanhead{II}
% CHECK roman numeral read

Efter Troels Lunds Opfattelse laa det i Forholdenes Natur
her i Danmark i Begyndelsen af det 19. Aarhundrede,
at Digtekunsten maatte blive „den kækkeste og kyndigste
Vejviser“ paa det aandelige Omraade. „Den gik i Spidsen og
udtalte ofte det betagende og forløsende Ord, medens Videnskaben
endnu tøvede, Kirken forkætrede.“ „Gennem
Digtningen er den aandelige Udvikling i Danmark lettest
foregaaet, medens Videnskab, Teologi og Filosofi fulgte
bagefter.“

I disse Ord er Værkets Grundtanke givet. Og det ligger i
de anførte Ord, at det Mørke, Lyset skal besejre, ikke blot
er det, der skyldes Ulykke, Smerte og Sorg, men ogsaa det,
der voldes af dogmatiske Overleveringers Magt til at spærre
Adgangen til det Lys, Menneskeaanden i sig har Kraft til at
naa. Værket giver saa et Galleri af Landets Digterskikkelser,
tildels med glimrende Karakteristik og med en Mængde
enkelte Træk. Troels Lund havde, baade som Literaturhistoriker
og som Kulturhistoriker, en stor Evne til at leve
sig ind i Fortidens Skikkelser og lade Fantasien udstyre
dem med en Rigdom af Træk. Angaaende Enkelthederne
spørger man --- deri maa man give hans ofte saa strenge
Kritikere Ret --- ofte, hvorfra han egentlig har dem, ---
og spørger ofte forgæves. Men hvad der er Hovedsagen for
os her, det er den Række af Lysseere, han vilde fremføre
for os. Det er umuligt at gengive Fremstillingen i det Enkelte;
vi maa lade os nøje med nogle faa Exempler.

Først er den danske Digterkonge, Adam Oehlenschläger,
at nævne. Hos ham kundgør sig allerede i Ungdomsdigtet
„Jesu Liv i den aarlige Natur“ et nyt Livssyn. Medens man
i Kirken ofte har oplyst sin Lære ved Billeder fra den ydre
Natur, saa bruger nu Digteren her omvendt Kirkelærens
Forestillinger som symbolske Begivenheder i Naturen. Han
appellerer til Læsernes religiøse Barnelærdom for at anskueliggøre
og levendegøre Naturens Liv for dem. Dette
betegner et Vendepunkt. Naturen begynder at komme til
% --- p. 132 ---
\setcounter{footnote}{0}%
\opage{132}sin Ret. Man følger dens Anvisning, og de kirkelige Forestillinger
benyttes kun som Støttepunkter for Fantasien.

Særlig dvæles der ved Kamma Rahbek, som af Ingemann
blev kaldt de danske Digteres Moder. Ikke blot fulgte hun
unge Digtere paa deres Vej med hjertelig Forstaaelse, men
hun gennemførte Poesien i sit Liv, var en Naturevangelist, i
stadig Udvikling og Stræben. I sit kærlige Sind og i det
Haab, der bygger paa Kærligheden, fandt hun Løsningen
af Livets Gaader. Hun var en religiøs Natur, men gik ikke
med til den Skærpelse af Ortodoxien, der kom i Romantikens
Følge.

Selve denne Ortodoxi, med den stærke Betoning af Djævle-
og Helvedeslære, blev overvunden gennem de store Lysseere.
Ligesom Danmarks største Jurist, A. S. Ørsted, i Tilslutning
til sin almindelige Opfattelse af, hvorledes Fortidens
Lovbestemmelser skulde anvendes paa nye Tiders
Forhold, hævdede, at kirkelige Bekendelsesskrifter, der
selv skyldte en bestemt Tidsalder deres Tilblivelse, maa fortolkes
efter de nye Tiders Trang og Erfaring, saaledes blev
endog Grundtvig, den nye Ortodoxis krasseste Profet, efterhaanden,
især under Indflydelse af sin hjertelige Deltagelse,
i det opvaagnende nationale Liv herhjemme, en stor Lysseer,
for hvem Djævlelæren vel stod som en historisk Forudsætning
for det nye Testamentes Lære, men ikke var en
for alle Tider gældende Læresætning. Derved blev han en
stor aandelig Befrier og Leder for saa mange i Folket, en
mægtig Kæmpe for sand Oplysning.

Men Højdepunktet af frit Livssyn naas, ifølge Troels
Lund, hos Ingemann, Hauch og Paludan-Müller.

Hvad Ingemann angaar, er især „Holgers Sang til Livet“
at nævne. Den er „en Livshymne, saa fuld og inderlig, som
den endnu ikke var hørt i Danmark“. I den fandt Troels
Lund det bedste Udtryk for sit eget lyse Livssyn, kulminerende
i det sidste Vers:

\begin{quote}
Men naar under Solen ej Sjælen, hvad den vil,\\
saa er der andre Sole og andre Stjerner til.\\
Og slukkes alle Sole og Stjerner engang,\\
dog Livets Kilde springer, hvor evig den sprang.
\end{quote}

% --- p. 133 ---
\setcounter{footnote}{0}%
\opage{133}Med Rette fremhæver Troels Lund den store uendelige
Horizont i Ingemanns Livsdigte. Medens Middelalderens
begrænsede Verdensbillede stadig beherskede Grundtvigs
Forestillinger, spores hos Ingemann, ligesom senere hos
Hauch og Paludan-Müller, den store Udvidelse af Verdensbilledet,
som var bleven Følgen af Kopernikanismen. Allerede
Holberg havde --- med megen Forsigtighed --- udtalt
sig om den Indflydelse, Kopernikanismen maatte faa paa
hele Verdensopfattelsen. Men jeg tror, at Troels Lund har
Ret i, at først hos Ingemann, Hauch og Paludan-Müller spores
denne Indflydelse i dansk Poesi og i den Religiøsitet,
som fandt sit Udtryk i dansk Poesi.

Hos Hauch fremtræder lignende Tanker paa inderligere
Maade og i Samklang med hans levende Natursans. Han
er Trofasthedens og Ærestrangens Digter. Paludan-Müllers
Tro er en Haabets Religion. For ham er Haabet ikke blot
Sjælens Øje, men Livskærnen i os. Digtet „Nattevagt“ (1837)
indeholder hans store Tro paa Lyset og Livet, trods alle
Illusioner og Begrænsninger, der møder os paa Vejen. Senere
har han i „Adam Homo“ beskrevet Menneskets Svaghed,
Selvbedrag og Usselhed, men ogsaa Kærlighedens Magt
til at hæve selv den usleste Sjæl opad mod Lyset: Alma, der
selv har Adgang til Salighedens Rige, følger dog trofast
Adam gennem Skærsilden.

Troels Lund's skønne Værk blev en Torso. Det afbrødes
midt i tredie Del ved hans Død, den Død, han gik imøde med
Frejdighed, som det sig en Lysseer bør. Denne tredie Del
skulde handle om Søren Kierkegaard, men naaede kun til at
berette om hans Ungdom. Fortsættelsen vilde have stillet
Troels Lund et stort Problem, thi Kierkegaard erklærede
ikke blot Romantikens Poesi for en blot Drøm, men hævdede
tillige, at den idylliske Opfattelse af Kristendommen,
som Romantikerne byggede paa, var et Bedrag. Jeg véd
ikke, hvorledes Troels Lund over for denne vældige Skikkelse
og denne kraftige Protest vilde have kunnet fastholde
og gennemføre sin Lære om det lyse Livssyns Sejr
i dansk Literatur. Den Skærpelse af Livsproblemerne, Kier%
% --- p. 134 ---
\setcounter{footnote}{0}%
\opage{134}kegaard bevirkede, krævede i hvert Tilfælde en alvorlig
Prøvelse. Maaske kan man af en Antydning i Slutningen
af andet Bind, formode, at Troels Lund væsentlig har set
Kierkegaards Betydning i hans Omstyrten af den synlige
Kirke. Saa kunde der blive større Plads for Digtningens
frie Flugt; men den maatte rigtignok saa ogsaa foregaa paa
andre Vilkaar end forhen.

\romanhead{III}
% CHECK roman numeral read

Kun nogle faa Ord maa jeg fra mit eget Standpunkt
føje til den Karakteristik, jeg har forsøgt at give af min
afdøde Vens Værk. Det er en stor Opgave, han havde sat
sig, at vise, hvorledes en høj Livspoesi udviklede sig hos
den danske Guldalders Digtere. Og med stor Taknemmelighed
erkender ogsaa jeg, hvor meget vi i denne Henseende
skylder disse Digtere. Men det var Romantikens
Tid. Man tog sig Realiteterne, baade i Natur og Historie,
let, og Poesien svævede højt over Virkelighedens Vande.
Det kan ikke siges, at den romantiske Digtning var Forgænger
for Videnskab og Kirkelære. Den gav i mange Henseender en Omtydning af de kirkelige Forestillinger, og
dens dristigste Flugt skyldtes Renaissancens udvidede Verdensbillede,
som først nu ret blev draget frem udenfor de
specielle Videnskabers Omraade. Og Digtning kan overhovedet
ikke i den Forstand være Forgænger for Videnskab,
at den skulde kunne forklare os noget, som Videnskaben
endnu ikke kan forklare. Den kan ikke udfylde
Videnskabens Huller, naar det gælder Forstaaelse eller Forklaring
af Virkeligheden, ligesaalidt som Videnskaben nogensinde
kan naa op i Højde med Digtningen, naar det
gælder at udtrykke den Stemning, den Sindsbevægelse, som
Livets store Modsætninger og Sammenstød kan vække.

Den Livspoesi, vi trænger til, og som enhver Tid trænger
til paa sin Maade, maa ikke, som Romantiken, svæve over
Virkeligheden, men have sin Rod i frie og dybe personlige
Erfaringer, udspringe af Sind, der er prøvede i Livets og
% --- p. 135 ---
\setcounter{footnote}{0}%
\opage{135}Tankernes Strid. Den maa ikke mene at kunne løse Problemer,
som Tanken er stanset ved, heller ikke maa den
være blot Omsætning i digterisk Form af overleverede Forestillinger,
men i dristigt Opsving skal den give store Symboler
paa den Livets Kamp, den Brydning af Lys og
Mørke, der vil vedvare, saa længe Livet varer. Lyset kan
sejre, men Sejren maa stedse tilkæmpes paany. Det var vel
egentlig ogsaa Troels Lunds Mening, skønt han ofte synes
at svælge i det rene Lys, hvor alle Modsætninger er fortonede.
Han har udtalt, at netop i de mørke Tider er det,
at Troen paa Lyset skal vise sig.

Jeg for mit Vedkommende er enig med ham i, at der i
hvert Tilfælde er to vidunderlige Exempler paa Livspoesi,
vi kan tage med os fra den danske Guldalder, --- Ingemann's
„Sang til Livet“ og Paludan-Müller's „Nattevagt“.
Den Mandighed, det første Digt udtrykker, og den dybe Enhedsfølelse,
det andet finder Ord for, er vigtige Elementer
i enhver kraftig og inderlig Følelse ved Livet. Har vi i Livet,
som vi kender det, og som vi har grundet over det med de
Tanker, der staar til vor Raadighed, fundet noget, som var
værd at arbejde og kæmpe for, og har vi sat vore bedste
Kræfter ind herpaa, saa behøver vi ingen ydre Støtter for
vort Haab. Vi kan da sige med en Digter fra en noget senere
Tid end den, fra hvilken Troels Lund hentede sine Exempler:

\begin{quote}
Kæmp for Alt, hvad du har kært,\\
dø, om saa det gælder!\\
Da er Livet ej saa svært,\\
Døden ikke heller!
\end{quote}
\piecerule

% --- p. 136 ---
\setcounter{footnote}{0}%

\opage{136}\essayhead{BOGEN OM DE VISES STEN}

Dette er en Bog om at blive voxen. Den begynder derfor
naturlig med Barnelivet. Og den Maade, hvorpaa Børnenes
Liv i den lille Landsby skildres, vil maaske for mange
Læsere være det Bedste i Bogen, det, hvortil de helst vender
tilbage. Ikke fordi Børnenes Liv skildres som idel
Glæde og Harmoni. Ogsaa dér finder vi baade Raahed og
Mildhed, og selv det mildeste Sind kan gribes af Rovdyrsvildhed,
naar Raaheden har forarget det. Men i Barnealderen
løses Vanskelighederne forholdsvis let. Horizonten er
begrænset, Opgaverne passer nogenlunde til Evnerne, og
Mangfoldigheden er ikke større, end at et modtageligt Sind
kan mene at se den som et samlet Hele. De, der ikke er
modtagelige paa denne Maade, mister derfor ikke Sammenhængen
i deres Sind. De møder hver Bølge, Livet sender
dem imøde, og svømmer videre paa dens Ryg. De spørger
ikke „Hvorfor?“ Deres Valgsprog er „Hvorfor ikke?“ ---
hvis de har noget. De følger den Overlevering, de er ført
ind i, og deres Tillid til den rokkes ikke, netop fordi de
tager Tingene, som de kommer, og ikke har nogen Trang
til at kaste et Blik ud over Livet idethele. Der skildres
saaledes i Bogen et ungt Par, Ellen Nielsen og Hans Olsen.
En Dag i en Skrivetime i Skolen kom de til at se op paa
hinanden og smilte derved, og siden den Dag har den ene
af dem kun let eller grædt, naar den Anden gjorde det.
En skønne Dag, efterat de var blevne voxne, kom de til at
tale sammen herom og blev enige om, at det altid vilde
gaa saaledes med dem. Saa vandrede de Haand i Haand
gennem Livet, i fælles Lykke og i fælles Arbejde. De tror
paa, hvad de har lært; hvorfor skulde de ikke det? Og
% --- p. 137 ---
\setcounter{footnote}{0}%
\opage{137}de er endda sikre paa, at selv om de ikke havde lært det,
havde de alligevel troet det: „det kommer over dem af sig
selv.“ For andre gaar det ikke saa let. Der er den milde og
hjælpsomme Holger Enke, som i Forbitrelse over, at den
unge Kvinde, hvis Renhed og Uskyld han beundrer, forføres
af en Vellystning, der tager Flugten, selv voldtager
og myrder hende, og saa efter et mislykket Selvmord villig
gaar i Tugthuset. Da han formedelst god Opførsel benaades,
vender han tilbage til Landsbyen, tager Folks Frygt
og Foragt som en Del af sin Straf, vinder tilsidst ny Agtelse
og forkynder i Bogens Slutning den Livsvisdom, han
er naaet til. Til ham vil vi siden vende tilbage.

Det er Degnesønnen og Præstesønnen, vi foreløbig maa
dvæle ved. Degnens Søn, Jens Dahl, stiller ved sin Personlighed
og sine Tanker Bogens egentlige Problem. Allerede
som Barn lærte han af egen Erfaring den store Forskel
mellem, hvad han kaldte, „det Aabne“ og „det Lukkede“.
For den, der er „aaben“, gælder Tids- og Stedsforskelle
ikke. Alt er nært. Hans lille Broder, der er død, er stadig
nærværende. Det er ikke muligt at udtrykke „Aabenheden“ i
Ord; det ene maatte jo saa følge paa det andet, hvilket er en
ligesaa stor Ufuldkommenhed, som at man maa tage Skridt
efter Skridt for at komme fra det ene Sted til det andet.
I „Himmelsproget“ siger vi Alt, uden et eneste Ord. Her
er alle Gaader løste --- fordi de slet ikke er blevne til. Men
naar man er „lukket“, maa man tænke sig om og gaa
Skridt for Skridt. Og saa kommer alle Vanskelighederne
og det besværlige Arbejde. Den syttenaarige Dreng følte
sig i Kraft af sin „Aabenhed“ at være et Geni, sikker paa,
at Alt, hvad der var „lukket“, vilde kunne gennemskues
og overvindes ud fra det „Aabne“. Og det advarede ham
ikke, at han faldt for den første Fristelse, han mødte. Men
da han senere kom tilbage til Landsbyen, havde han gjort
Erfaringer, der lod ham se paa sine gamle Kammerater
dels med Vemod, dels med Skinsyge, paa Grund af den
naive, uvilkaarlige Maade at tage Livet paa, de havde bevaret.
De gik fra Trin til Trin paa Livets Stige, og det ene
% --- p. 138 ---
\setcounter{footnote}{0}%
\opage{138}Trin fulgte saa af sig selv efter det andet. I deres Skikkelighed
og Ubevidsthed var de ikke langt fra „Himmelsprogets“
Verden. --- Før vi følger ham videre, maa vi fortælle
om hans Ven, Præstesønnen.

Christian Barnes havde tidlig mistet „Aabenheden“. Det
Indblik, han som Skoledreng i Købstaden havde faaet i,
hvad der kan skjule sig bag de Voxnes „Lukkethed“, havde
gjort stærkt Indtryk paa ham. Og endnu stærkere virkede
en hemmelig Last, han var kommen ind i. Hans Trang til
at analysere sig selv og Andre førte ham bort fra „Himmelsproget“.
Men skønt han havde forladt sin Barnetro, havde
han dog den Overbevisning, at der gaves en naturlig religiøs
Følelse, som ikke staar og falder med nogen bestemt Religionsform.
Og han holder herom et langt religionsfilosofisk
Foredrag for sin Ven Jens Dahl, da ogsaa denne er
kommen som Student til København. (Der forekommer sikkert
for mange Forelæsninger i Bogen. Men jeg overlader
de literære Kritikere at paatale det, hvis det skal paatales).
Hvad Christian Barnes søger (alt imedens han forbereder
sig til Magisterkonferens i Engelsk), er at blive karsk, at
faa det sunde Blik paa Livet, som fører til fri religiøs Følelse.
Det lykkes ham ogsaa under Indflydelse af en energisk
engelsk Clairvoyante, der sender ham til Amerika, for
at han som Cowboy kan faa at vide, hvad fysisk Træthed
er, og derefter erhverve sig et Overskud af Livskraft, der
kan benyttes til sund Tænken. Tillige lærer han hos hende
at samle sit Sind i Stilhed og Andagt. Derimod finder han
i Amerika ikke de Vises Sten, ja, han mener endog, at
den er ikke værd at have, efter at der med høj Visdom var
sørget for, at den var godt gemt. Da han er vendt hjem
fra Amerika, møder han igen den Kvinde, der som lille
Pige talte med ham over Havegærdet, naar han om Torsdagen
havde en Mellemtime. Hun havde ikke glemt ham.

Før han rejste til Amerika, havde han givet sin Ven
Jens Dahl det skæbnesvangre Raad at studere Mystik og
Teosofi. Han sagde til ham: „Du, som vil ind i det Aabne
og leve i Himmelsprogets Verden, der er skjult for almin%
% --- p. 139 ---
\setcounter{footnote}{0}%
\opage{139}delige Dødelige, er du ikke, ganske som disse [Mystikere
og Teosofer], En, der haaber at finde de Vises Sten?“ ---
Jens Dahl var forberedt til dette Studium gennem sine
barnlige Forsøg paa at forstaa „det Lukkede“ --- d. v. s.
det Mangfoldige, de overraskende Erfaringer, det Haarde
og Vanskelige i Livet --- ved „det Aabne“ --- d. v. s. det
uvilkaarlig opdukkende Helhedsbillede, som kan danne sig
i et barnligt Sind. Nu opgav han alle andre Studier for at
fordybe sig i de gamle Mystikere. Og han følte sig snart
hjemme i deres Verden. Det var, som om han her mødte
sit Eget igen, --- som var han kaldt tilbage til Barndommens
Dage.

Det er maaske ikke unødvendigt at forklare, hvad de
Vises Sten historisk betyder. Udtrykket stammer fra Alkymisterne,
der søgte efter et Stof, som kunde forvandles til
alle andre Stoffer, fremfor Alt til de kosteligste, Sølv og
Guld. Og dette vidunderlige Stof, troede man saa, kunde
tillige helbrede for Sygdom og forlænge Livet. Alle materielle
Goder kunde opnaas ved det. Alkymisterne var tillige
ofte Mystikere, og omvendt. Men de Mystikere, der er
Tale om i den foreliggende Bog, søgte ikke et Stof, men en
Tanke, Tanken om, hvad der ligger til Grund for Alt og
holder Verden sammen, --- om Noget, som Folkereligionerne
har søgt at udtrykke i deres Billedsprog, men som
man kun ret fatter gennem energisk aandelig Koncentration,
især naar man kan blive Et dermed og se Alt i Lyset
af det.

Jens Dahl mente, som sagt, her at genfinde, hvad han i
sin Barndoms „aabne“ Stunder havde oplevet. Dog var de
Oplevelser, Mystikerne talte om, af en højere og vanskeligere
Art, og først da han havde faaet en Anvisning til
daglige aandelige Øvelser i Retning af Samling og Selvfordybelse,
kom en Dag Ekstasen over ham. Efter en streng
% --- p. 140 ---
\setcounter{footnote}{0}%
\opage{140}Selvprøvelse, der lod ham gennemskue sin hele Usselhed,
saa han tilsidst sin Personlighed, som han hidtil kendte
den, synke hen i Intet, men følte sig til Gengæld som baaren
af en kærlig Magt, den samme, der bar den hele Verden,
og ikke blot baaren af den, men Et med den, uden
at nogen Tanke, noget Ord kunde udtrykke det Lys, der
gennemstrømmede hans hele Væsen. Han mente nu at have
erfaret, hvad hans Ven saa ofte havde talt om, --- den
religiøse Følelse, der ligger dybt i Menneskets Natur og i
sig selv er uafhængig af alle Dogmer og alle ydre Former.
Gennem sin stærke Koncentration fik han ogsaa Evne
til at paavirke Andre gennem Hypnotisering og Suggestion.

Dahl var overtydet om, at hvad han nu oplevede, var
noget rent Naturligt, betinget ved bestemte Aarsager. Det
maatte derfor, mente han, være muligt at fremkalde den
ekstatiske Tilstand igen. Men kun ufuldkomne og flygtige
Glimt dukkede op hos ham. De stadige krampagtige Øvelser
gjorde ham nervøs; den forcerede Anspændelse voldte
% PRINTED AS IS (read at the image, odd spots o1)
kun uhyggelige Halucinationer, og samtidig vaktes sanselige
Drifter, ligesom det allerede havde været Tilfældet i
hans Overgangsaar, netop naar han havde været mest
„aaben“. I denne Tilstand imponeredes han af den Ro og
Sikkerhed, hvormed hans teosofiske Venner gennemførte
deres Metoder. De gennemgik en Slags Forskole og kunde
derefter faa umiddelbar Undervisning af Mestre fra selve
Aandeverdenen, Udsendinge fra de store indiske Ledere.
Dahl gennemgik da ogsaa en ny Træning i ekstatiske
Drømme, i hvilke han færdedes paa „Astralplanet“ mellem
de Afdøde, hvis Skærsildstid han blev Vidne til. Snart kom
det dertil, at han ikke kendte Forskel paa Drøm og Vaagen.
Han fik Hallucinationer paa Gaden og maatte befordres
hjem af Politiet. I et Besindelsens Øjeblik indsaa han,
at han havde misbrugt de „aabne“ Tilstande, der var
blevne ham forundte, idet han kun havde søgt Nydelse i
dem --- ligesom han havde behandlet de Kvinder, han
havde mødt paa sin Vej, som Urter, af hvilke han havde
uddraget kostelig Elixir, medens de hver for sig Intet havde
% --- p. 141 ---
\setcounter{footnote}{0}%
\opage{141}været for ham. Forgæves havde den engelske Clairvoyante
formanet ham til ikke at stræbe saa meget efter at se Himmelens
Engle, men hellere se til at de Engle, han traf her
paa Jorden, fik deres fulde og rette Betydning i hans Liv.
Han havde nu engang, som han selv tilsidst indsaa, „gjort
Guds Have til en Lystgaard, og det var ham derfor forment
at arbejde i den“. --- En Dag, han sad ved Bredden af et
Vand, saa han paa den modsatte Bred en kvindelig Skikkelse
vinke. Han vilde gaa lige over til hende, men Vandet
lukkede sig over ham.

Spørger man, hvilken menneskelig Interesse, Skildringen
af Jens Dahls Udvikling og Undergang har, maa den utvivlsomt
søges i, at der i et enkelt Exempel fremstilles de Farer,
der truer ved Overgangen fra Barndommens og den
første Ungdoms Drømme til Livets Virkelighed og Livets
Krav. Der er, som vi har hørt, dem, der paa en for dem
selv umærkelig Maade, glider over i Livets Virkelighed.
% PRINTED AS IS (read at the image, odd spots o1)
Drøm og Maal virker hos dem som mildt besjælende Krafte,
og det kommer ikke til Krise, end sige til Undergang. Men
der er Naturer, for hvem Barndommens „Aabenhed“, dens
naive Tro paa Letheden i at løse Livets Gaader fører til
haarde og skæbnesvangre Sammenstød. De ser ikke, at
det kun er gennem alvorligt Arbejde, et Arbejde, der maa
vælges i Overensstemmelse med personlig Trang og Evne,
at Omsætningen fra de drømte Idealer til Virkelighed kan
ske. Og det er da ikke blot den foreliggende Virkelighed,
der skal bøje sig for det drømte Ideal; det er nok saa ofte
dette Ideal, der skal antage nye Former efter de nye Erfaringer,
som paanøder sig. Her føres en indre Kamp, der
ikke altid mærkes af Andre, men som kan være af stor Betydning
og præge det inderste Liv. Paa alle Omraader kan
der blive Anledning til en saadan Kamp, --- i Videnskab,
Kunst, Religion, i Forholdet til enkelte Mennesker og til
Samfundet, og ikke mindst paa det Omraade, der til alle
% --- p. 142 ---
\setcounter{footnote}{0}%
\opage{142}Tider har været et Yndlingstema for Digtere, det erotiske
Forhold mellem Mand og Kvinde. Bogen om de Vises Sten
antyder da ogsaa en Parallel mellem Jens Dahls Forhold
til Livets Virkelighed og hans særlige Forhold til de Kvinder,
Livet fører ham imøde. Det kunde her være undersøgt,
om der ikke var en direkte Sammenhæng mellem
Mystik og Erotik. Hos Jens Dahl optræder Erotikken snarest
som en Modsætning til Mystikken, en Art Kontrastvirkning.

Men Læseren kan ikke undgaa at spørge sig selv, hvorfor
han skal føres ind i al denne Mystik og Teosofi, der
gør det nødvendigt for Forfatteren at lade En af sine Personer
holde Foredrag, for at Læseren kan være à jour med,
hvad der skal fremstilles. Frembyder Livet ikke Forhold
nok, som kan benyttes til at belyse den betydningsfulde
Overgang, der er Bogens egentlige Tema? Det er muligt
(derom kan jeg ikke dømme), at Mystik, Yogisme og Teosofi
spiller en saa stor Rolle i Nutidens Menneskeliv, at det
kunde ligge nær for en Forfatter at vælge sit Emne her. Og
det kan ikke nægtes, at Forfatteren giver karakteristiske
Billeder i de Mystikere og Teosofer, han bringer den unge
Mand i Berøring med. Det er levende Skikkelser hver for
sig, og de tager hver for sig deres Religion paa individuel
Maade. Det er ovenikøbet en af Teosoferne, der udtaler det
mest forløsende Ord om de Vises Sten. Det er en stakkels
Haandværker, der har ladet sin Kone og sine Børn i Stikken
og har kombineret sin Teosofi med Stræben efter at
frembringe et perpetuum mobile. Han kommer tilsidst til
at indse Umuligheden af begge Dele, og han vender tilbage
til sit Haandværk, hvor Arbejdet kan og maa gaa Stykke
for Stykke, Skridt for Skridt. Og naar han spørges, om han
ikke tror paa de Vises Sten, svarer han: den er en Milepæl
paa Vejen, --- og spørges der, hvilken Milepæl, er Svaret:
den næste paa Vejen og altid den næste! Det er Voltaire's
gamle Slagord om at dyrke sin Have, kun at den
brave Teosof har indset, at Haven ikke altid ligger ved
samme Milepæl. 

% --- p. 143 ---
\setcounter{footnote}{0}%
\opage{143}Bogens Betydning vilde i mine Øjne have været større,
hvis den havde bevæget sig mere i Livets Dagslys, hvor
tilsidst alle virkelige Opgaver ligger, selv om vi kan behøve
Nattens Grublen og Nattens Stilhed for ret at løse dem. Og
selv om man netop vil skildre, hvorledes man kan drømme
sig bort fra det virkelige Liv, behøver man ikke at hente
sine Exempler fra Teosofien. Der er desværre Exempler
nok endda. ---

Hvad Skildringen af de mystiske Fænomener angaar, er
der et Par Bemærkninger, jeg ikke kan undlade at gøre.
Det burde i Fremstillingen have spillet en større Rolle, at
Ekstasen, selv om der er gjort forberedende Øvelser, dog
stedse kommer uvilkaarlig, som en Overraskelse, og at det
er haabløst og selvmodsigende at bygge sit Liv paa Opnaaeligheden
af den Art Tilstande. Et nøjere Studium af
Mystikerbiografier fører desuden til den Formodning, at
ekstatiske Tilstande er knyttede til visse Perioder i Livet
og falder bort efter deres Slutning. Hos de berømte kirkelige
Mystikere træder dette ikke saa stærkt frem, fordi de
med stor Lethed glider over fra Oplevelse til Tydning. Og
da de henter deres Tydning fra de kirkelige Forestillinger
som selvfølgelige, lægger de ikke altid stærk Vægt paa
Forskellen mellem de egentlige Ekstaser og de livfulde og
betagende Stemninger, i hvilke de kan komme, naar den
kirkelige Overleverings mægtige Billeder stærkt griber dem.
Anderledes vil \textit{den} være stillet, der uden saadan Støtte vil
leve sig ind i Mystikkens Verden. For ham vil Modsætningen
være uudholdelig mellem, hvad man kunde kalde de store
og de smaa Øjeblikke, og den vil kun kunne overvindes,
naar der sættes et Arbejde ind paa i Erfaringens Dal at
gennemføre, hvad der er skuet oppe paa Højderne. Men
det er netop et saadant Arbejde, som moderne Mystikere
og Teosofer skyr. Mystikken har bragt store og gode Ting
med sig, navnlig ved at føre bort fra blind Optagethed af
Dogmer. Den har ofte netop ført fra Dogmet til Livet. Men
den kan selv blive en Fare, og den er det maaske for
Mange i vore Dage. Saa meget vi end maa haabe, at Kon%
% --- p. 144 ---
\setcounter{footnote}{0}%
\opage{144}centration, Begejstring og Selvfordybelse aldrig vil ophøre
at øve deres befrugtende Indflydelse paa det menneskelige
Aandsliv, saa er det ikke altid heldigt, at de antager Ekstasens
Form, og de vil vel ogsaa sjeldnere og sjeldnere
gøre det.

Der er en Person i Bogen, hvis Barndomshistorie ikke
meddeles. Det er „Kandidaten“. Lejlighedsvis faar vi at
vide, at han har foretaget Rejser rundt om i Verden. Hans
Eventyr paa denne Rejse skal ikke altid have været uskyldige;
men han kan dog sige, at han intetsteds har gjort
noget, der stred mod vedkommende Lands Love. Nu har
han slaaet sig ned i den Landsby, hvorfra han stammer, og
har købt sin afdøde Faders Hus. Han finder Egnen at være
Paradisets Have, naar man blot ser paa den med de rette
Øjne. Hver Enkelt maa skaffe sig disse Øjne ad sin egen
Vej. Med dette Blik paa Livet vandrer Kandidaten saa
rundt i Egnen og hjælper Store og Smaa paa Gled, naar
det kniber for dem. En barsk Missionær gør han latterlig
i Befolkningens Øjne. Et Par Ægtefolk hjælper han paa
en snedig Maade ud over den „Gemytternes Uoverensstemmelse“,
der truer deres Samliv. En ung Kvinde, der, overvældet
af Livets Ondskab, vil tage Livet af sig, redder han
i sidste Øjeblik. Han styrker Folk i deres religiøse Opfattelse,
naar han mærker, at den virkelig stemmer med
deres Tarv og Trang. Præsten i Sognet, som iøvrig er ham
meget velsindet, siger om ham, at han altid har den
Religion, som den, han netop talte med, har Trang til
at have. Det er den sokratiske Kunst, han øver og trods
den Medfølelse, der rører sig hos ham, er det ham en Nydelse
at øve denne Kunst, og han kan „grine“, naar den
er lykkedes godt. Denne Landsbyens Sokrates er et i høj
Grad oplivende, og ofte tillige rørende Element i Fortællingen.
Man glæder sig altid, naar han tager Affære. Husker
jeg ret, forekommer han allerede i en af Forfatterens tidligere
Fortællinger. 

% --- p. 145 ---
\setcounter{footnote}{0}%
\opage{145}Naar Kandidaten ikke altid tog sig sin Sjælesørgervirksomhed
højtidelig, saa var der dog én Skikkelse, overfor
hvilken han trods al sin Kunst og al sin Menneskekundskab
kom til at føle sig ringe. Det var Holger Enke, Voldtægtsmanden
og Morderen, ham, der efter at være kommen
ud af Tugthuset gik sin Pilgrimsgang i sin Hjemstavn,
hvor hans Daad var øvet, og langsomt arbejdede sig frem
trods Mistanke og Ringeagt. Hvad han maatte lide, var ham
en Lindring, og i ekstatisk Betagethed erklærede han endog,
at han villig gik ind under evig Fordømmelse, hvis det
skulde være saa. Tilgivelse søgte han ikke. Kandidaten
drog heraf den Slutning, at Holger havde overvundet Helvede,
var kommen ud over Modsætningen mellem Synd og
Naade og hvilede i den Forvisning, at Alt hvad han havde
været og øvet, var Genstand for Guddommens evige Forstaaelse.
Den 139. Psalme gav ham hans Livsfilosofi: Fra
hans Tilblivelsesstund var Alt skrevet i Guds Bog. Dagene
var bestemte, før de kom, ogsaa Alt, hvad der skete i Dagene.
Denne evige Tanke, Kilden til de store Love for Tingenes
Gang i Verden, bøjede han sig for, og hvad den lærte
ham, anvendte han ikke blot paa sig selv, men ogsaa paa
Andre. Da han saa mødte den Mand, som i sin Tid havde
forført den unge Pige, han beundrede og elskede, og som
derved var bleven Aarsagen til hans Forbrydelse og Ulykke,
lod han ham rolig gaa sin Vej med de Ord: „Jeg har ikke
skabt Dig.“

I Bogens sidste Kapitel udvikler Holger sin hele Pilgrimsfilosofi
for Kandidaten: Selv om man kom op i Himmelen,
vilde man intetsteds finde Gud, netop fordi han er overalt
og gennemtrænger Alt. Naar Mennesket stanser paa sin Vej
og mener selv at være Noget, gør han derved Gud til Intet,
og enhver Tanke, vi gør os om Gud, fjerner os fra ham.

Utvivlsomt hører vi her snarere Forfatteren end Holger
Enke. Det er Forfatteren, der i Spinoza's Aand omsætter
den 139. Psalme til et spekulativt System. Den virkelige
Holger vilde være stanset forinden og have holdt sig til de
store Billeder, hans Religion fremførte for ham. Maaske
% --- p. 146 ---
\setcounter{footnote}{0}%
\opage{146}vilde han, uden at være sig det bevidst, have tydet sine
Billeder paa sin egen Maade, paa Grundlag af hvad han
havde oplevet, forbrudt, lidt og forstaaet. Men abstrakte
Begreber som Væren og Ikkeværen havde han neppe tumlet
med. Og Landsbyens Sokrates vilde være bleven greben af
hans Sinds Dybde og hans Aands Kraft uden det.

Jeg har stillet mig den Opgave ved denne Anmeldelse at
undersøge, om det er lykkedes Forfatteren at omsætte den
Forstaaelse af menneskeligt Aandsliv, særlig af dets religiøse
Former, han utvivlsomt besidder, i virkelig Poesi.
Jeg tror ikke, det overalt er lykkedes ham. Min Opgave
har gjort, at jeg har holdt mig til Hovedpersonerne og
deres Tanker. Men rundt om dem staar en Række Skikkelser,
hver med sin Livsbane og sin Maade at tage Livet
paa. Og her har jeg fundet ægte Poesi. Naar jeg igen tager
Bogen frem, hvad jeg utvivlsomt vil, vil jeg først og fremmest
opsøge disse Bipersoner, som optræder dels i Landsbyen,
dels i Købstaden, dels i København. I Skildringen af
dem lægger der sig en Menneskekundskab og en Menneskekærlighed
for Dagen, som for mig betinger den Stemning,
Bogen efterlader. Det laa kun udenfor min Opgave,
desværre ogsaa udenfor min Evne, at karakterisere disse
Skikkelser.

\piecerule

% --- p. 147 ---
\setcounter{footnote}{0}%

\opage{147}\essayhead{ANKER LARSEN: „FOR AABEN DØR“}

Opmærksomme Læsere af Anker Larsens tidligere Bøger
vil denne Bog, saa meget Nyt og Frisk den end bringer,
ikke volde nogen Overraskelse. Det Emne, som i hans tidligere
Bøger i manges Øjne blev behandlet med størst Kærlighed
og Dygtighed, lykkelige Børns Liv paa Landet, spiller
ogsaa i den nye Bog en afgørende Rolle, skønt det optræder
og virker mere umiddelbart. Barnets fulde Hengivelse
i sine Tilstande og Stemninger, dets inderlige Opgaaen
i, hvad enkelte Øjeblikke kan give, dets Uberørthed
af Kritik og Refleksion gør det for Forfatteren til et Ideal
for de Voxne, hvis Sind ofte tynges eller indsnevres af
Vanskeligheder og Problemer, som ikke altid har deres
gode Grund. Vejen til Helvede er i vore Dage ofte brolagt
med ubegrundede Bekymringer.

Allerede i Forfatterens maaske mest læste Bog, „De Vises
Sten“, og ikke mindre i „Martha og Maria“, stod Barnesindet
ikke blot som digterisk Emne, men blev fremholdt
som Udtryk for en Tilstand, i hvilken vi kunde aabne os
helt og føle os helt fri. Han tager allerede der det gamle Ord
op, at vi skal blive Børn paany. Grundtanken i „De vises
Sten“ er given i Modsætningen mellem det Aabne og det
Lukkede. For den, der er „aaben“, gælder Tids- og Stedforskelle
ikke. Han lever helt i et værdifuldt Nu, og alt samler
sig for ham til en Helhed. I de „lukkede“ Tilstande er man
derimod betagen af Afstande, af Hindringer, af en Mangfoldighed,
der ikke kan samles til Harmoni.

Allerede ved sin Titel minder den nye Bog om, hvad der
var et Hovedemne i de foregaaende. Forskellen er, at dette
Emne nu tages mere umiddelbart op, tildels i selvbiografisk
% --- p. 148 ---
\setcounter{footnote}{0}%
\opage{148}Form. Man følger Forfatterens Fremstilling med stor Interesse,
og den er af ikke ringe psykologisk Betydning. Anker
Larsen beretter, hvorledes der paa et vist Tidspunkt i hans
Liv, efterat han havde prøvet baade Teologi og Teosofi, og
efter at han var bleven hildet i, hvad han kalder „Dagbladskultur“,
fødtes en Trang hos ham til at gennemleve paany,
hvad der havde gjort hans Barndom saa lys, saa fredfyldt
og saa helstøbt. Mere og mere fyldte Længselen efter den
svundne Tid hans Sind. „Det første Skridt paa Vejen hjem“
kom han til at gøre en Dag under en Vandring i en Skov.
Han følte sig paa én Gang fuld af Sundhed og Kraft og berust
af Skovstemningen. I denne paradisiske Stemning dukkede
Billedet af en yndig Skovsti op i hans Sjæl. Det fremgaar
af hans Beretning, at dette Billede først dukkede op,
efter at Stemningen havde udfoldet sig. Glæden ved Billedet
forøgede da selve den Stemning, hvis Symbol det blev. Han
havde været mæt af at se paa, hvad der omgav ham og havde
taget Øjnene til sig, og først da han igen saa op, saa han
for sig „en lille Skovsti saa frisk, ren og eventyrlig, at den
maatte være en af Gangene i Paradisets Have“. Først senere
blev han klar over, at Stien hørte hjemme i hans Barndomsegn.
Der var egentlig ingen særlig Lighed mellem den
Sti, han gik paa, og den, der nu stod for ham. --- Jeg tror,
at hans Beskrivelse maa være nøjagtig. Det har været Stemningen,
der aldeles overvejende blev Aarsag til Billedets Opdukken,
og Ligheden mellem de to Stier har i hvert Tilfælde
været af underordnet Betydning. Derimod har der
været Lighed i hans Stemning nu og hans Stemninger fra
hans Barndomstid, dog saaledes, at hin ikke naaede disse i
fuldt Liv, førend det Syn, der havde været knyttet til en
Barndomsoplevelse, uvilkaarlig dukkede op. Den Trang, der
i enhver stærk Stemning rører sig efter Udtryk, en Slags
aandelig Mimik, førte til en Art Vision, idet Synet af det
for Øjeblikket Nærværende ikke tilstrækkelig fyldestgjorde
den.

De ældre Mystikere skelner skarpt mellem Oplevelserne
selv og de Billeder, i hvilke disse Oplevelser finder Udtryk.
% --- p. 149 ---
\setcounter{footnote}{0}%
\opage{149}Man kan fra en Lighed mellem de Billeder, der dukker op
hos forskellige Mystikere, slutte tilbage til en Lighed i deres
Oplevelser. I min Bog „Oplevelse og Tydning“ har jeg særlig
fremdraget dette i den Sammenligning, jeg foretager
mellem to kvindelige Mystikere, den spanske Nonne Teresa
i det 16. Aarhundrede og den mærkelige protestantiske Mystiker
i det 20. Aarhundrede, som af hendes psykologiske
Skriftefader kaldes Cécile V. Der er Analogi mellem Teresa's
Billede af Jernet, der gløder i Baalet, men ellers ligger som
en sort Masse, og Cécile's Billede af Planken i Havvandet,
der snart straaler belyst af Aftensolen, snart blot ses som
et sort Bræt. De to Mystikere har begge følt Livets og Sindets
Modsætninger saa stærkt, at disse Modsætninger uvilkaarlig
fandt symbolsk Udtryk i deres Visioner. Paa dette
Punkt er der en karakteristisk Forskel mellem Anker Larsen's visionære
Oplevelser og hine Mystikeres, en Modsætning,
jeg senere skal vende tilbage til. Det er naturligvis ogsaa
hos ham en skarpt følt Modsætning, der afføder Trangen
og Stemningen. Men det har, efter hans Beretning,
netop været Higen efter at komme ud over Splittelse, efter
Samling og efter helt og frit Liv, der har ført ham over i
de Tilstande, hans Visioner symboliserer, og han har ikke
følt Trang til at lade selve de Modsætninger, han stræber
ud over, finde Udtryk i de visionære Symboler. Han higer
efter Barndommens naive, udelte Liv, og han finder Tilfredsstillelse
i Dvælen ved et Billede, der faar ham til at
glemme alt Andet, ogsaa Striden og Deltheden. „Virkelighedens
dybe Lykke“, siger han, „opfylder mig; min Tilstand
bredte sig og blev Et med alle ensartede Tilstande“.
--- Og saa vandrer han fri og glad henad sin genfødte
Skovsti.

Forfatteren lægger stærk Vægt paa, at hvad han saa for
sig, ikke var Erindringer, der kunde fremkaldes efter Ønske.
Glimtene kom af sig selv, og selve Tilstanden havde
Præg af en Nærværelse, ikke af Fantasi. Tilstanden fik hos
ham sit Præg af at være fremkaldt ved Længsel efter Barndommens
hellige Enfold, og Glimtene virkede som Manel%
% --- p. 150 ---
\setcounter{footnote}{0}%
\opage{150}ser om at begynde Livet forfra. De havde som han udtrykker
sig, en villiesbestemmende Magt, satte ham et Maal,
der ikke kan ændres. Dette Maal er at tage Livet op, saaledes
som det nu engang foreligger, --- at ville sin Skæbne.
Møder man Modstand, saa betragtes den som et Middel, et
Redskab til at overvinde sig selv. Den Frihed, der brød igennem
ved det enkelte Glimt, bliver først virkelig ved at fastholdes
i Tidens Løb.

Efterhaanden fik Glimtene en varigere Karakter og førte
til et Klarsyn, der hævede over Tid og Rum og aabnede
Vejen for Evighed. „Det evige Liv“, siger Anker Larsen,
„er ikke et Liv paa den anden Side Døden. Det evige Liv
er en Tilstand i Menneskets Sjæl; det er Livets dybeste
Virkelighed, som folder sig ud i hele hans Væsen og for
hans undrende Øjne. Ak, alle Ord er ubrugelige til dets
Beskrivelse.“ Det syntes, som om ogsaa Tingene i Verden
nu aabnede sig for ham og viste ham deres hele sande Virkelighed.
Den gamle „lukkede Skikkelse“ var kun som en
Forklædning, ikke deres egentlige Virkelighed. Til den
„lukkede“ Skikkelse hører især Modsætningen mellem Liv
og Død. Den falder bort for det aabne Blik. Forfatteren
minder om, at han i en af sine Bøger lader en Dreng sidde
under et Hyldetræ og tale „Himmelsprog“ med sin afdøde
Broder. Den Dreng har befundet sig „i en indre Tilstand,
hvor alle hans Oplevelser var nærværende“: „Der gives
ikke Afstand; alting er \textit{her.} Der gives ikke Fortid og Fremtid,
alting er \textit{nu}.“ „Døden er slet ikke til.“ „Den, der har
oplevet Nuet, ser intet bundløst Svælg mellem „Liv“ og
„Død“.“ „Jeg er ikke interesseret i Langvarighed, men i
Evighed.“

Den Tanke, Anker Larsen her anslaar, minder om Shelley's
„Sensitive Plant“, hvor Digteren hævder, at Skønhed
og Kærlighed er hævede over Modsætninger mellem Liv og
Død. Den for faa Aar siden afdøde dybsindige engelske
Tænker Herbert Bradley udtalte sin Tilslutning til Shelley's
skønne Digt, og han føjer til, hvad der vil tiltale vor Forfatter,
at det Paradis, i hvilket hine Ting gror, ikke er
% --- p. 151 ---
\setcounter{footnote}{0}%
\opage{151}nogen Egn eller Sted i Tid og Rum, men at det er \textit{her,} ---
at det er overalt, hvor et endeligt Væsen er hævet op til det
højere Liv, som nu er vaagen Virkelighed. ---

Den Aabenhed, der vandtes, alt eftersom varig Lysning
afløste de enkelte Glimt, affødte ikke blot Visioner, men
ogsaa Trang til Udraab, til Ord, --- saa klar Forfatteren
ellers var over Ordenes Utilstrækkelighed overfor det Bedste
og Højeste. Han fandt Luft i Udraabet „Gud!“ Dette
Ord indeholdt for ham „Jubel, Væren, Nærværelse, personligt
Forhold“. Hvad der skulde udtrykkes, var, at „der kom
noget Mægtigt og Godt til ham“. Det var som et dybt sjæleligt
Aandedrag, --- som om Varme, Styrke og Tryghed
strømmede ind i Sjælen. Men end ikke Gudernes Sprog vilde
kunne udtrykke saadan Oplevelse, og Forfatteren vil ikke
drøfte de forskellige teologiske og spekulative Gudsbegreber.
„Her er,“ siger han, „Tale om et praktisk og faktisk
Liv, og ikke om Livsanskuelser.“ Hvad der er oplevet, afføder,
som vi har hørt, Trang og Mod til at lade det Evige
(det Nu, som \textit{er}) gennemtrænge Livet i Tiden (det Nu, som
\textit{sker),} klare og lutre de begrænsede Formaal, saa at Arbejdet
staar som et Gode, ikke blot øves af Hensyn til Nytten.
Selve den saaledes fremkaldte Livsfølelse bliver Andagt,
uden at særegne Andagtstimer behøves. ---

Anker Larsen føler sig beslægtet med Mystiken, saaledes
som den optræder indenfor de forskellige Religioner. Han
henviser ofte til de Skrifter, gennem hvilke vi lærer den at
kende, som til sine „Familiedokumenter“. Dog er hans Mystik
af en anden Art end den, der fremtræder hos dem, han
nævner som sine aandelige Slægtninge. Den er en Art Friluftsreligion,
og „den aabne Dør“ kan ofte minde om „den
grønne Mark“, hvor en af vore fineste og djærveste Lyrikere
fandt Udtryk for sine Tanker om Livet. Og alligevel
synes det, som han vil drage sig ud af Livet og dets Kamp;
han bliver for let færdig med Livet i Forandringens Verden
ved at tale om „Dagbladskultur“ og „Døgninteresser“.
Og det er intet Tilfælde, at de Billeder, hans Stemninger
faar Luft i, er hvilende Billeder. „Skovstien“ ligger stille,
% --- p. 152 ---
\setcounter{footnote}{0}%
\opage{152}og der er intet Arbejde at gøre der. Det hænger, som allerede
bemærket, sammen med den Længsel efter Barndomsidyllen,
der var et Hovedmotiv til hans Mystik.

Og dog viste det sig, at hvad han oplever i de højeste
Øjeblikke ikke er ham nok. Det evige Nu er jo for ham
ikke fuldt ud virkeligt; det trænger til at „virkeliggøres“ i
Tiden. Det er altsaa (hvad iøvrigt de fleste Mystikere glemmer)
ikke blot det Timelige, der trænger til det Evige; ogsaa
det Omvendte gælder. Og det er en Hovedvanskelighed
for al Mystik (og for al Spekulation), at selv om der er fundet
en Vej fra Tid til Evighed, er der ingen funden fra Evigheden
tilbage til Tiden. I de store, samlende Øjeblikke er
det, som om Livet og alle Ting aabner sig. Men hvorfor er
de da ellers lukkede, --- og hvorfor er vort eget Sind saa
ofte lukket? Hvorledes kan den lukkede Tilstand bringes
til at ophøre? ---

Anker Larsen vil give os Oplevelser, ikke Livsanskuelser,
altsaa ingen Tydning. Han har for stor Foragt for Livsanskuelser,
der i hans Øjne kun er Frugter af ørkesløs Spekulation.
Beskrivelse og Tydning kan dog være vanskelige
at skille ad. Han bruger selv Billeder og Tanker, der stiller
Problemer, som man i Længden ikke kan overse. Men det
maa siges til hans Ære, at hvad vi hører fra ham gennem
den aabne Dør, gennemgaaende bærer Ægthedens og Selvoplevelsens
Præg. Med intellektuel Redelighed har han sørget
for, at hverken overleverede Meninger eller ufordøjede
Reflexioner kommer til at hæmme Oplevelsernes Oprindelighed.
Og han holder al forhastet Tydning borte, saavidt
han formaar. Han har her et Fortrin for Ludvig Feilberg,
der som sidste Tydning af sine Iagttagelser havde en hel
teologisk Fysik, som faa Læsere vist er blevne ret opmærksomme
paa.

Efter sin egen Erklæring er Anker Larsen ikke paavirket
af Feilberg. Dog er der paa mange Punkter et Slægtskab.
Feilberg's „Mulighedsværdier“ henstaar ligesaa sondrede
fra Livets store Arbejde som Anker Larsens „evige Nuer“.
Fælles for begge er dog fremfor Alt det friske, uhildede
% --- p. 153 ---
\setcounter{footnote}{0}%
\opage{153}Blik og Evnen til at give os umiddelbare Oplevelser. Her
kan det dog ikke nægtes, at Feilberg er Mesteren. Han var
en psykologisk Botaniker uden Lige, uovertræffelig i at
gribe og gengive de Glimt, i hvilke saa meget, vi i Regelen
gaa ligegyldige og tankeløse forbi, kan blive til for os. Men
han vil glæde sig i sin Grav over, at der er kommen en
værdig Efterfølger, som ad sin egen Vej er ført til at opdage
værdifulde Træk af Livets uvilkaarlige Udfoldelse og
til at værne om deres Betydning.

\piecerule
% --- p. 154 ---
\setcounter{footnote}{0}%

\opage{154}\essayhead{LUDVIG HOLSTEIN: „DEN GRØNNE MARK“}

Denne lille Bog kommer som et forfriskende Pust fra
Natur og Menneskeliv, vederkvægende efter al den lidet
dybtgaaende Tale, man har kunnet høre herhjemme i den
sidste Tid, en Tale, i hvilken hverken det sunde og frie
Menneskeliv eller Forskningens Arbejde er komne til deres
Ret. Bogen kommer fra den af vore Digtere, som mest har
levet i Naturen, som har fyldt vor Fantasi med de fineste
og friskeste Naturbilleder, og som nu godtgør, at han kan
se paa menneskeligt Sjæleliv og menneskelig Historie med
samme klare Øjne som paa Naturen omkring os. Vi har i
høj Grad en saadan Digter Behov, thi vi gaar ofte, Lærde
og Ulærde, uforstaaende ikke blot forbi Naturlivet omkring
os, men ogsaa forbi det indre Liv i Menneskers Sind.

Ludvig Holstein stiller sig paa rent menneskelig Grund.
Al Religion er for ham tilsidst fremgaaet af menneskelige
Sjæle og deres Oplevelser. Her er den grønne Mark, hvor
alle Aandslivets Planter, de religiøse som alle andre, gror.
Her voxer Livets Religion frem, med sin dybe Rod i menneskelig
Livsfølelse, i hvilken alle dybtgaaende Erfaringer
nedfældes. Det er Digterens Kald at tolke Livet, ikke mindst
hvor det rører sig dybest. Og Troen selv er den største af
alle Digtere.

I denne sidste Sætning mærkes Forfatterens personlige
Stade. Den siger i sig selv ikke nok, thi en Spotter kunde
ogsaa udtale den. Men den udtales her med Alvor og fin
Forstaaelse for religiøse Følelsers og Tankers uvilkaarlige
Udfoldelse under Livets Kamp, saavel som med aaben Sans
og god Forstaaelse for, hvad Religionshistorien kan lære
os. --- Jeg tør maaske minde om, at jeg ofte, første Gang i
% --- p. 155 ---
\setcounter{footnote}{0}%
\opage{155}min Bog om Kierkegaard (1893), har kaldt Religion for
Livspoesi. Men --- ak! --- hvad kan det nytte, at en Filosof
taler om Poesi! Der maa en Digter til. Og nu er han kommen
som en Digter skal, ikke med Forsøg paa at stoppe
Teologiens Huller med Lyrik, men med Livets og Naturens
friske Grønt om sin Pande. Lad os da se, hvorledes han
lægger sine Tanker til Rette. Jeg kan desværre ikke fremstille
dem med den rige Poesi, hvormed hans Tale fremtræder,
skønt han her, for første Gang, skriver i Prosa. ---

Holstein kræver „Studier i Marken“. Gaar vi Iagttagelsens
Vej, opdager vi let, hvor mange forskellige Anskuelser
der findes blandt Mennesker, selv blandt dem, der regner
sig til samme Konfession. Enhver skaber sit Paradis
efter sit Behov, og der er god Brug for de mange Boliger,
Jesus talte om. Men først naar vi forlader Menneskeverdenen
og gaar ud i den frie Natur, staar vi overfor det Mysterium,
der indeholder Nøglen til alle Mysterier, først og
fremmest til Myten om Moderen og Barnet, al Religions
Grundmyte, stammende fra Naturens eget Hjerte: „Hvert
groende Græsstraa er et Mysterium, ja Mysteriet selv. Som
Hyrderne staar vi Ansigt til Ansigt med Alnaturens Tilblivelsesmysterium.“
Oehlenschläger har allerede set det i
„Jesu Liv i den aarlige Natur“: „Hvert Aar, naar Taagerne
flygte hen, da fødes det lille Barn Jesus igen.“ --- Og herude,
fortsætter Holstein, paa den grønne Mark skal Maria
og Barnet forblive, Menneskemoderen, Tilblivelsesmysteriet
selv, ufatteligt og evigt, Virkelighed og Illusion i Et, Urmyten,
af hvis Skød alle Religioner er rundne! --- Ingen
Dualisme! I Menneskelivet som i Naturen kundgør sig den
store Fødsels-, Døds- og Fornyelsesproces.

Som Holstein her uvilkaarlig kommer over fra Natur til
Historie formedelst begges inderlige Sammenhæng, saa
sporer man ogsaa hans historiske Sans, idet de fororientalske
Religioners Myter om Gudemoderen og hendes Søn
og de dertil knyttede Mysterier kommer til at staa som Forgængere
for de kristelige Forestillinger. Hans historiske
Sans viser sig ogsaa i den Bemærkning, at Kirken som
% --- p. 156 ---
\setcounter{footnote}{0}%
\opage{156}Kirke naaede sit Højdepunkt i det 13. og 14. Aarhundrede.
Protestantismen er for ham en Halvhed, naar den vil optræde
med Autoritet. Den er kun konsekvent, naar den bygger
paa Samvittighedens Frihed. Ude paa den grønne Mark
tror vi paa Livets naturlige Vækst, baaren af Villie og ledet
af ærlig Tænken. Vi synger Psalmer --- men ikke alle kirkelige
Psalmer kan vi bruge. Vi ser klart, at Tiden gaar lange
Veje, og vi véd ikke selv, til hvilken Bygning vi samler
Sten. Men kigge bag Tæppet kan vi ikke. ---

Holstein er jo ikke den eneste af vore Digtere, der har
taget Ordet i Drøftelsen af det religiøse Spørgsmaal. Før
ham har Helge Rode givet sit Indlæg, ikke saa meget som
Digter (skønt mere end han vil være ved), men som Kritiker
af Videnskabernes Stilling overfor Kirken. Holstein
bebrejder Rode, at han fastholder den kristelige Dualisme,
der skyldes Indflydelse af Paulus' Tankegang, men indvender
især imod ham, at hans Standpunkt er „okkult Rationalisme“.
Dette sidste Udtryk karakteriserer træffende en
Opfattelse, der vil godtgøre Miraklers Mulighed ved at henvise
til den nyere Tids for læge Sjæle overraskende Opfindelser.
De gamle rationalistiske Præsters Tankegang var
den samme som Helge Rodes. --- Jeg betragter det som
overordentlig betydningsfuldt, at en Digter her har imødegaaet
en anden Digter.

Noget uklar staar Holstein overfor den saakaldte Bibelkritik
(der burde kaldes Bibelforskning, i Analogi med
Homer- og Platonforskning). Undertiden taler han om „Bibelkritiken“
i al Almindelighed, undertiden om „en vis Bibelkritik“,
som afvises med særligt Eftertryk, og undertiden
om „den objektive Bibelkritik“. Der var dog al Grund til at
holde sig til denne sidste alene, der er den egentlig videnskabelige,
og der derfor, som al sand Videnskab, har klart
Blik for de Vanskeligheder, der stadig staar tilbage. Denne
Art af Forskning vil samstemme i Holsteins Ytring om, at
det, hvad Kristendommens Oprindelse angaar, „neppe nogensinde
vil være muligt at udrede alle Traadene“.

Siden Talen er om Bibelforskningen, vil jeg bemærke, at
% --- p. 157 ---
\setcounter{footnote}{0}%
\opage{157}mange Forskere ikke vil være enige med Holstein i den
Vægt, han lægger paa, at Kristendommens Udvikling, særlig
i paulinsk og johanneisk Retning, skulde kunne forstaas
ud fra Sammenhængen mellem det gamle og det nye
Testamente. Utvivlsomt har visse orientalske Indflydelser
her gjort sig gældende, hvad det ikke vilde være svært for
Holstein at gaa ind paa, da han lægger saa megen Vægt paa
„Urmyten“ om Moderen og Sønnen. ---

Hvad den fremtidige Udvikling angaar, gør Holstein
først opmærksom paa, at der her i Landet er indtraadt en
Krise, fordi den Tid er forbi, i hvilken „dansk Aandslivs
Højrenaissance uhæmmet af Kirkens Dogmer kunde blomstre
under dens Tag“, saa at „Kirken, sine Dogmer til
Trods, blev sin Tids Formidler mellem Kristendom og
Aandsliv. Siden er Aandslivet gaaet sine egne Veje, og den
lutherske Kirke staar som en Anakronisme“ uden Sammenhæng
med Aandslivet i det hele. Dog ønsker Holstein ikke
Kirken skilt fra Staten, da det vil avle Præsteherredømme
og snæver Ortodoksi. Det gælder, om Kirken vil sætte Aanden
over Dogmet. Holsteins sidste Ord er: Lad Klokkerne
ringe og Dogmerne sove!

Med dette sidste Ord er vistnok ogsaa den virkelige Tilstand
i vort Folk udtrykt. Man vil ikke af med Klokkerne,
hører dem gerne ringe og stemmes ved dem, ligesom idethele
ved den liturgiske Del af Gudstjenesten. Men man vil
ikke høre for meget om Dogmerne. Historien maa vise,
hvor længe denne aandelige Tilstand kan holde sig. Baade
de, der staar udenfor Folkekirken, og de, der staar indenfor
den, vil med Spænding følge den videre historiske Udvikling
her. Og alle vil de have Grund til at være taknemmelige
for den Henvisning til Livets grønne Mark, som en af
vore ypperste Digtere i den foreliggende Bog har givet os.

\piecerule

% --- p. 158 ---
\setcounter{footnote}{0}%

\opage{158}\essayhead{SHAKESPEARE’S „STORMEN“}

For alle, som beundrer Shakespeare og lever i hans Værker,
forestaar der en stor Begivenhed, idet det kongelige
Teater nu [1926] vil opføre det Drama, der sætter Kronen paa
hans Livsarbejde. Jeg for mit Vedkommende har deri fundet
Tilskyndelse til paany at fordybe mig i dette vidunderlige
Værk og har faaet Lyst til at nedskrive noget om, hvorledes
jeg forstaar det. Jeg er ikke Shakespeareforsker, kun
Shakespearelæser, og jeg maa paa Forhaand udtale, hvor
meget jeg skylder de udmærkede engelske Shakespearefortolkere,
navnlig Eduard Dowden og A. C. Bradley. Det hører
til mine bedste Erindringer fra mine Besøg i England,
at jeg i Efteraaret 1904 paa Universitetet i Oxford
hørte Bradley's Foredrag om Shakespeare som Menneske
(„Shakespeare the man“), der senere er trykt i hans „Oxford
Lectures“.

Jeg vil forsøge at give Dramaets Indhold i korte Træk.

Prospero, Hertug af Milano, blev ved Svig fordreven fra
sit Land af sin egen Broder, til hvem han havde overladt
Landets Styrelse, medens han selv fordybede sig i sine Studier.
Nu lever han med sin Datter paa en Ø langt ude i
Havet, stadig fordybet i sin Gransken og optagen af at undervise
Datteren, Miranda, der i sin spæde Alder blev reven
bort fra alt andet menneskeligt Samkvem end Faderens.
Prospero har som Faust fordybet sig i Magi, men hans
Tankeliv har under hans Ophold paa Øen vundet en saadan
Magt, at han kan beherske Elementerne uden som Faust
at behøve en saa betænkelig Medhjælper som Mefistofeles.
Dog har han to tjenende Væsener, som han har fundet paa
Øen. Det ene, Ariel, er et æterisk Væsen, som af en Hex,
% --- p. 159 ---
\setcounter{footnote}{0}%
\opage{159}der forhen herskede paa Øen, var bleven holdt bunden i
haard Pine, indtil Prospero befriede ham. Ariel bevæger
sig med Tankens Fart, handler rask og klogt i sin Mesters
Aand og kan give ham gode Raad. Ja, skønt han er et æterisk
Væsen, uden menneskeligt Kød og Blod, kan han sige
Prospero, hvad højsindet Menneskelighed byder; han véd,
hvorledes han vilde handle, hvis han var et Menneske. Midt
i al Iver i sin Velgørers Tjeneste længes han dog efter at
blive løst af denne Tjeneste for igen at blive Et med Naturens
reneste Elementer. Den anden Tjener, Caliban, er en
Mellemting mellem Dyr og Djævel, en Søn af Hexen. Prospero
har forgæves søgt at vække menneskeligt Sind hos
ham. Han kan ikke glemme, at han er Søn af Øens tidligere
Herskerinde. Han vedbliver at være dyrisk og hævngerrig
og bliver derfor behandlet strengt, især efterat han havde
forsøgt at voldtage Miranda. Det menneskelige Sprog, Prospero
har lært ham, bruger han mest til at forbande sin
Herre. Det eneste menneskelige Træk hos ham er hans
Glæde ved at høre Musik. Ogsaa han længes efter Frihed,
men det er raa Elementers Tøjlesløshed, han higer efter.
Han tvinges til at trælle for Prospero's materielle Fornødenheder,
medens Ariel arbejder for Gennemførelsen af Mesterens
ideale Planer.

Den Magi, Prospero gennem sin Gransken er bleven indviet
i, benytter han nu til at faa sine Fjender i sin Magt og
derved Befrielse og Lykke for sig og sin Datter. I en af
ham fremkaldt Storm strander hans forræderiske Broder,
Antonio, og dennes Hjælper, Kong Alonzo af Neapel, paa
Øen. Ariel sørger for, at ingen mister Livet. Og nu strejfer
de Skibbrudne i forskellige Hold omkring paa Øen. En liflig
Musik leder Kongens Søn Ferdinand til Prosperos Hule,
hvor han og Miranda strax betages af hinanden. Et af de
yndefuldeste Elskovsforhold, Shakespeare har skildret, udfolder
sig nu. Det volder Prospero den største Glæde: „Saa
glad som de kan jeg ikke blive, og dog kunde intet gøre mig
større Glæde!“ Dog sætter han Prinsen paa en haard Prøve
gennem strengt Arbejde, inden han giver sit Minde.

% --- p. 160 ---
\setcounter{footnote}{0}%
\opage{160}Da den vigtigste Gruppe af de Skibbrudne, Antonio,
Alonzo og deres Hofmænd, vides at være under Ariels Magt,
lader Prospero i sin store Glæde over det unge Pars Lykke
Ariel fremtrylle en Række Scener, i hvilke Gudinder bringer
deres Lykønskninger til Pagten, og Nymfer og Høstfolk
danser. Og saa kommer en af de mest karakteristiske Scener
i Stykket. Midt under det fremtryllede Skuespil farer
Prospero pludselig op og byder Aanderne forsvinde. Pludselig
var den Tanke faret igennem ham, at alt ikke endnu
var sikkert. Han havde under Optagetheden af de Syner,
der fremtryllede sig, glemt, at Uhyret Caliban ude i Skoven
havde truffet to Mænd fra Skibet, der havde reddet et godt
Forraad af Vin, og at nu alle tre var paa Vej for at angribe
Prospero og bemægtige sig hans Datter, ved hvis Side den
ene af de Sammensvorne, Neapelkongens Køgemester, i
hvem Caliban saa en Aabenbaring af sin fordums Gud, saa
skulde herske over Øen. Intet Under, at Prospero blev bestyrtet
og oprørt, da han kom til at tænke paa, hvad der
var i Gære imod ham, medens han var begyndt at fejre det
Opgør, hvor hans Højsind overfor hans Fjender skulde
volde Fred og Glæde. Hans Datter havde aldrig set ham
saa forbitret. Men han fatter sig og beroliger sig og de andre
med den Tanke, at ikke blot saadanne Syner, de nylig
havde havt, maatte opløses i intet, men at det vilde gaa
saaledes med alt Stort og Herligt, med selve vor Jord.
Vort Liv er jo kun som et af de Drømmesyner, der nys
var fremme.

Det gaar nu op for Prospero, at hans Møje med at opdrage
Caliban har været forgæves, og at han maa gaa strengt
frem, før han lader Højsindet raade. Ariel bringer de brutale
Sammensvorne i hans Magt, de faar deres Straf, og da
Prospero senere sejler bort fra Øen, lader han Caliban tilbage.
Forbedret sig har Uhyret ikke, men han har indset,
hvilken Dumhed han begik ved at se en Gud i en drukken
Køgemester. --- Paa et Teater i London saa man tilsidst Skibet
med Prospero og hans Fjender, der nu havde faaet Tilgivelse,
sejle bort fra Kysten, medens Caliban liggende paa
% --- p. 161 ---
\setcounter{footnote}{0}%
\opage{161}en Fjældkant følger dem med Øjnene. Han har aabenbart
ikke forstaaet en Smule af, hvad der er sket. ---

Efterhaanden samles nu alle de andre Skibbrudne ved
Prosperos Hule, førte af Ariel, snart under liflig Musik,
„den bedste Hjælper for et uroligt Sind“, snart under truende
Lyde. Kongen og Antonio vaagner op af den Tryllesøvn,
Ariel har bragt dem i, og Prospero modtager dem med
Opfordring til at glemme, hvad der er voldet og sket, og
oprette det saa godt som muligt. Forbindelsen mellem Kongens
Søn og Miranda skal besegle Forsoningen. Hvis der
har været nogen Tvivl hos Prospero om, hvorledes han
skulde behandle dem, der saa helt er i hans Magt, saa optræder
Ariel som hans gode Genius. Ariel mener, at hvis
hans Mester havde set disse Fjenders Lidelser, vilde han
ynkes over dem; det vilde Ariel selv, hvis han havde været
et menneskeligt Væsen! Prospero, der som Menneske endnu
bedre end Ariel kan føle deres Elendighed, følger hans Opfordring.
Hans Fjenders arge Færd har saaret ham i hans
Inderste, men Højsindet vinder Magt over Forbitrelsen og
Aandskraften over Hævntrangen. Hans Følelse af sin Magt
over Mennesker som over Naturen. Bevidstheden om, at
hans Planer er lykkedes, og om at Oprejsningen, der er
vunden ved hans egen Kraft, er nær, giver ham Overlegenhed
til at gaa „den sjeldnere Vej“, den, der fører ud over
Trangen til at gengælde ondt med ondt. Han har lært, at
det onde i Naturen og i Menneskelivet ikke kan forsvinde
af Verden, men han har tillige lært, at det kan overvindes
med Taalmodighed, uden at man behøver at svigte det
ædleste, der rummes i Ens eget Sind. ---

Shakespeare var ikke ene om at betragte Højsind som
den ypperste Karakteregenskab. I Renaissancetidens Tænkning
kommer denne Idé atter og atter frem. Hos en Række
af Filosofer har den vaagnende Naturforsken og det udvidede
Blik over menneskelige Forhold ført til at se, at
Selviskhed og Hævntrang er en Begrænsning af Sindets
Kraft, tilsidst en Svaghed. Denne Tankegang fuldendes senere
i Spinozas Lære om det frie Menneske, der lever ud af
% --- p. 162 ---
\setcounter{footnote}{0}%
\opage{162}sig selv, og hvis Kraft ikke blot formaar at hævde sig selv,
men vise sin Overlegenhed gennem at hjælpe andre. Generositas
som Udtryk for aandelig Kraft er en Hovedtanke i
Renaissancens Etik. Det er ikke min Mening, at Shakespeare
skulde være paavirket af Tidens Filosofer, men gennem
sine Erfaringer og gennem Arbejdet i sin Kunst er
han ad sin egen Vej kommen til samme Resultat som hans
Tids Tænkere, og han har formaaet at udtrykke det i et
Værk, som ingen uden han kunde skabe. Det er Troen paa
Livet, der fører ud over Sorg som over Had. Da Kong Alonzo
trods Forsoningen ikke kan glemme, hvad der er sket, stanser
Prospero ham med de Ord: „Lad os ikke bebyrde vor
Erindring med en Sorg, hvis Tid er forbi!“ --- Den Tankegang,
der havde ladet Prospero med Sindsro tænke paa, at
ikke blot Fantasiens Paradiser, men ogsaa alt Menneskeværk,
ja, Jorden selv, vi lever paa, er Opløsning underlagt,
fører ham til at erkende at ogsaa den Sorg, der er voldet
ham, og den Brøde, som er øvet mod ham, er denne Lov
underlagt.

Shakespeareforskerne er enige om, at den aandelige Højhed,
der skildres i „Stormen“ og har Fred og Forsoning til
Følge, idethele er karakteristisk for den sidste Periode af
Shakespeares Digtervirksomhed. Samme Stemning raader
i „Cymbeline“ og i „Et Vintereventyr“. Trods al Forbrydelse
og Forurettelse sejrer Trofasthed og Kærlighed. Imogen
og Hermione er Naturer som Prospero, skønt de ingen
anden Magi raader over end den taalmodige Kærlighed.
Alle disse Dramaer betegner en ejendommelig Periode i
Shakespeare's Udvikling, som iøvrig ikke blot er karakteriseret
ved den overlegne Mildhed, men ogsaa ved den ganske
særegne Ynde, med hvilken unge Kvinder (Miranda,
Imogen, Perdita) skildres. Hvad der er ejendommeligt for
„Stormen“ i Forhold til de andre Dramaer fra samme Tid,
er, at Shakespeare gennem Prospero ligefrem og tydeligere
end i andre Dramaer udtaler den Livsopfattelse, han har
vundet, og at han nu bebuder at ville ophøre med sine Dig%
% --- p. 163 ---
\setcounter{footnote}{0}%
\opage{163}terværker. Han tager Afsked med sine Tilhørere, idet han
lader Prospero bryde sin Tryllestav, begrave den dybt i Jorden,
og sænke sin magiske Bog ned i Havet. Livets Modsætninger
har nu fortonet sig for ham. Han har ført sin
aandelige Kamp med de Erfaringer og Oplevelser, der faldt
i hans Lod, hvad enten de greb ind i hans eget personlige
Liv eller ikke. Han har kæmpet som Menneske og som
Kunstner, og „Stormen“ viser, at han har vundet Sejr i
begge Egenskaber. Nu gaar han tilbage til Stratford on
Avon, hvor han med den praktiske Sans, der stedse synes at
have virket ved Siden af hans kunstneriske Trang, havde
beredt sig et smukt Hjem. Han, der efter Sagnet maatte
flygte som ung Mand fra sin Fødeby, ejede nu en af de
stateligste Gaarde der, og som Ariel glædede sig til sin Frihed
i Æteren, har Shakespeare glædet sig til igen at færdes
ved Avons Bred. En udmærket Shakespeareforsker har da
ogsaa i Shakespeare's senere Produktion mærket Duften af
Engene ved Avon. (For mit Vedkommende har omvendt
en Vandring over disse Enge gjort Indtrykket af Shakespeares
smaa Naturbilleder mere friske og levende, end
den blotte Læsning vilde formaa).

Den Shakespeare, vi møder i den sidste Periode af hans
Digtervirksomhed, og mest udpræget i „Stormen“, maa ses
i Sammenhæng med den, der tidligere, omkring 40 Aar
gammel, digtede Hamlet, Othello, Macbeth og Lear, disse
vældige Tragedier, i hvilke Menneskelivets Modsætninger
kommer saa skarpt og hensynsløst frem, og i hvilke Undergang
ofte synes at være det sidste Ord. Som Bradley har
vist, handler disse Tragedier for en væsentlig Del om de
Skuffelser og den Modgang, aabne og frejdige Naturer kan
møde paa deres Vej. En saadan Natur var Hamlet før Faderens
Mord og Moderens Hastværksbryllup, før Vennernes
Forræderi og Ofelia's Spioneri, en saadan var ogsaa
Othello, før Jago djævelsk vakte hans Tvivl, og Lear, før
Døtrenes Utaknemmelighed viste sig. En saadan var fremdeles
Macbeth før Hexesøstrenes Spaadom og hans Hustrus
% --- p. 164 ---
\setcounter{footnote}{0}%
\opage{164}Opæggelse. Og Bradley har træffende lagt Vægt paa, at Shakespeare
selv af Samtidige karakteriseres som en aaben, fri
og fin Natur. Han maa have gjort Erfaringer om, hvad
Livet kan bringe af Ulykke og Ondskab, selv om det ikke
umiddelbart greb ind i hans egen Skæbne. Hans Fantasi
er i hvert Tilfælde bleven greben deraf, og han har ikke
hvilet, før han i store og stærke Træk fik skildret, hvad han
saa. Men han har ikke tabt Troen paa det Godes Realitet.
Den hævder han i selve det Stores og Ædles Fald, og det
skildrer han som Trofasthed og Taalmodighed i Skikkelser
som Kent og Horatio. Han staar midt i Striden og lader
Modsætningerne prelle mod hinanden. Den samme aandelige
Kraft, der i „Stormen“ bar ham ud over Livets Modsætninger,
havde gjort ham det muligt i hine store Tragedier
at lade disse Modsætninger træde op imod hinanden
i deres hele Styrke. Det er Frygt og Medlidenhed, der kommer
til at raade, ikke den forklarede og inderlige Forstaaelse
og den taalmodige Udholdenhed, der fører Ordet i „Stormen“.
Ogsaa i „Kong Lear“ bliver der Tale om Undergang.
Da den gamle Konge kommer ind med Cordelias Lig, raaber
Kent: „Er det Alts Ophør, som er varslet?“, --- Edgar:
„eller Billede af den Rædsel, der kan ventes?“ --- og Albanerhertugen:
„Undergang og Ophør?“ (V, 3, 263ff.). Men
nu, i „Stormen“, er det nok med bevæget Sind, men ikke i
Rædsel, at Prospero tænker paa Undergangen. Der raader
nu en forklaret Stemning i Digterens Sind. Vi vilde ikke
glæde os ved „Stormen“, hvis vi ikke ved den samme Digters
Haand havde gennemlevet de tragiske Stemninger, han
i sin første Manddom følte Trang til at faa Udtryk for, og
der vilde ikke have været den store Vederhæftighed i Digterens
sidste Forkyndelse, hvis han ikke havde været i
Skærsilden, før han færdedes i Paradiset.

Det er ikke muligt uden videre at benytte Udbyttet af andres
Livserfaring, saa at man holder sig til den forklarede
Form, i hvilken deres Livsopfattelse kan fremtræde, efterat
de har prøvet Livet. Vi maa alle begynde forfra. Og saa ren
% --- p. 165 ---
\setcounter{footnote}{0}%
\opage{165}og klar end den Atmosfære er, vi føres ind i af Prospero og
Ariel, saa vil der dog være dem, der, naar de trænger til
Støtte under Livets Kamp, hellere tyer til „Kong Lear“ end
til „Stormen“. Selv om det Ædle og Store vies til Undergang,
er det dog opløftende at vide, at det har været og kan
være i Verden. Og det kan være styrkende at tage en Dyst
med de Tanker, der vækkes ved Lears og Cordelias Skæbne,
før man prøver paa at leve i „Stormen“s Æterverden. Stort
er det at have skænket os et Værk som „Stormen“, men
større er det, at dette Værk først blev skabt, efterat det var
levet, følt og tænkt i Stridens og Modsætningernes Verden.
Og det har han gjort, hvem Goethe kaldte „William, Stern
der höchsten Höhe“. Han svigter os ikke, hverken under
Kampen eller i Nederlaget, eller naar Sejren er vunden.

\piecerule

\end{document}
